\chapter{Sequ\^encias e s\'eries} \label{seq:chapter}

%%%%%%%%%%%%%%%%%%%%%%%%%%%%%%%%%%%%%%%%%%%%%%%%%%%%%%%%%%%%%%%%%%%%%%%%%%%%%%

\section{Sequ\^encias e limites}
\label{sec:seqsandlims}

%mbxINTROSUBSECTION

\sectionnotes{2,5 aulas}

Em essência, a análise trata de limites. O tipo mais básico de limite é o
limite de uma sequência de números reais.
Já vimos sequências serem usadas informalmente. Vamos apresentar a definição
formal.

\begin{defn}
Uma \emph{\myindex{sequ\^encia}} (de números reais) é uma função
$x \colon \N \to \R$. Em geral, em vez de $x(n)$, denotamos por $x_n$ o
$n$-ésimo termo da sequência.
Para denotar uma sequência, escrevemos%
\footnote{É comum usar $\{ x_n \}$ ou $\{ x_n \}_n$ por brevidade.}%
\glsadd{not:sequence}%
\begin{equation*}
\{ x_n \}_{n=1}^\infty.
\end{equation*}

Uma sequência $\{ x_n \}_{n=1}^\infty$ é \emph{limitada}\index{sequ\^encia limitada} se
a função correspondente é limitada. Isto é, existe $B \in \R$ tal que
\begin{equation*}
\sabs{x_n} \leq B \qquad \text{para todo } n \in \N.
\end{equation*}
Em outras palavras, a sequência $\{x_n\}_{n=1}^\infty$ é limitada sempre
que o conjunto $\{ x_n : n \in \N \}$ é limitado.
Definimos de modo análogo os termos
\emph{limitada inferiormente}\index{sequ\^encia limitada inferiormente} e
\emph{limitada superiormente}\index{sequ\^encia limitada superiormente}.
\end{defn}

Quando precisamos apresentar uma sequência concreta, muitas vezes damos
cada termo por meio de uma fórmula em função de $n$.
Por exemplo, $\{ \nicefrac{1}{n} \}_{n=1}^\infty$ representa a sequência
$1, \nicefrac{1}{2}, \nicefrac{1}{3}, \nicefrac{1}{4},
\nicefrac{1}{5}, \ldots$.
A sequência $\{ \nicefrac{1}{n} \}_{n=1}^\infty$ é limitada (basta tomar
$B=1$). Por outro lado, a sequência $\{ n \}_{n=1}^\infty$ representa
$1,2,3,4,\ldots$, e não é limitada (por quê?).

Embora a notação de uma sequência seja semelhante\footnote{Em \cite{BS},
usa-se $(x_n)_{n=1}^\infty$ para denotar uma sequência, em vez de
$\{ x_n \}_{n=1}^\infty$, como em \cite{Rudin:baby}. Ambas as notações são
comuns.} à de um conjunto, as noções são distintas. Por exemplo, a sequência
$\bigl\{ {(-1)}^n \bigr\}_{n=1}^\infty$ é a sequência
$-1,1,-1,1,-1,1,\ldots$, enquanto o conjunto de seus valores, a
\emph{imagem da sequência}\index{imagem de uma sequ\^encia}, é simplesmente
o conjunto $\{ -1, 1 \}$. Escrevemos esse conjunto como
$\bigl\{ {(-1)}^n : n \in \N \bigr\}$.
% When ambiguity can arise, we
%use the words \emph{sequence} or \emph{set} to distinguish the two
%concepts.

Outro exemplo de sequência é a \emph{\myindex{sequ\^encia constante}}.
Trata-se da sequência $\{ c \}_{n=1}^\infty = c,c,c,c,\ldots$, formada por
um único valor constante $c \in \R$, repetido indefinidamente.

%We now get to the idea of a
%\emph{limit of a sequence}\index{limit!of a sequence}.
%We will see in \propref{prop:limisunique}
%that a limit, if it exists, is unique.
%It makes sense to talk about \emph{the} limit of a sequence.

\begin{defn}
A sequência $\{ x_n \}_{n=1}^\infty$ \emph{converge} para um número
$x \in \R$ se, para todo $\epsilon > 0$, existe $M \in \N$ tal que
$\sabs{x_n - x} < \epsilon$ para todo $n \geq M$.
O número $x$ é chamado de \emph{limite} da sequência.
Provaremos que o limite $x$ é único e, por isso, escrevemos%
\footnote{No texto, isso pode ser renderizado como $\lim_{n\to\infty} x_n$.}
\glsadd{not:limseq}%
\begin{equation*}
\lim_{n\to \infty} x_n \coloneqq x .
\end{equation*}

Uma sequência que converge é chamada de \emph{convergente}\index{sequ\^encia convergente}.
Caso contrário, dizemos que a sequência \emph{diverge} ou que é
\emph{divergente}\index{sequ\^encia divergente}.
\end{defn}

Em breve, na \propref{prop:limisunique}, mostraremos que, se existir, o
limite $x$ é sempre único. Faz sentido falar \emph{do} limite de uma
sequência; basta mostrar que ela converge para um único número.
Nos próximos exemplos, vamos supor que já provamos a unicidade dos limites.

Intuitivamente, dizer que o limite é $x$ significa que, a partir de certo
índice, todos os termos da sequência ficam próximos de $x$. Mais
precisamente, podemos nos aproximar arbitrariamente do limite, desde que
avancemos o suficiente na sequência. Isso não significa que chegaremos ao
limite. É possível, e bastante comum, que nenhum termo $x_n$ da sequência
seja igual ao limite $x$.
Ilustramos o conceito na \figureref{figsequenceconvergence}. Primeiro,
pensamos na sequência como um gráfico, pois ela é uma função definida em
$\N$. Em seguida, também a representamos como uma sequência de pontos
identificados na reta real.

\begin{myfigureht}
\myincludepdft{sequence-convergence_full}{%
Dois diagramas.
O primeiro é o gráfico de uma sequência: os valores estão no eixo vertical
e o índice, no eixo horizontal. No eixo vertical, marca-se o número x e
traça-se uma reta horizontal. Uma faixa de x menos épsilon a x mais épsilon está
sombreada. Até o índice 7, indicado por $M$, a maioria dos termos da
sequência está fora da faixa sombreada. Todos os termos seguintes estão
dentro dela.
No segundo diagrama, o eixo x é horizontal, com o número x marcado e os
valores x sub 1, x sub 2, até x sub 10 indicados por marcas.}
\caption{Ilustração da convergência.
Acima, mostramos os dez primeiros termos da sequência como um gráfico, com
$M$ e o intervalo em torno do limite $x$ assinalados.
Abaixo, os termos da mesma sequência estão marcados na reta real.
\label{figsequenceconvergence}}
\end{myfigureht}

Quando escrevemos $\lim_{n\to\infty} x_n = x$ para algum número real $x$,
estamos afirmando duas coisas: primeiro, que $\{ x_n \}_{n=1}^\infty$ é
convergente; segundo, que seu limite é $x$.

A definição acima é uma das mais importantes da análise e é necessário
compreendê-la perfeitamente. O ponto-chave é que, dado \emph{qualquer}
$\epsilon > 0$, podemos encontrar um $M$. Esse $M$ pode depender de
$\epsilon$; portanto, só o escolhemos depois de conhecer $\epsilon$.
Vamos ilustrar a convergência com alguns exemplos.

\begin{example}
A sequência constante $1,1,1,1,\ldots$ converge para 1. Para todo
$\epsilon > 0$, tome $M = 1$. Temos
$\sabs{x_n-x} = \sabs{1-1} < \epsilon$ para todo n.
\end{example}

\begin{example}
Afirmação: \emph{A sequência $\{ \nicefrac{1}{n} \}_{n=1}^\infty$ é convergente e}
\begin{equation*}
\lim_{n\to \infty} \frac{1}{n} = 0 .
\end{equation*}
Demonstração: Dado $\epsilon > 0$, encontre $M \in \N$ tal que
$0 < \nicefrac{1}{M} < \epsilon$
(a \hyperref[thm:arch:i]{propriedade arquimediana} em ação).
Para todo $n \geq M$,
\begin{equation*}
\sabs{x_n - x}
=
\abs{\frac{1}{n} - 0} = \abs{\frac{1}{n}} = \frac{1}{n} \leq \frac{1}{M} < \epsilon .
\end{equation*}
\end{example}

\begin{example}
A sequência $\bigl\{ {(-1)}^n \bigr\}_{n=1}^\infty$ é divergente.
Demonstração: Se existisse um limite $x$, então, para $\epsilon = \frac{1}{2}$,
deveria existir um $M$ que satisfizesse a definição. Suponha que tal $M$
exista. Para um $n$ par com $n \geq M$, calculamos
\begin{equation*}
\nicefrac{1}{2} > \sabs{x_n - x}  = \sabs{1 - x}
\qquad \text{e} \qquad
\nicefrac{1}{2} > \sabs{x_{n+1} - x}  = \sabs{-1 - x} .
\end{equation*}
Obtemos então uma contradição:
\begin{equation*}
2 = \babs{1 - x - (-1 -x)} \leq
\sabs{1 - x} + \sabs{-1 -x} < \nicefrac{1}{2} + \nicefrac{1}{2} = 1 .
\end{equation*}
\end{example}

\begin{prop} \label{prop:limisunique}
Uma sequência convergente tem um único limite.
\end{prop}

A demonstração desta proposição apresenta uma técnica útil em análise.
Muitas demonstrações seguem o mesmo esquema geral. Queremos mostrar que certa
quantidade é zero. Usando a desigualdade triangular, escrevemos essa
quantidade como a soma de duas outras e estimamos cada uma por números
arbitrariamente pequenos.

\begin{proof}
%NOTE: should be word for word the same as 7.3.3
Suponha que $\{ x_n \}_{n=1}^\infty$ tenha limites $x$ e $y$.
Tome um $\epsilon > 0$ arbitrário.
Pela definição, encontre $M_1$ tal que, para todo $n \geq M_1$,
$\sabs{x_n-x} < \nicefrac{\epsilon}{2}$. De modo semelhante, encontre
$M_2$ tal que, para todo $n \geq M_2$,
$\sabs{x_n-y} < \nicefrac{\epsilon}{2}$.
Agora tome $n$ tal que $n \geq M_1$ e $n \geq M_2$, e estime
\begin{equation*}
\begin{split}
\sabs{y-x}
& =
\babs{x_n-x - (x_n -y)} \\
& \leq
\sabs{x_n-x} + \sabs{x_n -y} \\
& <
\frac{\epsilon}{2} + \frac{\epsilon}{2} = \epsilon .
\end{split}
\end{equation*}
Como $\sabs{y-x} < \epsilon$ para todo $\epsilon > 0$, temos
$\sabs{y-x} = 0$ e $y=x$. Portanto, o limite (se existir) é único.
\end{proof}

\begin{prop}
Uma sequência convergente $\{ x_n \}_{n=1}^\infty$ é limitada.
\end{prop}

\begin{proof}
Suponha que $\{ x_n \}_{n=1}^\infty$ convirja para $x$. Então existe
$M \in \N$ tal que, para todo $n \geq M$,
$\sabs{x_n - x} < 1$. Para $n \geq M$,
\begin{equation*}
\begin{split}
\sabs{x_n} & = \sabs{x_n - x + x}
\\
& \leq \sabs{x_n - x} + \sabs{x}
\\
& < 1 + \sabs{x} .
\end{split}
\end{equation*}
O conjunto $\bigl\{ \sabs{x_1}, \sabs{x_2}, \ldots, \sabs{x_{M-1}} , 1+\sabs{x} \bigr\}$
é finito; portanto, seja
\begin{equation*}
B \coloneqq \max \bigl\{ \sabs{x_1}, \sabs{x_2}, \ldots, \sabs{x_{M-1}}, 1+\sabs{x} \bigr\} .
\end{equation*}
Então, para todo $n \in \N$,
\begin{equation*}
\sabs{x_n} \leq B. \qedhere
\end{equation*}
\end{proof}

A sequência $\bigl\{ {(-1)}^n \bigr\}_{n=1}^\infty$ mostra que a recíproca
não vale. Uma sequência limitada não é necessariamente convergente.

\begin{example}
Vamos mostrar que $\left\{ \frac{n^2+1}{n^2+n} \right\}_{n=1}^\infty$
converge e que
\begin{equation*}
\lim_{n\to\infty} \frac{n^2+1}{n^2+n} = 1 .
\end{equation*}

Dado $\epsilon > 0$, encontre $M \in \N$ tal que $\frac{1}{M} < \epsilon$.
Então, para todo $n \geq M$,
\begin{equation*}
\begin{split}
%\abs{\frac{n^2+1}{n^2+n} - 1} & =
%\abs{\frac{n^2+1 - (n^2+n)}{n^2+n}} \\
\abs{\frac{n^2+1}{n^2+n} - 1}  =
\abs{\frac{n^2+1 - (n^2+n)}{n^2+n}}
& =
\abs{\frac{1 - n}{n^2+n}} \\
& =
\frac{n-1}{n^2+n} \\
& \leq 
\frac{n}{n^2+n} 
 =
\frac{1}{n+1}  \\
& \leq \frac{1}{n}
\leq \frac{1}{M} < \epsilon .
\end{split}
\end{equation*}
Portanto,
$\lim_{n\to\infty} \frac{n^2+1}{n^2+n} = 1$.
Este exemplo mostra que, às vezes, para obter o que queremos, precisamos
descartar algumas informações e assim conseguir uma estimativa mais simples.
\end{example}

\subsection{Sequ\^encias mon\'otonas}

O tipo mais simples de sequência é a sequência monótona. Verificar que uma
sequência monótona converge é tão fácil quanto verificar que ela é limitada.
Também é fácil encontrar o limite de uma sequência monótona convergente,
desde que possamos determinar o supremo ou o ínfimo de um conjunto contável
de números.

\begin{defn}
A sequência $\{ x_n \}_{n=1}^\infty$ é \emph{monótona não decrescente}\index{sequ\^encia mon\'otona n\~ao decrescente} se $x_n \leq x_{n+1}$ para todo $n \in \N$.
%
A sequência $\{ x_n \}_{n=1}^\infty$ é \emph{monótona não crescente}\index{sequ\^encia mon\'otona n\~ao crescente} se $x_n \geq x_{n+1}$ para todo $n \in \N$.
%
Se uma sequência é monótona não decrescente ou monótona não crescente,
podemos simplesmente dizer que ela é \emph{monótona}\index{sequ\^encia mon\'otona}.\footnote{Alguns autores usam o termo
\emph{monotônica}\index{sequ\^encia monot\^onica}.}
\end{defn}

Por exemplo, $\{ n \}_{n=1}^\infty$ é monótona não decrescente,
$\{ \nicefrac{1}{n} \}_{n=1}^\infty$ é monótona não crescente,
a sequência constante $\{ 1 \}_{n=1}^\infty$ é tanto monótona não
decrescente quanto monótona não crescente, e
$\bigl\{ {(-1)}^n \bigr\}_{n=1}^\infty$ não é monótona.
Os primeiros termos de uma sequência monótona não decrescente de exemplo
estão representados na \figureref{figsequenceincreasing}.

\begin{myfigureht}
\myincludegraphics{sequence-increasing}{%
Um gráfico dos primeiros termos de uma sequência, com o índice no eixo
horizontal. Cada termo da sequência é igual ou maior que o anterior; por
isso, os pontos que representam os valores sobem à medida que avançamos para
a direita no gráfico.}
\caption{Os primeiros termos de uma sequência monótona não decrescente,
representados em um gráfico.\label{figsequenceincreasing}}
\end{myfigureht}

\begin{thm}[Teorema da convergência monótona] \label{thm:monotoneconv}
\index{teorema da converg\^encia mon\'otona}%
Uma sequência monótona $\{ x_n \}_{n=1}^\infty$ é limitada se, e somente se,
é convergente.

Além disso, se $\{ x_n \}_{n=1}^\infty$ é monótona não decrescente e limitada, então
\begin{equation*}
\lim_{n\to \infty} x_n = \sup \{ x_n : n \in \N \} .
\end{equation*}
Se $\{ x_n \}_{n=1}^\infty$ é monótona não crescente e limitada, então
\begin{equation*}
\lim_{n\to \infty} x_n = \inf \{ x_n : n \in \N \} .
\end{equation*}
\end{thm}

\begin{proof}
Considere uma sequência monótona não decrescente
$\{ x_n \}_{n=1}^\infty$. Suponha, primeiro, que a sequência seja limitada,
isto é, que o conjunto $\{ x_n : n \in \N \}$ seja limitado. Seja
\begin{equation*}
x \coloneqq \sup \{ x_n : n \in \N \} .
\end{equation*}
Seja $\epsilon > 0$ arbitrário. Como $x$ é o supremo, deve existir pelo
menos um $M \in \N$ tal que $x_{M} > x-\epsilon$.
Como $\{ x_n \}_{n=1}^\infty$ é monótona não decrescente, é fácil ver
(por \hyperref[induction:thm]{indução}) que $x_n \geq x_{M}$ para todo
$n \geq M$. Portanto, para todo $n \geq M$,
\begin{equation*}
\sabs{x_n-x} = x-x_n \leq x-x_{M} < \epsilon  .
\end{equation*}
Assim, $\{ x_n \}_{n=1}^\infty$ converge para $x$.
Portanto, uma sequência monótona não decrescente e limitada converge.
Na outra direção, já provamos que uma sequência convergente é limitada.

A demonstração para sequências monótonas não crescentes fica como exercício.
\end{proof}

Uma sequência monótona não decrescente $\{ x_n \}_{n=1}^\infty$ é sempre
limitada inferiormente, pois $x_1 \leq x_2 \leq \cdots \leq x_n$ para todo
$n$; assim, $x_1$ é uma cota inferior. Portanto, para verificar se uma
sequência monótona não decrescente é limitada, basta verificar se é limitada
superiormente. De modo semelhante, uma sequência monótona não crescente é
sempre limitada superiormente; portanto, basta verificar se é limitada
inferiormente.

\begin{example}
Considere a sequência $\bigl\{ \frac{1}{\sqrt{n}} \bigr\}_{n=1}^\infty$.

A sequência é limitada inferiormente, pois
$\frac{1}{\sqrt{n}} > 0$ para todo $n \in \N$.
Vamos mostrar que ela é monótona não crescente. Começamos com
$\sqrt{n+1} \geq \sqrt{n}$ (por que isso é verdade?). Dessa desigualdade,
obtemos
\begin{equation*}
\frac{1}{\sqrt{n+1}} \leq \frac{1}{\sqrt{n}} .
\end{equation*}
Logo, a sequência é monótona não crescente e limitada inferiormente (e,
portanto, limitada). Pelo \thmref{thm:monotoneconv}, a sequência é
convergente e
\begin{equation*}
\lim_{n\to \infty} \frac{1}{\sqrt{n}}
=
\inf \left\{ \frac{1}{\sqrt{n}} : n \in \N \right\} .
\end{equation*}
Já sabemos que o ínfimo é maior ou igual a 0, pois 0 é uma cota inferior.
Seja $b \geq 0$ tal que $b \leq \frac{1}{\sqrt{n}}$ para todo $n$.
Elevando ambos os lados ao quadrado, obtemos
\begin{equation*}
b^2 \leq \frac{1}{n} \qquad \text{para todo } n \in \N.
\end{equation*}
Já vimos que isso implica $b^2 \leq 0$ (como consequência da
\hyperref[thm:arch:i]{propriedade arquimediana}). Como também temos
$b^2 \geq 0$, segue que $b^2 = 0$ e, portanto, $b = 0$.
Assim, $b=0$ é a maior das cotas inferiores, e
$\lim\limits_{n\to\infty} \frac{1}{\sqrt{n}} = 0$.
\end{example}

\begin{example}
Uma ressalva: mostrar que uma sequência monótona é limitada para poder usar
o \thmref{thm:monotoneconv} pode ser difícil.
A sequência
$\{ 1 + \nicefrac{1}{2} + \cdots + \nicefrac{1}{n} \}_{n=1}^\infty$
é monótona não decrescente e cresce lentamente; na verdade, cresce cada vez
mais devagar à medida que $n$ aumenta. Quando chegarmos às séries, veremos
que essa sequência não tem cota superior e, portanto, não converge. De modo
algum é evidente que essa sequência não tem cota superior.
\end{example}

Um exemplo comum do aparecimento de sequências monótonas é a proposição a
seguir. Deixamos a demonstração como exercício.

\begin{prop} \label{prop:supinfseq}
Seja $S \subset \R$ um conjunto limitado e não vazio.
Então existem sequências monótonas $\{ x_n \}_{n=1}^\infty$ e
$\{ y_n \}_{n=1}^\infty$ tais que $x_n, y_n \in S$ e
\begin{equation*}
\sup\,S = \lim_{n\to \infty} x_n \qquad \text{e} \qquad \inf\,S =
\lim_{n\to\infty} y_n .
\end{equation*}
\end{prop}

\subsection{Cauda de uma sequência}

\begin{defn}
Para uma sequência $\{ x_n \}_{n=1}^\infty$,
a \emph{cauda-$K$} (em que $K \in \N$), ou simplesmente a
\emph{cauda},\index{cauda de uma sequ\^encia} de
$\{ x_n \}_{n=1}^\infty$ é a sequência que começa em $K+1$, geralmente
escrita como
\begin{equation*}
\{ x_{n+K} \}_{n=1}^\infty
\qquad \text{ou} \qquad \{ x_n \}_{n=K+1}^\infty .
\end{equation*}
\end{defn}

Por exemplo, a cauda-$4$ de $\{ \nicefrac{1}{n} \}_{n=1}^\infty$ é
$\nicefrac{1}{5}, \nicefrac{1}{6}, \nicefrac{1}{7}, \nicefrac{1}{8},
\ldots$. A cauda-$0$ de uma sequência é a própria sequência.
A convergência e o limite de uma sequência dependem apenas de sua cauda.

\begin{prop}
Seja $\{ x_n \}_{n=1}^\infty$ uma sequência. Então, as seguintes
afirmações são equivalentes:
\begin{enumerate}[(i)]
\item \label{prop:ktail:i}
A sequência $\{ x_n \}_{n=1}^\infty$ converge.
\item \label{prop:ktail:ii}
A cauda-$K$ $\{ x_{n+K} \}_{n=1}^\infty$ converge para todo $K \in \N$.
\item \label{prop:ktail:iii}
A cauda-$K$ $\{ x_{n+K} \}_{n=1}^\infty$ converge para algum $K \in \N$.
\end{enumerate}
Além disso, se um desses limites existe (e, portanto, todos), então, para todo $K \in \N$,
\begin{equation*}
\lim_{n\to \infty} x_n = \lim_{n \to \infty} x_{n+K} .
\end{equation*}
\end{prop}

\begin{proof}
É claro que
\ref{prop:ktail:ii} implica \ref{prop:ktail:iii}.
Portanto, mostraremos primeiro que
\ref{prop:ktail:i}
implica
\ref{prop:ktail:ii},
e, em seguida, mostraremos que
\ref{prop:ktail:iii}
implica
\ref{prop:ktail:i}. Isto é,
%mbxSTARTIGNORE
\begin{equation*}
\begin{tikzcd}[row sep=0pt]
& \text{\ref{prop:ktail:ii}} \ar[dr, Rightarrow] & \\
\text{\ref{prop:ktail:i}} \ar[ur, Rightarrow, "\text{a demonstrar}"] & &
\text{\ref{prop:ktail:iii}} \ar[ll, Rightarrow, "\text{a demonstrar}"]
\end{tikzcd}
\end{equation*}
%mbxENDIGNORE
%mbxlatex \begin{center}
%mbxlatex \myincludepdft{tail_of_sequence_cd}{Um diagrama com
%mbxlatex três setas de implicação. A primeira seta, de (i) para (ii),
%mbxlatex está identificada com "a demonstrar". A segunda seta, de (ii) para (iii),
%mbxlatex não está identificada. A terceira seta, de (iii) para (i), está
%mbxlatex identificada com "a demonstrar".}
%mbxlatex \end{center}
No processo, também mostraremos que os limites são iguais.

Comecemos mostrando que \ref{prop:ktail:i} implica \ref{prop:ktail:ii}.
Suponha que $\{x_n \}_{n=1}^\infty$ convirja para algum $x \in \R$.
Seja $K \in \N$ arbitrário e defina $y_n \coloneqq x_{n+K}$. Queremos
mostrar que $\{ y_n \}_{n=1}^\infty$ converge para $x$.
Dado $\epsilon > 0$, existe $M \in \N$ tal que
$\sabs{x-x_n} < \epsilon$ para todo $n \geq M$.
Observe que $n \geq M$ implica $n+K \geq M$. Portanto, para todo
$n \geq M$,
\begin{equation*}
\sabs{x-y_n} = \sabs{x-x_{n+K}} < \epsilon .
\end{equation*}
Consequentemente, $\{ y_n \}_{n=1}^\infty$ converge para $x$.

Passemos a \ref{prop:ktail:iii} implica \ref{prop:ktail:i}.
Seja dado $K \in \N$, defina $y_n \coloneqq x_{n+K}$ e suponha que
$\{ y_n \}_{n=1}^\infty$ convirja para $x \in \R$.
Isto é, dado $\epsilon > 0$, existe $M' \in \N$ tal que
$\sabs{x-y_n} < \epsilon$ para todo $n \geq M'$.
Seja $M \coloneqq M'+K$. Então $n \geq M$ implica $n-K \geq M'$.
Assim, sempre que $n \geq M$,
\begin{equation*}
\sabs{x-x_n} = \sabs{x-y_{n-K}} < \epsilon.
\end{equation*}
Portanto, $\{ x_n \}_{n=1}^\infty$ converge para $x$.
\end{proof}

Em última análise, o limite não depende de como a sequência começa; depende
apenas de sua cauda. O início da sequência pode ser arbitrário.

Por exemplo, a sequência definida por $x_n \coloneqq \frac{n}{n^2+16}$ é
decrescente a partir de $n=4$ (antes disso, é crescente). Isto é:
$\{ x_n \}_{n=1}^\infty =
\nicefrac{1}{17},
\nicefrac{1}{10},
\nicefrac{3}{25},
\nicefrac{1}{8},
\nicefrac{5}{41},
\nicefrac{3}{26},
\nicefrac{7}{65},
\nicefrac{1}{10},
\nicefrac{9}{97},
\nicefrac{5}{58},\ldots$, e
\begin{equation*}
\nicefrac{1}{17} <
\nicefrac{1}{10} <
\nicefrac{3}{25} <
\nicefrac{1}{8} >
\nicefrac{5}{41} >
\nicefrac{3}{26} >
\nicefrac{7}{65} >
\nicefrac{1}{10} >
\nicefrac{9}{97} >
\nicefrac{5}{58} > \ldots .
\end{equation*}
Se descartarmos os três primeiros termos e considerarmos a cauda de índice 3, ela
será decrescente. Deixamos a demonstração como exercício. Como a cauda de índice 3
é monótona e limitada inferiormente por zero, ela é convergente; portanto,
a sequência também é convergente.

\subsection{Subsequ\^encias}

Às vezes, é útil considerar apenas alguns termos de uma sequência.
Uma subsequência de $\{ x_n \}_{n=1}^\infty$ é uma sequência que contém
apenas alguns dos termos de $\{ x_n \}_{n=1}^\infty$, na mesma ordem.

\begin{defn}
Seja $\{ x_n \}_{n=1}^\infty$ uma sequência.
Seja $\{ n_i \}_{i=1}^\infty$ uma sequência estritamente crescente de
números naturais, isto é, $n_i < n_{i+1}$ para todo $i \in \N$ (em outras
palavras, $n_1 < n_2 < n_3 < \cdots$).
A sequência
\begin{equation*}
\{ x_{n_i} \}_{i=1}^\infty
\end{equation*}
é chamada de \glsadd{not:subsequence}
\emph{\myindex{subsequ\^encia}} de $\{ x_n \}_{n=1}^\infty$.
\end{defn}

Assim, a subsequência é a sequência $x_{n_1},x_{n_2},x_{n_3},\ldots$.
Por exemplo, a sequência
$\{ \nicefrac{1}{3i} \}_{i=1}^\infty
= \nicefrac{1}{3},\nicefrac{1}{6},\nicefrac{1}{9},\nicefrac{1}{12},\ldots$
é uma subsequência de
$\{ \nicefrac{1}{n} \}_{n=1}^\infty
= 1, \nicefrac{1}{2},\nicefrac{1}{3},\nicefrac{1}{4},\ldots$.
Para ver como essas duas sequências se encaixam na definição, tome
$n_i \coloneqq 3i$; isto é, $\{ n_i \}_{i=1}^\infty$ é a sequência
$3,6,9,12,\ldots$.
Os números $x_{n_i}$ da subsequência devem vir da sequência original.
Assim, a sequência $1,0,\nicefrac{1}{3},0,
\nicefrac{1}{5},\ldots$ não é uma subsequência de $\{ \nicefrac{1}{n} \}_{n=1}^\infty$.
Além disso, a ordem deve ser preservada: a sequência
$1,\nicefrac{1}{3},\nicefrac{1}{2},\nicefrac{1}{5},\ldots$
não é uma subsequência de $\{ \nicefrac{1}{n} \}_{n=1}^\infty$.

Uma cauda de uma sequência é um tipo especial de subsequência. Para uma
subsequência arbitrária, temos a seguinte proposição sobre convergência,
que é mais fraca do que o resultado para caudas.

\begin{prop} \label{prop:seqtosubseq}
Se $\{ x_n \}_{n=1}^\infty$ é uma sequência convergente, então toda
subsequência $\{ x_{n_i} \}_{i=1}^\infty$ também converge, e
\begin{equation*}
\lim_{n\to \infty} x_n = 
\lim_{i\to \infty} x_{n_i} .
\end{equation*}
\end{prop}

\begin{proof}
Suponha que $\lim_{n\to \infty} x_n = x$. Então, para todo
$\epsilon > 0$, existe $M \in \N$ tal que, para todo $n \geq M$,
\begin{equation*}
\sabs{x_n - x} < \epsilon .
\end{equation*}
Não é difícil provar (faça isso!) por \hyperref[induction:thm]{indução} que
$n_i \geq i$ para todo $i \in \N$. Portanto, $i \geq M$ implica $n_i \geq M$.
Assim, para todo $i \geq M$,
\begin{equation*}
\sabs{x_{n_i} - x} < \epsilon ,
\end{equation*}
e a demonstração está concluída.
\end{proof}

\begin{example}
A existência de uma subsequência convergente não implica que a própria
sequência seja convergente.
Considere a sequência $0,1,0,1,0,1,\ldots$. Isto é,
$x_n = 0$ se $n$ é ímpar, e $x_n = 1$ se $n$ é par. A sequência
$\{ x_n \}_{n=1}^\infty$ é divergente; no entanto, a subsequência
$\{ x_{2i} \}_{i=1}^\infty$ converge para 1 e a subsequência
$\{ x_{2i+1} \}_{i=1}^\infty$ converge para 0. Compare com
\propref{seqconvsubseqconv:prop}.
\end{example}

\subsection{Exerc\'icios}

\begin{exnote}
Nos exercícios a seguir, sinta-se à vontade para usar o que sabe de cálculo
para encontrar o limite, se ele existir. Mas você deve \emph{demonstrar} que
encontrou o limite correto ou provar que a sequência é divergente.
\end{exnote}

\begin{exercise}
A sequência
$\{ 3n \}_{n=1}^\infty$
é limitada? Prove ou refute.
\end{exercise}

\begin{exercise}
A sequência
$\{ n \}_{n=1}^\infty$
é convergente? Se for, qual é o limite?
\end{exercise}

\begin{exercise}
A sequência
$\left\{ \dfrac{{(-1)}^n}{2n} \right\}_{n=1}^\infty$
é convergente? Se for, qual é o limite?
\end{exercise}

\begin{exercise}
A sequência
$\{ 2^{-n} \}_{n=1}^\infty$
é convergente? Se for, qual é o limite?
\end{exercise}

\begin{exercise}
A sequência
$\left\{ \dfrac{n}{n+1} \right\}_{n=1}^\infty$
é convergente? Se for, qual é o limite?
\end{exercise}

\begin{exercise}
A sequência
$\left\{ \dfrac{n}{n^2+1} \right\}_{n=1}^\infty$
é convergente? Se for, qual é o limite?
\end{exercise}

\begin{exercise} \label{exercise:absconv}
Seja $\{ x_n \}_{n=1}^\infty$ uma sequência.
\begin{enumerate}[a)]
\item Mostre que $\lim\limits_{n\to\infty} x_n = 0$ (isto é, que o limite
existe e é zero) se, e somente se, $\lim\limits_{n\to\infty} \sabs{x_n} = 0$.
\item Dê um exemplo em que $\{ \sabs{x_n} \}_{n=1}^\infty$ converge e
$\{ x_n \}_{n=1}^\infty$ diverge.
\end{enumerate}
\end{exercise}

\begin{exercise}
A sequência
$\left\{ \dfrac{2^n}{n!} \right\}_{n=1}^\infty$
é convergente? Se for, qual é o limite?
\end{exercise}

\begin{exercise}
Mostre que a sequência
$\left\{ \dfrac{1}{\sqrt[3]{n}} \right\}_{n=1}^\infty$ é monótona e limitada.
Em seguida, use o \thmref{thm:monotoneconv} para encontrar o limite.
\end{exercise}

\begin{exercise}
Mostre que a sequência
$\left\{ \dfrac{n+1}{n} \right\}_{n=1}^\infty$
é monótona e limitada. Em seguida, use o
\thmref{thm:monotoneconv} para encontrar o limite.
\end{exercise}

\begin{exercise}
Complete a demonstração do \thmref{thm:monotoneconv} para sequências
monótonas não crescentes.
\end{exercise}

\begin{exercise}
Prove a \propref{prop:supinfseq}.
\end{exercise}

\begin{exercise}
Seja $\{ x_n \}_{n=1}^\infty$ uma sequência monótona convergente. Suponha
que exista $k \in \N$ tal que
\begin{equation*}
\lim_{n\to \infty} x_n = x_k .
\end{equation*}
Mostre que $x_n = x_k$ para todo $n \geq k$.
\end{exercise}

\begin{exercise}
\pagebreak[2]
Encontre uma subsequência convergente da sequência
$\bigl\{ {(-1)}^n \bigr\}_{n=1}^\infty$.
\end{exercise}

\begin{exercise}
Seja $\{x_n\}_{n=1}^\infty$ uma sequência definida por
\begin{equation*}
x_n \coloneqq 
\begin{cases}
n               & \text{se } n \text{ é ímpar} , \\
\nicefrac{1}{n} & \text{se } n \text{ é par} .
\end{cases}
\avoidbreak
\end{equation*}
\begin{enumerate}[a)]
\item A sequência é limitada? (prove ou refute)
\item Existe uma subsequência convergente? Se existir, encontre-a.
\end{enumerate}
\end{exercise}

\begin{exercise}
\pagebreak[2]
Seja $\{ x_n \}_{n=1}^\infty$ uma sequência.
Suponha que existam duas subsequências convergentes,
$\{ x_{n_i} \}_{i=1}^\infty$ e $\{ x_{m_i} \}_{i=1}^\infty$. Suponha que
\begin{equation*}
\lim_{i\to\infty} x_{n_i} = a
\qquad \text{e} \qquad
\lim_{i\to\infty} x_{m_i} = b,
\end{equation*}
em que $a \neq b$. Prove que $\{ x_n \}_{n=1}^\infty$ não é convergente,
sem usar a \propref{prop:seqtosubseq}.
\end{exercise}

\begin{exercise}[Desafio]
Encontre uma sequência $\{ x_n \}_{n=1}^\infty$ tal que, para todo
$y \in \R$, exista uma subsequência $\{ x_{n_i} \}_{i=1}^\infty$ que
converge para $y$.
\end{exercise}

\begin{exercise}[F\'acil]
Seja $\{ x_n \}_{n=1}^\infty$ uma sequência e seja $x \in \R$.
Suponha que, para todo $\epsilon > 0$, exista $M$ tal que
$\sabs{x_n-x} \leq \epsilon$ para todo $n \geq M$.
Mostre que $\lim\limits_{n\to\infty} x_n = x$.
\end{exercise}

\begin{exercise}[F\'acil]
Seja $\{ x_n \}_{n=1}^\infty$ uma sequência e seja $x \in \R$ tais que
exista $k \in \N$ para o qual $x_n = x$ para todo $n \geq k$.
Prove que $\{ x_n \}_{n=1}^\infty$ converge para $x$.
\end{exercise}

\begin{exercise}
Seja $\{ x_n \}_{n=1}^\infty$ uma sequência e defina a sequência
$\{ y_n \}_{n=1}^\infty$ por
$y_{2k} \coloneqq x_{k^2}$ e $y_{2k-1} \coloneqq x_k$ para todo $k \in \N$.
Prove que $\{ x_n \}_{n=1}^\infty$ converge se, e somente se,
$\{ y_n \}_{n=1}^\infty$ converge.
Além disso, prove que, se ambas convergem, então
$\lim\limits_{n\to\infty} x_n = \lim\limits_{n\to\infty} y_n$.
\end{exercise}

\begin{exercise}
Mostre que a cauda de índice 3 da sequência definida por
$x_n \coloneqq \frac{n}{n^2+16}$ é monótona não crescente.
Dica: suponha que $n \geq m \geq 4$ e considere o numerador da expressão
$x_n-x_m$.
\end{exercise}

\begin{exercise}
Suponha que $\{ x_n \}_{n=1}^\infty$ seja uma sequência cujas
subsequências $\{ x_{2n} \}_{n=1}^\infty$, $\{ x_{2n-1} \}_{n=1}^\infty$ e
$\{ x_{3n} \}_{n=1}^\infty$ sejam todas convergentes. Mostre que
$\{ x_n \}_{n=1}^\infty$ é convergente.
\end{exercise}

\begin{exercise}
Suponha que $\{ x_n \}_{n=1}^\infty$ seja uma sequência monótona não
decrescente que tenha uma subsequência convergente.
Mostre que $\{ x_n \}_{n=1}^\infty$ é convergente.
Observação: assim, a \propref{prop:seqtosubseq} é uma afirmação do tipo
\myquote{se, e somente se} para sequências monótonas.
\end{exercise}

%%%%%%%%%%%%%%%%%%%%%%%%%%%%%%%%%%%%%%%%%%%%%%%%%%%%%%%%%%%%%%%%%%%%%%%%%%%%%%

\sectionnewpage
\section{Fatos sobre limites de sequências}
\label{sec:factslimsseqs}

%mbxINTROSUBSECTION

\sectionnotes{2--2,5 aulas; sequências definidas recursivamente podem ser
omitidas sem prejuízo}

Nesta seção, apresentamos alguns resultados básicos sobre limites de
sequências.
Começamos examinando como as sequências se relacionam com desigualdades.

\subsection{Limites e desigualdades}

Um lema básico sobre limites e desigualdades é o chamado teorema do confronto.
Ele permite demonstrar a convergência de sequências em casos difíceis
quando encontramos duas outras sequências convergentes mais simples que
\myquote{comprimem} a sequência original.

\begin{lemma}[Teorema do confronto]\index{teorema do confronto} \label{squeeze:lemma}
Sejam $\{ a_n \}_{n=1}^\infty$,
$\{ b_n \}_{n=1}^\infty$ e
$\{ x_n \}_{n=1}^\infty$ sequências tais que
\begin{equation*}
a_n \leq x_n \leq b_n \quad \text{para todo } n \in \N.
\end{equation*}
Suponha que $\{ a_n \}_{n=1}^\infty$ e $\{ b_n \}_{n=1}^\infty$ sejam convergentes e que
\begin{equation*}
\lim_{n\to \infty} a_n
=
\lim_{n\to \infty} b_n .
\end{equation*}
Então $\{ x_n \}_{n=1}^\infty$ é convergente e
\begin{equation*}
\lim_{n\to \infty} x_n
=
\lim_{n\to \infty} a_n
=
\lim_{n\to \infty} b_n .
\end{equation*}
\end{lemma}

\begin{proof}
Seja $x \coloneqq \lim_{n\to\infty} a_n = \lim_{n\to\infty} b_n$.
Seja dado $\epsilon > 0$.
Encontre $M_1$ tal que, para todo $n \geq M_1$, tenhamos
$\sabs{a_n-x} < \epsilon$, e $M_2$ tal que, para todo $n \geq M_2$,
tenhamos $\sabs{b_n-x} < \epsilon$. Defina $M \coloneqq \max \{M_1, M_2 \}$.
Suponha que $n \geq M$.
Em particular,
$x - a_n < \epsilon$, ou seja,
$x - \epsilon < a_n$. De modo análogo, $b_n < x + \epsilon$.
Reunindo essas desigualdades, obtemos
\begin{equation*}
x - \epsilon < a_n \leq x_n \leq b_n < x + \epsilon .
\end{equation*}
Em outras palavras, $-\epsilon < x_n-x < \epsilon$, ou seja,
$\sabs{x_n-x} < \epsilon$.
Portanto, $\{x_n\}_{n=1}^\infty$ converge para $x$.
Veja a \figureref{figsqueeze}.
\end{proof}

\begin{myfigureht}
\myincludepdft{figsqueeze}{%
Diagrama da reta real. Os pontos estão marcados na reta,
da esquerda para a direita: x menos épsilon, a sub n, x, x sub n,
b sub n e x mais épsilon. As setas destacam que a distância
de x menos épsilon até x é épsilon e que a distância
de x até x mais épsilon também é épsilon.}
\caption{Demonstração do teorema do confronto em figura.\label{figsqueeze}}
\end{myfigureht}

\begin{example}
Uma aplicação do
\hyperref[squeeze:lemma]{teorema do confronto} é calcular limites de
sequências usando limites que já conhecemos. Por exemplo, considere
a sequência $\bigl\{ \frac{1}{n\sqrt{n}} \bigr\}_{n=1}^\infty$.
Como $\sqrt{n} \geq 1$ para todo $n \in \N$, temos
\begin{equation*}
0 \leq \frac{1}{n\sqrt{n}} \leq \frac{1}{n}
\end{equation*}
para todo $n \in \N$. Já sabemos que $\lim_{n\to\infty} \nicefrac{1}{n} = 0$.
Logo, usando no teorema do confronto
a sequência constante $\{ 0 \}_{n=1}^\infty$ e a sequência
$\{ \nicefrac{1}{n} \}_{n=1}^\infty$, concluímos que
\begin{equation*}
\lim_{n\to\infty} \frac{1}{n\sqrt{n}} = 0 .
\end{equation*}
\end{example}

Quando existem, os limites preservam desigualdades não estritas.

\begin{lemma} \label{limandineq:lemma}
Sejam $\{ x_n \}_{n=1}^\infty$ e $\{ y_n \}_{n=1}^\infty$
sequências convergentes tais que
\begin{equation*}
x_n \leq y_n \quad \text{para todo } n \in \N .
\end{equation*}
Então
\begin{equation*}
\lim_{n\to\infty} x_n \leq
\lim_{n\to\infty} y_n .
\end{equation*}
\end{lemma}

\begin{proof}
Sejam $x \coloneqq \lim_{n\to\infty} x_n$ e $y \coloneqq \lim_{n\to\infty} y_n$.
Seja dado $\epsilon > 0$. Encontre $M_1$ tal que, para todo $n \geq M_1$,
tenhamos $\sabs{x_n-x} < \nicefrac{\epsilon}{2}$. Encontre $M_2$ tal que,
para todo $n \geq M_2$, tenhamos
$\sabs{y_n-y} < \nicefrac{\epsilon}{2}$. Em particular,
para qualquer $n \geq \max\{ M_1, M_2 \}$, temos
$x-x_n < \nicefrac{\epsilon}{2}$ e
$y_n-y < \nicefrac{\epsilon}{2}$. Somamos essas desigualdades e
obtemos
\begin{equation*}
y_n-x_n+x-y < \epsilon, \qquad \text{ou} \qquad
y_n-x_n < y-x+ \epsilon .
\end{equation*}
Como $x_n \leq y_n$, temos
$0 \leq y_n-x_n$ e, portanto, $0 < y-x+ \epsilon$.
Em outras palavras,
\begin{equation*}
x-y < \epsilon .
\end{equation*}
Como $\epsilon > 0$ era arbitrário, concluímos que
$x-y \leq 0$.
Portanto, $x \leq y$.
\end{proof}

O próximo corolário decorre da aplicação de sequências constantes em
\lemmaref{limandineq:lemma}. A demonstração fica como exercício.

\begin{samepage}
\begin{cor} \label{limandineq:cor}
\leavevmode
\begin{enumerate}[(i)]
\item Se $\{ x_n \}_{n=1}^\infty$ é uma sequência convergente tal que $x_n \geq 0$ para todo
$n \in \N$,
então
\begin{equation*}
\lim_{n\to\infty} x_n \geq 0.
\end{equation*}
\item
Sejam $a,b \in \R$ e
$\{ x_n \}_{n=1}^\infty$ uma sequência convergente tal que
\begin{equation*}
a \leq x_n \leq b \quad \text{para todo } n \in \N .
\end{equation*}
Então
\begin{equation*}
a \leq \lim_{n\to\infty} x_n \leq b.
\end{equation*}
\end{enumerate}
\end{cor}
\end{samepage}

Em \lemmaref{limandineq:lemma} e \corref{limandineq:cor}, não podemos simplesmente substituir
todas as desigualdades não estritas por desigualdades estritas. Por exemplo,
sejam $x_n \coloneqq \nicefrac{-1}{n}$ e $y_n \coloneqq \nicefrac{1}{n}$.
Então $x_n < y_n$, $x_n < 0$ e $y_n > 0$ para todo $n$. No entanto, essas desigualdades
não são preservadas pela operação de limite, pois
$\lim_{n\to\infty} x_n = \lim_{n\to\infty} y_n = 0$.
O ensinamento deste exemplo é que desigualdades estritas podem se tornar não estritas
ao aplicarmos limites; se sabemos que $x_n < y_n$ para todo $n$,
podemos concluir apenas que
\begin{equation*}
\lim_{n \to \infty} x_n \leq
\lim_{n \to \infty} y_n .
\end{equation*}
Esse ponto é uma fonte comum de erros.

\subsection{Continuidade das operações algébricas}

Os limites interagem bem com as operações algébricas.

\begin{prop} \label{prop:contalg}
Sejam $\{ x_n \}_{n=1}^\infty$ e $\{ y_n \}_{n=1}^\infty$ sequências convergentes.
\begin{enumerate}[(i)]
\item \label{prop:contalg:i}
A sequência $\{ z_n \}_{n=1}^\infty$, em que $z_n \coloneqq x_n + y_n$, converge e
\begin{equation*}
\lim_{n \to \infty} (x_n + y_n) =
\lim_{n \to \infty} z_n =
\lim_{n \to \infty} x_n +
\lim_{n \to \infty} y_n .
\end{equation*}
\item \label{prop:contalg:ii}
A sequência $\{ z_n \}_{n=1}^\infty$, em que $z_n \coloneqq x_n - y_n$, converge e
\begin{equation*}
\lim_{n \to \infty} (x_n - y_n) =
\lim_{n \to \infty} z_n =
\lim_{n \to \infty} x_n -
\lim_{n \to \infty} y_n .
\end{equation*}
\item \label{prop:contalg:iii}
A sequência $\{ z_n \}_{n=1}^\infty$, em que $z_n \coloneqq x_n y_n$, converge e
\begin{equation*}
\lim_{n \to \infty} (x_n y_n) =
\lim_{n \to \infty} z_n =
\left( \lim_{n \to \infty} x_n \right)
\left( \lim_{n \to \infty} y_n \right) .
\end{equation*}
\item \label{prop:contalg:iv}
Se $\lim_{n\to\infty} y_n \neq 0$ e $y_n \neq 0$ para todo $n \in \N$, então
a sequência $\{ z_n \}_{n=1}^\infty$, em que $z_n \coloneqq \dfrac{x_n}{y_n}$, converge e
\begin{equation*}
\lim_{n \to \infty} \frac{x_n}{y_n} =
\lim_{n \to \infty} z_n =
\frac{\lim_{n \to \infty} x_n}{\lim_{n \to \infty} y_n} .
\end{equation*}
\end{enumerate}
\end{prop}

\begin{proof}
Comecemos por \ref{prop:contalg:i}.
Suponhamos que $\{ x_n \}_{n=1}^\infty$ e $\{ y_n \}_{n=1}^\infty$ sejam sequências convergentes e
escrevamos $z_n \coloneqq x_n + y_n$. Sejam $x \coloneqq \lim_{n\to\infty} x_n$,
$y \coloneqq \lim_{n\to\infty} y_n$ e $z \coloneqq x+y$.
Seja $\epsilon > 0$ dado.
Encontremos um $M_1$ tal que, para todo $n \geq M_1$,
valha
$\sabs{x_n - x} < \nicefrac{\epsilon}{2}$.
Encontremos um $M_2$ tal que, para todo $n \geq M_2$,
valha
$\sabs{y_n - y} < \nicefrac{\epsilon}{2}$. Seja $M \coloneqq \max \{ M_1, M_2 \}$.
Para todo $n \geq M$, temos
\begin{equation*}
\begin{split}
\sabs{z_n - z} &=
\babs{(x_n+y_n) - (x+y)} \\
& =
\sabs{x_n-x + y_n-y} \\
& \leq
\sabs{x_n-x} + \sabs{y_n-y} \\
& <
\frac{\epsilon}{2} +
\frac{\epsilon}{2}
= \epsilon.
\end{split}
\end{equation*}
Assim, provamos \ref{prop:contalg:i}.
A demonstração de \ref{prop:contalg:ii} é quase idêntica e fica como
exercício.

Passemos a
\ref{prop:contalg:iii}.
Suponhamos novamente que $\{ x_n \}_{n=1}^\infty$ e $\{ y_n \}_{n=1}^\infty$ sejam sequências convergentes e
escrevamos $z_n \coloneqq x_n y_n$. Sejam $x \coloneqq \lim_{n\to\infty} x_n$,
$y \coloneqq \lim_{n\to\infty} y_n$ e $z \coloneqq xy$.
Seja $\epsilon > 0$ dado.
Seja $K \coloneqq \max\{ \sabs{x}, \sabs{y}, \nicefrac{\epsilon}{3} , 1 \}$.
Encontremos um $M_1$ tal que, para todo $n \geq M_1$,
valha
$\sabs{x_n - x} < \frac{\epsilon}{3K}$.
Encontremos um $M_2$ tal que, para todo $n \geq M_2$,
valha
$\sabs{y_n - y} < \frac{\epsilon}{3K}$. Seja $M \coloneqq \max \{ M_1, M_2 \}$.
Para todo $n \geq M$, temos
\begin{equation*}
\begin{split}
\sabs{z_n - z} &=
\babs{(x_ny_n) - (xy)} \\
& =
\babs{(x_n-x+x)(y_n-y+y) - xy} \\
& =
\babs{(x_n-x)y + x(y_n-y) +(x_n-x)(y_n-y)} \\
& \leq
\babs{(x_n-x)y} + \babs{x(y_n - y)} +
\babs{(x_n-x)(y_n-y)} \\
& =
\sabs{x_n -x}\sabs{y} +
\sabs{x}\sabs{y_n -y} +
\sabs{x_n -x}\sabs{y_n -y}
\\
& <
\frac{\epsilon}{3K} K +
K \frac{\epsilon}{3K} +
\frac{\epsilon}{3K}
\frac{\epsilon}{3K}
\qquad \qquad \text{(observe que } \tfrac{\epsilon}{3K} \leq 1
\text{ e }
K \geq 1\text{)}
\\
& \leq
\frac{\epsilon}{3} + \frac{\epsilon}{3} + \frac{\epsilon}{3}
= \epsilon .
\end{split}
\end{equation*}

Por fim, examinemos
\ref{prop:contalg:iv}.
Demonstraremos a seguinte afirmação mais simples:

Afirmação: \emph{Se $\{ y_n \}_{n=1}^\infty$ é uma sequência convergente tal que
$\lim_{n\to\infty} y_n \neq 0$ e $y_n \neq 0$ para todo $n \in \N$, então
$\{ \nicefrac{1}{y_n} \}_{n=1}^\infty$ converge e}
\begin{equation*}
\lim_{n\to\infty} \frac{1}{y_n} = \frac{1}{\lim_{n\to\infty} y_n} .
\end{equation*}

Uma vez demonstrada a afirmação, tomamos a sequência
$\{ \nicefrac{1}{y_n} \}_{n=1}^\infty$,
multiplicamo-la pela sequência $\{ x_n \}_{n=1}^\infty$ e aplicamos o item
\ref{prop:contalg:iii}.

Demonstração da afirmação: Seja $\epsilon > 0$ dado.
Seja $y \coloneqq \lim_{n\to\infty} y_n$.
Como $\sabs{y} \neq 0$, temos
$\min \left\{ \sabs{y}^2\frac{\epsilon}{2}, \, \frac{\sabs{y}}{2} \right\} > 0$.
Encontremos um $M$ tal que, para todo $n \geq M$,
valha
\begin{equation*}
\sabs{y_n - y} < \min \left\{ \sabs{y}^2\frac{\epsilon}{2}, \, \frac{\sabs{y}}{2}
\right\} .
\end{equation*}
Para todo $n \geq M$, temos
$\sabs{y - y_n} < \nicefrac{\sabs{y}}{2}$ e, portanto,
\begin{equation*}
\sabs{y} =
\sabs{y - y_n + y_n } \leq
\sabs{y - y_n} + \sabs{ y_n } < \frac{\sabs{y}}{2} + \sabs{y_n}.
\end{equation*}
Subtraindo $\nicefrac{\sabs{y}}{2}$ de ambos os lados, obtemos
$\nicefrac{\sabs{y}}{2} < \sabs{y_n}$, ou, em outras palavras,
\begin{equation*}
\frac{1}{\sabs{y_n}} < \frac{2}{\sabs{y}} .
\end{equation*}
Concluímos a demonstração da afirmação:
\begin{equation*}
\begin{split}
\abs{\frac{1}{y_n} - \frac{1}{y}} &=
\abs{\frac{y - y_n}{y y_n}} \\
& =
\frac{\sabs{y - y_n}}{\sabs{y} \sabs{y_n}} \\
& \leq
\frac{\sabs{y - y_n}}{\sabs{y}} \, \frac{2}{\sabs{y}} \\
& <
\frac{\sabs{y}^2 \frac{\epsilon}{2}}{\sabs{y}} \, \frac{2}{\sabs{y}}
= \epsilon .
\end{split}
\avoidbreak
\end{equation*}
E a demonstração está concluída.
\end{proof}

Substituindo constantes por sequências, obtemos vários corolários simples.
Se $c \in \R$ e $\{ x_n \}_{n=1}^\infty$ é uma sequência convergente, então,
por exemplo,
\begin{equation*}
\lim_{n \to \infty} c x_n =
c \left( \lim_{n \to \infty} x_n \right) \qquad
\text{e}
\qquad
\lim_{n \to \infty} (c + x_n) =
c + \lim_{n \to \infty} x_n .
\end{equation*}
De modo semelhante, obtemos igualdades para subtração e divisão por constantes.

Como podemos passar limites através da multiplicação, podemos demonstrar (como exercício)
que $\lim_{n\to\infty} x_n^k = {(\lim_{n\to\infty} x_n)}^k$ para todo $k \in \N$.
Isto é, podemos passar limites
por potências. Vejamos se podemos fazer o mesmo com raízes.

\begin{prop}
\pagebreak[2]
Seja $\{ x_n \}_{n=1}^\infty$ uma sequência convergente tal
que $x_n \geq 0$ para todo $n \in \N$.
Então
\begin{equation*}
\lim_{n\to\infty} \sqrt{x_n} =
\sqrt{ \lim_{n\to\infty} x_n } .
\end{equation*}
\end{prop}

É claro que, para formular essa afirmação, precisamos aplicar
\corref{limandineq:cor} e mostrar
que
$\lim_{n\to\infty} x_n \geq 0$, para que possamos calcular a raiz quadrada sem
preocupação.

\begin{proof}
Seja $\{ x_n \}_{n=1}^\infty$ uma sequência convergente e seja $x \coloneqq
\lim_{n\to\infty} x_n$.
Como acabamos de mencionar, $x \geq 0$.

Suponhamos primeiro que $x=0$. Seja $\epsilon > 0$ dado.
Então existe um $M$ tal que, para todo $n \geq M$, vale
$x_n = \sabs{x_n} < \epsilon^2$, ou, em outras palavras, $\sqrt{x_n} < \epsilon$.
Logo,
\begin{equation*}
\abs{\sqrt{x_n} - \sqrt{x}} =
\sqrt{x_n} < \epsilon.
\end{equation*}

Agora suponhamos que $x > 0$ (e, portanto, $\sqrt{x} > 0$).
\begin{equation*}
\begin{split}
\abs{\sqrt{x_n}-\sqrt{x}} &=
\abs{\frac{x_n-x}{\sqrt{x_n}+\sqrt{x}}} \\
&=
\frac{1}{\sqrt{x_n}+\sqrt{x}}
\sabs{x_n-x} \\
& \leq
\frac{1}{\sqrt{x}}
\sabs{x_n-x} .
\end{split}
\end{equation*}
Deixamos o restante da demonstração para o leitor.
\end{proof}

Uma demonstração semelhante vale para a raiz $k$-ésima. Isto é, também
obtemos
$\lim_{n\to\infty} x_n^{1/k} = {( \lim_{n\to\infty} x_n )}^{1/k}$. Deixamos essa demonstração para o leitor
como um exercício desafiador.

Também podemos querer passar o limite através do módulo.
A recíproca desta proposição não vale; veja o
item b) do \exerciseref{exercise:absconv}.

\begin{prop}
Se $\{ x_n \}_{n=1}^\infty$ é uma sequência convergente, então $\{ \sabs{x_n} \}_{n=1}^\infty$
também converge e
\begin{equation*}
\lim_{n\to\infty} \sabs{x_n} =
\abs{\lim_{n\to\infty} x_n} .
\end{equation*}
\end{prop}

\begin{proof}
Basta observar a desigualdade triangular reversa:
\begin{equation*}
\big\lvert \sabs{x_n} - \sabs{x} \big\rvert \leq \sabs{x_n-x} .
\end{equation*}
Portanto, se $\sabs{x_n -x}$ puder ser arbitrariamente pequeno, então
$\big\lvert \sabs{x_n} - \sabs{x} \big\rvert$ também poderá.
Deixamos os detalhes para o leitor.
\end{proof}

Vejamos um exemplo que reúne as proposições acima. Como
$\lim_{n\to\infty} \nicefrac{1}{n} = 0$, temos
\begin{equation*}
\lim_{n\to \infty}
\abs{\sqrt{1 + \nicefrac{1}{n}} - \nicefrac{100}{n^2}} =
\abs{\sqrt{1 + \Bigl(\lim_{n\to\infty} \nicefrac{1}{n}\Bigr)}
- 100 \Bigl(\lim_{n\to\infty} \nicefrac{1}{n}\Bigr)\Bigl(\lim_{n\to\infty} \nicefrac{1}{n}\Bigr)} = 1.
\end{equation*}
Isto é, o limite do lado esquerdo existe porque o limite do lado direito existe. Você realmente deve ler a igualdade acima da direita para a esquerda.

Por outro lado, é preciso aplicar as proposições com cuidado.
Por exemplo, reescrevendo primeiro a expressão com denominador comum,
obtemos
\begin{equation*}
\lim_{n\to \infty} \left( \frac{n^2}{n+1} - n \right)
= -1 .
\end{equation*}
No entanto,
$\bigl\{ \frac{n^2}{n+1} \bigr\}_{n=1}^\infty$ e
$\{n\}_{n=1}^\infty$ não são convergentes;
portanto,
$\Bigl(\lim\limits_{n\to\infty} \frac{n^2}{n+1}\Bigr) -
\Bigl(\lim\limits_{n\to\infty} n\Bigr)$ não faz sentido.

\subsection{Sequências definidas recursivamente}

Agora que sabemos trocar limites e operações algébricas, podemos
calcular os limites de muitas sequências.
Uma dessas classes é a das sequências definidas recursivamente, isto é, sequências em que
o termo seguinte é calculado por uma fórmula que depende de um número fixo
de termos anteriores.

\begin{example}
Seja $\{ x_n \}_{n=1}^\infty$ definida por $x_1 \coloneqq 2$ e
\begin{equation*}
x_{n+1} \coloneqq x_n - \frac{x_n^2-2}{2x_n} .
\end{equation*}
Primeiro, precisamos descobrir se essa sequência está bem definida; precisamos demonstrar que nunca
dividimos por zero.
Depois, precisamos descobrir se a sequência converge. Só então
podemos tentar encontrar seu limite.

Comecemos demonstrando que, para todo $n$,
$x_n$ existe e $x_n > 0$
(assim, a sequência está bem definida e é limitada inferiormente).
Demonstremos isso por \hyperref[induction:thm]{indução}. Sabemos que
$x_1 = 2 > 0$. Para o passo indutivo, suponhamos que $x_n$ exista e
que $x_n > 0$. Então
\begin{equation*}
x_{n+1} = x_n - \frac{x_n^2-2}{2x_n} =
\frac{2x_n^2 - x_n^2+2}{2x_n} =
\frac{x_n^2+2}{2x_n} .
\end{equation*}
Sempre vale $x_n^2+2 > 0$,
e, como
$x_n > 0$, temos $x_{n+1} = \frac{x_n^2+2}{2x_n} > 0$.

Em seguida, mostraremos que a sequência é monótona decrescente. Se demonstrarmos que
$x_n^2-2 \geq 0$ para todo $n$, então $x_{n+1} \leq x_n$ para todo $n$.
É claro que $x_1^2-2 = 4-2 = 2 > 0$. Para um $n$ arbitrário, temos
\begin{equation*}
x_{n+1}^2-2 =
{\left( \frac{x_n^2+2}{2x_n} \right)}^2 - 2
=
\frac{x_n^4+4x_n^2+4 - 8x_n^2}{4x_n^2}
=
\frac{x_n^4-4x_n^2+4}{4x_n^2}
=
\frac{{\left( x_n^2-2 \right)}^2}{4x_n^2} .
\end{equation*}
Como quadrados são não negativos,
$x_{n+1}^2-2 \geq 0$ para todo $n$. Portanto,
$\{ x_n \}_{n=1}^\infty$ é monótona decrescente e limitada (pois $x_n > 0$ para todo $n$); assim,
o limite existe. Resta encontrar esse limite.

Escrevamos
\begin{equation*}
2x_nx_{n+1} = x_n^2+2 .
\end{equation*}
Como $\{ x_{n+1} \}_{n=1}^\infty$ é a cauda de ordem 1 de $\{ x_n \}_{n=1}^\infty$, ela converge para o
mesmo limite. Seja $x \coloneqq \lim_{n\to\infty} x_n$. Tomemos o limite dos
dois lados e obtemos
\begin{equation*}
2x^2 = x^2+2 ,
\end{equation*}
ou $x^2 = 2$. Como $x_n > 0$ para todo $n$, temos $x \geq 0$ e, portanto, $x = \sqrt{2}$.
\end{example}

Talvez você já tenha visto a sequência acima. Trata-se do
\emph{método de Newton}%
\footnote{%
Nomeado em homenagem ao físico e matemático inglês
\href{https://en.wikipedia.org/wiki/Isaac_Newton}{Isaac Newton}
(1642--1726/7).}
para encontrar a raiz quadrada de 2. Esse método aparece com frequência na
prática e converge muito rapidamente. Usamos o fato de que
$x_1^2 -2 >0$, embora isso não fosse estritamente necessário para provar a convergência
considerando uma cauda da sequência.
A sequência converge desde que $x_1 \neq 0$, embora, para um $x_1$
negativo, obtivéssemos $x=-\sqrt{2}$. Substituindo o 2 no numerador,
obtemos a raiz quadrada de qualquer número positivo. Essas afirmações ficam como
exercício.

No entanto, é preciso ter cuidado. Antes de calcular qualquer limite, devemos
verificar se a sequência converge. Vejamos um exemplo.

\begin{example}
Suponhamos que $x_1 \coloneqq 1$ e $x_{n+1} \coloneqq x_n^2+x_n$.
Se presumíssemos, sem verificar, que o limite existe (chamemo-lo de $x$), obteríamos
a equação $x = x^2+x$, da qual poderíamos
concluir $x=0$. No entanto, não é difícil
demonstrar que $\{ x_n \}_{n=1}^\infty$ não é limitada e, portanto, não converge.

O que devemos notar neste exemplo é que o método ainda funciona, mas
depende do valor inicial $x_1$. Se tomarmos $x_1 \coloneqq 0$,
a sequência converge e o limite é de fato 0.
Um ramo inteiro da matemática, chamado de dinâmica, trata precisamente dessas
questões.
Veja
\exerciseref{exercise:convergentinitialvalues}.
\end{example}

\subsection{Alguns critérios de convergência}

Nem sempre é necessário voltar à definição de convergência
para demonstrar que uma sequência converge. Primeiro, apresentamos um critério simples.
A ideia principal é que
$\{ x_n \}_{n=1}^\infty$ converge para $x$ se, e somente se,
$\bigl\{ \sabs{ x_n - x } \bigr\}_{n=1}^\infty$ converge para zero.

\begin{prop} \label{convzero:prop}
Seja $\{ x_n \}_{n=1}^\infty$ uma sequência.
Suponhamos que exista um $x \in \R$
e uma sequência convergente $\{ a_n \}_{n=1}^\infty$
tais que
\begin{equation*}
\lim_{n\to\infty} a_n = 0
\end{equation*}
e
\begin{equation*}
\sabs{x_n - x} \leq a_n
\qquad
\text{para todo $n \in \N$.}
\end{equation*}
Então $\{ x_n \}_{n=1}^\infty$ converge e $\lim\limits_{n\to\infty} x_n = x$.
\end{prop}

\begin{proof}
Seja $\epsilon > 0$ dado. Note que $a_n \geq 0$
para todo $n$. Encontremos um $M \in \N$ tal que, para
todo $n \geq M$, valha
$a_n = \sabs{a_n - 0} < \epsilon$. Então, para todo $n \geq M$,
temos
\begin{equation*}
\sabs{x_n - x} \leq a_n < \epsilon . \qedhere
\end{equation*}
\end{proof}

Em outras palavras, estudar quando uma sequência tem limite é
o mesmo que estudar quando outra sequência converge para zero.
Pode ser difícil decidir se uma dada sequência converge, mas
para algumas sequências há critérios de convergência fáceis de aplicar.
Vejamos um deles.
Primeiro, calcularemos o limite de uma sequência específica.

\begin{prop}
Seja $c > 0$.
\begin{enumerate}[(i)]
\item
Se $c < 1$, então
\begin{equation*}
\lim_{n\to\infty} c^n = 0.
\end{equation*}
\item
Se $c > 1$, então $\{ c^n \}_{n=1}^\infty$ não é limitada.
\end{enumerate}
\end{prop}

\begin{proof}
Consideremos primeiro $c < 1$. Como $c > 0$, obtemos $c^n > 0$ para todo $n \in \N$ por
\hyperref[induction:thm]{indução}. Como $c < 1$, temos
$c^{n+1} < c^n$ para todo $n$.
Assim, a sequência $\{ c^n \}_{n=1}^\infty$ é
limitada inferiormente e
decrescente. Portanto, é convergente.
Seja $x \coloneqq \lim_{n\to\infty} c^n$. A cauda de ordem 1, $\{ c^{n+1} \}_{n=1}^\infty$,
também converge para $x$. Tomando o limite dos dois lados de $c^{n+1} = c \cdot
c^n$, obtemos $x = cx$, ou $(1-c)x=0$. Segue que $x=0$, pois $c \neq 1$.

Agora consideremos $c > 1$.
Seja $B > 0$ arbitrário.
Como $\nicefrac{1}{c} < 1$, a sequência $\bigl\{ {(\nicefrac{1}{c})}^n \bigr\}_{n=1}^\infty$
converge para $0$.
Portanto, para algum $n$ suficientemente grande, temos
\begin{equation*}
\frac{1}{c^n} =
{\left(\frac{1}{c}\right)}^n < \frac{1}{B} .
\end{equation*}
Em outras palavras, $c^n > B$, e $B$ não é cota superior de
$\{ c^n \}_{n=1}^\infty$. Como $B$ era arbitrário, $\{ c^n \}_{n=1}^\infty$ não é limitada.
\end{proof}

Na proposição acima, a
razão entre o $(n+1)$-ésimo e o $n$-ésimo termo é $c$. Generalizamos
esse resultado simples para uma classe maior de sequências.
O lema a seguir voltará a aparecer quando estudarmos séries.

\begin{lemma}[Teste da razão para sequências]\index{teste da razão para sequências}
\label{seq:ratiotest}
Seja $\{ x_n \}_{n=1}^\infty$ uma sequência tal que $x_n \neq 0$ para todo $n$ e tal que o limite
\begin{equation*}
L \coloneqq \lim_{n\to\infty} \frac{\sabs{x_{n+1}}}{\sabs{x_n}}
\qquad
\text{exista.}
\end{equation*}
\begin{enumerate}[(i)]
\item
Se $L < 1$, então $\{ x_n \}_{n=1}^\infty$ converge e $\lim\limits_{n\to\infty} x_n = 0$.
\item
Se $L > 1$, então $\{ x_n \}_{n=1}^\infty$ não é limitada (logo, diverge).
\end{enumerate}
\end{lemma}

Se $L$ existe, mas $L=1$, o lema não fornece nenhuma conclusão.
Não podemos chegar a nenhuma
conclusão com base apenas nessa informação. Por exemplo,
a sequência $\{ \nicefrac{1}{n} \}_{n=1}^\infty$ converge para zero, mas $L=1$.
A sequência constante $\{ 1 \}_{n=1}^\infty$ converge para 1, não para zero, e
$L=1$. A sequência $\bigl\{ {(-1)}^n \bigr\}_{n=1}^\infty$ também não converge, e $L=1$.
Por fim, a sequência $\{ n \}_{n=1}^\infty$ não é limitada, mas, novamente, $L=1$.
O enunciado do lema pode ser um pouco fortalecido; veja os
exercícios \ref{exercise:strongerratiotest1} e
\ref{exercise:strongerratiotest2}.

\begin{proof}
Suponhamos que $L < 1$.
Como
$\frac{\sabs{x_{n+1}}}{\sabs{x_n}} \geq 0$ para todo $n$, temos $L \geq 0$. Escolhamos
$r$ tal que $L < r < 1$.
Queremos comparar a sequência $\{ x_n \}_{n=1}^\infty$ com a sequência $\{ r^n \}_{n=1}^\infty$. A ideia é que,
embora a razão $\frac{\sabs{x_{n+1}}}{\sabs{x_n}}$ não seja menor que $L$ a partir de algum ponto,
ela será eventualmente menor que $r$, que ainda é menor que 1.
A ideia intuitiva da demonstração está ilustrada na \figureref{figratseq}.
\begin{myfigureht}
\myincludepdft{figratseq}{%
Um diagrama da reta real com o ponto L à esquerda, depois
um pouco à direita o ponto r e, ainda mais à direita, o ponto 1.
Traços menores marcam pontos que se agrupam em torno de L e são todos
menores que r.}
\caption{Os traços curtos representam as razões
$\frac{\sabs{x_{n+1}}}{\sabs{x_n}}$
que convergem para $L < 1$.\label{figratseq}}
\end{myfigureht}

Como $r-L > 0$, existe um $M \in \N$ tal que, para
todo $n \geq M$, temos
\begin{equation*}
\abs{\frac{\sabs{x_{n+1}}}{\sabs{x_n}} - L} < r-L .
\end{equation*}
Portanto, para $n \geq M$,
\begin{equation*}
\frac{\sabs{x_{n+1}}}{\sabs{x_n}} - L < r-L
\qquad \text{ou} \qquad
\frac{\sabs{x_{n+1}}}{\sabs{x_n}} < r .
\end{equation*}
Para $n > M$ (isto é, para $n \geq M+1$),
escrevemos
\begin{equation*}
\sabs{x_n} =
\sabs{x_M}
\frac{\sabs{x_{M+1}}}{\sabs{x_{M}}}
\frac{\sabs{x_{M+2}}}{\sabs{x_{M+1}}}
\cdots
\frac{\sabs{x_{n}}}{\sabs{x_{n-1}}}
<
\sabs{x_M}
r r \cdots r = \sabs{x_M} r^{n-M} = (\sabs{x_M} r^{-M}) r^n .
\end{equation*}
A sequência $\{ r^n \}_{n=1}^\infty$ converge para zero; portanto,
$\sabs{x_M} r^{-M} r^n$ também converge para zero. Pela \propref{convzero:prop},
a cauda de ordem $M$,
$\{x_n\}_{n=M+1}^\infty$, converge para zero e, portanto, $\{x_n\}_{n=1}^\infty$ também converge para zero.

Agora suponhamos que $L > 1$. Escolhamos
$r$ tal que $1 < r < L$. Como $L-r > 0$,
existe um $M \in \N$ tal que, para
todo $n \geq M$,
\begin{equation*}
\abs{\frac{\sabs{x_{n+1}}}{\sabs{x_n}} - L} < L-r .
\end{equation*}
Portanto,
\begin{equation*}
\frac{\sabs{x_{n+1}}}{\sabs{x_n}} > r .
\end{equation*}
Novamente, para $n > M$,
escrevemos
\begin{equation*}
\sabs{x_n} =
\sabs{x_M}
\frac{\sabs{x_{M+1}}}{\sabs{x_{M}}}
\frac{\sabs{x_{M+2}}}{\sabs{x_{M+1}}}
\cdots
\frac{\sabs{x_{n}}}{\sabs{x_{n-1}}}
>
\sabs{x_M}
r r \cdots r = \sabs{x_M} r^{n-M} = (\sabs{x_M} r^{-M}) r^n .
\end{equation*}
A sequência $\{ r^n \}_{n=1}^\infty$ não é limitada (pois $r > 1$), logo
$\{x_n\}_{n=1}^\infty$ não pode ser limitada (se $\sabs{x_n} \leq B$ para todo $n$, então
$r^n < \frac{B}{\sabs{x_M}} r^{M}$ para todo $n > M$, o que é impossível).
Consequentemente, $\{ x_n \}_{n=1}^\infty$ não pode convergir.
\end{proof}

\begin{example}
Uma aplicação simples do lema acima é demonstrar que
\begin{equation*}
\lim_{n\to\infty} \frac{2^n}{n!} = 0 .
\end{equation*}

Demonstração:
Calculemos
\begin{equation*}
\frac{2^{n+1} / (n+1)!}{2^n/n!}
=
\frac{2^{n+1}}{2^n}\frac{n!}{(n+1)!}
=
\frac{2}{n+1} .
\end{equation*}
Não é difícil ver que $\bigl\{ \frac{2}{n+1} \bigr\}_{n=1}^\infty$ converge para zero.
A conclusão segue do lema.
\end{example}

\begin{example} \label{example:nto1overn}
Uma aplicação um pouco mais complicada (e útil) do teste da razão é demonstrar que
\begin{equation*}
\lim_{n\to\infty} n^{1/n} = 1 .
\end{equation*}

Demonstração:
Seja $\epsilon > 0$ dado. Consideremos a sequência
$\bigl\{ \frac{n}{{(1+\epsilon)}^n} \bigr\}_{n=1}^\infty$. Calculemos
\begin{equation*}
\frac{(n+1)/{(1+\epsilon)}^{n+1}}{n/{(1+\epsilon)}^{n}}
=
\frac{n+1}{n} \frac{1}{1+\epsilon} .
\end{equation*}
O limite de $\frac{n+1}{n} = 1+\frac{1}{n}$ quando $n \to \infty$ é 1, logo
\begin{equation*}
\lim_{n\to \infty} \frac{(n+1)/{(1+\epsilon)}^{n+1}}{n/{(1+\epsilon)}^{n}}
=
\frac{1}{1+\epsilon} < 1 .
\end{equation*}
Portanto, $\bigl\{ \frac{n}{{(1+\epsilon)}^n} \bigr\}_{n=1}^\infty$ converge para 0.
Em particular,
existe um $M$ tal que, para $n \geq M$, temos
$\frac{n}{{(1+\epsilon)}^n} < 1$, ou
$n < {(1+\epsilon)}^n$, ou
$n^{1/n} < 1+\epsilon$. Como $n \geq 1$, temos $n^{1/n} \geq 1$ e,
portanto, $0 \leq n^{1/n}-1 < \epsilon$. Consequentemente,
$\lim\limits_{n\to\infty} n^{1/n} = 1$.
\end{example}

\subsection{Exercícios}

\begin{exercise}
Demonstre \corref{limandineq:cor}. Dica: use sequências constantes
e \lemmaref{limandineq:lemma}.
\end{exercise}

\begin{exercise}
Demonstre o item \ref{prop:contalg:ii} da \propref{prop:contalg}.
\end{exercise}

\begin{exercise}
Demonstre que, se $\{ x_n \}_{n=1}^\infty$ é uma sequência convergente e $k \in \N$, então
\begin{equation*}
\lim_{n\to\infty} x_n^k =
{\left( \lim_{n\to\infty} x_n \right)}^k .
\end{equation*}
Dica: use \hyperref[induction:thm]{indução}.
\end{exercise}

\begin{exercise}
Suponha que $x_1 \coloneqq \frac{1}{2}$ e $x_{n+1} \coloneqq x_n^2$. Mostre que
$\{ x_n \}_{n=1}^\infty$ converge e encontre
$\lim_{n\to\infty} x_n$. Dica: não se pode dividir por zero!
\end{exercise}

\begin{exercise}
Seja $x_n \coloneqq \frac{n-\cos(n)}{n}$. Use o
\hyperref[squeeze:lemma]{lema do confronto} para mostrar que
$\{ x_n \}_{n=1}^\infty$ converge e determinar seu limite.
\end{exercise}

\begin{exercise}
Sejam $x_n \coloneqq \frac{1}{n^2}$ e $y_n \coloneqq \frac{1}{n}$. Definamos
$z_n \coloneqq \frac{x_n}{y_n}$ e
$w_n \coloneqq \frac{y_n}{x_n}$. As sequências $\{ z_n \}_{n=1}^\infty$ e $\{ w_n \}_{n=1}^\infty$
convergem? Quais são os limites? Podemos aplicar a \propref{prop:contalg}?
Por que sim ou por que não?
\end{exercise}

\begin{exercise}
Verdadeiro ou falso? Demonstre ou encontre um contraexemplo:
Se $\{ x_n \}_{n=1}^\infty$ é uma sequência
tal que $\{ x_n^2 \}_{n=1}^\infty$ converge, então $\{ x_n \}_{n=1}^\infty$ converge.
\end{exercise}

\begin{exercise}
Mostre que
\begin{equation*}
\lim_{n\to\infty} \frac{n^2}{2^n} = 0 .
\end{equation*}
\end{exercise}

\begin{exercise}
Suponha que $\{ x_n \}_{n=1}^\infty$ seja uma sequência, $x \in \R$ e $x_n \neq x$ para
todo $n \in \N$. Suponha que exista o limite
\begin{equation*}
L \coloneqq \lim_{n \to \infty} \frac{\sabs{x_{n+1}-x}}{\sabs{x_n-x}}
\end{equation*}
e que $L < 1$. Mostre que $\{ x_n \}_{n=1}^\infty$ converge para $x$.
\end{exercise}

\begin{exercise}[Desafiador]
Seja $\{ x_n \}_{n=1}^\infty$ uma sequência convergente tal
que $x_n \geq 0$ e $k \in \N$.
Mostre que
\begin{equation*}
\lim_{n\to\infty} x_n^{1/k} =
{\left( \lim_{n\to\infty} x_n \right)}^{1/k} .
\end{equation*}
Dica: Primeiro, trate o caso em que o limite é zero.
Em seguida, encontre uma expressão $q$ tal que
$\frac{x_n^{1/k}-x^{1/k}}{x_n-x} =
\frac{1}{q}$.
\end{exercise}

\begin{exercise}
Seja $r > 0$. Mostre que, partindo de um $x_1 \neq 0$ arbitrário, a sequência
definida por
\begin{equation*}
x_{n+1} \coloneqq x_n - \frac{x_n^2-r}{2x_n}
\end{equation*}
converge para $\sqrt{r}$ se $x_1 > 0$ e para $-\sqrt{r}$ se $x_1 < 0$.
\end{exercise}

\begin{samepage}
\begin{exercise}
Sejam $\{ a_n \}_{n=1}^\infty$ e $\{ b_n \}_{n=1}^\infty$ sequências.
\begin{enumerate}[a)]
\item
Suponha que $\{ a_n \}_{n=1}^\infty$ seja limitada e que $\{ b_n \}_{n=1}^\infty$ convirja para 0.
Mostre que $\{ a_n b_n \}_{n=1}^\infty$ converge para 0.
\item
Dê um exemplo em que $\{ a_n \}_{n=1}^\infty$ não seja limitada, $\{ b_n \}_{n=1}^\infty$ convirja para
0 e $\{ a_n b_n \}_{n=1}^\infty$ não seja convergente.
\item
Dê um exemplo em que $\{ a_n \}_{n=1}^\infty$ seja limitada, $\{ b_n \}_{n=1}^\infty$ convirja para
algum $x \neq 0$ e $\{ a_n b_n \}_{n=1}^\infty$ não seja convergente.
\end{enumerate}
\end{exercise}
\end{samepage}

\begin{samepage}
\begin{exercise}[Fácil] \label{exercise:strongerratiotest1}
Demonstre a seguinte versão mais forte do
\lemmaref{seq:ratiotest}, o teste da razão.
Suponha que $\{ x_n \}_{n=1}^\infty$ seja uma sequência tal que $x_n \neq 0$ para todo
$n$.
\begin{enumerate}[a)]
\item
Demonstre que, se existir um $r < 1$ e um $M \in \N$ tais que
\begin{equation*}
\frac{\sabs{x_{n+1}}}{\sabs{x_n}} \leq r \quad \text{para todo } n \geq M,
\end{equation*}
então $\{ x_n \}_{n=1}^\infty$ converge para $0$.
\item
Demonstre que, se existir um $r > 1$ e um $M \in \N$ tais que
\begin{equation*}
\frac{\sabs{x_{n+1}}}{\sabs{x_n}} \geq r \quad \text{para todo } n \geq M,
\end{equation*}
então $\{ x_n \}_{n=1}^\infty$ não é limitada.
\end{enumerate}
\end{exercise}
\end{samepage}

\begin{exercise} \label{exercise:convergentinitialvalues}
Suponha que $x_1 \coloneqq c$ e $x_{n+1} \coloneqq x_n^2+x_n$.
Mostre que $\{ x_n \}_{n=1}^\infty$ converge se, e somente se, $-1 \leq c \leq 0$; nesse caso,
converge para 0.
\end{exercise}

\begin{exercise}
Demonstre que $\lim\limits_{n \to \infty} {(n^2+1)}^{1/n} = 1$.
\end{exercise}

\begin{exercise}
Demonstre que $\bigl\{ {(n!)}^{1/n} \bigr\}_{n=1}^\infty$
não é limitada.
Dica: mostre que, para todo $C > 0$, a sequência $\left\{ \frac{C^n}{n!}
\right\}_{n=1}^\infty$ converge para zero.
\end{exercise}

%%%%%%%%%%%%%%%%%%%%%%%%%%%%%%%%%%%%%%%%%%%%%%%%%%%%%%%%%%%%%%%%%%%%%%%%%%%%%%

\sectionnewpage

\section{Limite superior, limite inferior e Bolzano--Weierstrass}
\label{sec:bw}

%mbxINTROSUBSECTION

\sectionnotes{1--2 aulas; a demonstração alternativa do teorema de Bolzano--Weierstrass é opcional}

Nesta seção, estudamos sequências limitadas e suas subsequências.
Definimos o que chamamos de limite superior e limite inferior
de uma sequência limitada, que são os maiores e menores
limites possíveis de suas subsequências.
Além disso, demonstramos o
teorema de Bolzano--Weierstrass%
\footnote{%
Nomeado em homenagem ao matemático tcheco
\href{https://en.wikipedia.org/wiki/Bernard_Bolzano}{Bernhard Placidus Johann Nepomuk Bolzano}
(1781--1848) e ao matemático alemão
\href{https://en.wikipedia.org/wiki/Karl_Weierstrass}{Karl Theodor Wilhelm Weierstrass}
(1815--1897).}, uma
ferramenta indispensável em análise, que garante a existência de
subsequências convergentes.

Demonstramos que toda sequência convergente é limitada; ainda assim,
há muitas sequências limitadas que divergem. Por exemplo,
a sequência $\bigl\{ {(-1)}^n \bigr\}_{n=1}^\infty$ é limitada,
mas divergente. No entanto, nem tudo está perdido: ainda podemos
calcular certos limites de sequências limitadas divergentes.

\subsection{Limites superior e inferior}

Há maneiras de construir sequências monótonas a partir de qualquer sequência e,
dessa forma, obtemos o
\emph{\myindex{limite superior}}\index{limsup} e o
\emph{\myindex{limite inferior}}\index{liminf}. Esses limites sempre existem para sequências
limitadas.

Se uma sequência $\{ x_n \}_{n=1}^\infty$ é limitada, então
o conjunto $\{ x_k : k \in \N \}$ é limitado. Para todo $n$,
o conjunto $\{ x_k : k \geq n \}$ também é limitado (pois é um subconjunto); assim,
tomamos seu supremo e seu ínfimo.

\begin{defn} \label{liminflimsup:def}
Seja $\{ x_n \}_{n=1}^\infty$ uma sequência limitada. Definimos as sequências
$\{ a_n \}_{n=1}^\infty$
e $\{ b_n \}_{n=1}^\infty$ por
$a_n \coloneqq \sup \{ x_k : k \geq n \}$ e
$b_n \coloneqq \inf \{ x_k : k \geq n \}$.
Definimos, quando os limites existem,\glsadd{not:limsupseq}\glsadd{not:liminfseq}
\begin{equation*}
\limsup_{n \to \infty} x_n \coloneqq \lim_{n \to \infty} a_n ,
\qquad
\liminf_{n \to \infty} x_n \coloneqq \lim_{n \to \infty} b_n .
\end{equation*}
\end{defn}

Para uma sequência limitada, liminf e limsup sempre existem (veja abaixo). É possível
definir liminf e limsup para sequências não limitadas se permitirmos $\infty$
e $-\infty$, e faremos isso mais adiante nesta seção.
Não é difícil generalizar os resultados a seguir para
incluir sequências não limitadas; primeiro, porém, nos restringimos às
limitadas.

\begin{prop}
Seja $\{ x_n \}_{n=1}^\infty$ uma sequência limitada. Sejam $a_n$ e $b_n$ como na
definição acima.
\begin{enumerate}[(i)]
\item
A sequência $\{ a_n \}_{n=1}^\infty$ é limitada e monótona decrescente,
e $\{ b_n \}_{n=1}^\infty$ é limitada e monótona crescente. Em particular,
$\liminf\limits_{n\to\infty} x_n$ e $\limsup\limits_{n\to\infty} x_n$ existem.
\item
$\displaystyle \limsup_{n \to \infty} x_n = \inf \{ a_n : n \in \N \}$
e
$\displaystyle \liminf_{n \to \infty} x_n = \sup \{ b_n : n \in \N \}$.
\item
$\displaystyle \liminf_{n \to \infty} x_n \leq \limsup_{n \to \infty} x_n$.
\end{enumerate}
\end{prop}

\begin{proof}
Vejamos por que $\{ a_n \}_{n=1}^\infty$ é uma sequência decrescente. Como $a_n$ é a menor cota superior
de $\{ x_k : k \geq n \}$, também é
uma cota superior do subconjunto $\{ x_k : k \geq n+1 \}$. Portanto,
$a_{n+1}$, a menor cota superior de
$\{ x_k : k \geq n+1 \}$, deve ser menor ou igual a $a_n$,
a menor cota superior de
$\{ x_k : k \geq n \}$.
Isto é,
$a_n \geq a_{n+1}$ para todo $n$. De modo semelhante (veja o exercício),
$\{ b_n \}_{n=1}^\infty$ é uma sequência crescente.
Fica como exercício mostrar que,
se $\{ x_n \}_{n=1}^\infty$ é limitada, então $\{ a_n \}_{n=1}^\infty$ e
$\{ b_n \}_{n=1}^\infty$ também são limitadas.

O segundo item decorre do fato de que as sequências
$\{ a_n \}_{n=1}^\infty$ e $\{ b_n \}_{n=1}^\infty$ são monótonas e limitadas.

Para o terceiro item, observe que $b_n \leq a_n$, pois o $\inf$ de um conjunto não vazio
é menor ou igual ao seu $\sup$. As sequências $\{ a_n \}_{n=1}^\infty$
e $\{ b_n \}_{n=1}^\infty$
convergem, respectivamente, para o limsup e o liminf.
Apliquemos a \lemmaref{limandineq:lemma} para obter
\begin{equation*}
\lim_{n\to \infty} b_n \leq \lim_{n\to \infty} a_n.  \qedhere
\end{equation*}
\end{proof}
\begin{myfigureht}
\myincludegraphics{sequence-limsupliminf_an_bn}{%
Um gráfico de uma sequência de pontos x sub n
com 50 termos, isto é, para n de 1 a 50.
No meio, há duas linhas horizontais tracejadas:
a de cima está marcada como lim sup quando n tende ao infinito de x sub n,
e
a de baixo está marcada como lim inf quando n tende ao infinito de x sub n.
Os pontos ficam tanto acima quanto abaixo das linhas, mas parecem
se agrupar cada vez mais perto da região entre elas, à
medida que avançamos para a direita. Duas sequências adicionais estão marcadas:
embaixo, com losangos, há uma sequência menor ou igual
a todos os pontos e constante até alcançar um ponto. Então ela sobe
e permanece constante até alcançar outro ponto, e isso se repete.
Assim, parece haver fileiras horizontais que terminam quando
o losango chega à mesma posição de um ponto, enquanto todos os demais pontos
ficam acima dos losangos.
Todos os losangos estão abaixo da linha tracejada inferior.
Um padrão semelhante, mas invertido, aparece com círculos por cima:
há uma fileira de círculos até que ela alcance um ponto; então desce
e o processo se repete; todos os círculos ficam acima ou sobre os pontos.
Todos os círculos estão acima da linha tracejada superior.}
\caption{Os primeiros 50 termos de uma sequência de exemplo. Os termos $x_n$ da sequência são
marcados com pontos
%mbxSTARTIGNORE
(\raisebox{0.25ex}{\tiny$\bullet$}),
%mbxENDIGNORE
%mbxlatex ($\boldsymbol{\cdot}$),
$a_n$ são marcados com
círculos ($\circ$), e
$b_n$ são marcados com losangos ($\diamond$).\label{sequence-limsupliminf_an_bn}}
\end{myfigureht}

\begin{example} \label{example:liminfsupex}
Seja $\{ x_n \}_{n=1}^\infty$ definida por
\begin{equation*}
x_n \coloneqq
\begin{cases}
\frac{n+1}{n} & \text{se } n \text{ é ímpar,} \\
0             & \text{se } n \text{ é par.}
\end{cases}
\end{equation*}
Calculemos o $\liminf$ e o $\limsup$ dessa sequência. Veja também a
\figureref{sequence-limsupliminf_an_bn-example}. Primeiro, o
limite inferior:
\begin{equation*}
\liminf_{n\to\infty} x_n =
\lim_{n\to\infty}
\bigl(
\inf \{ x_k : k \geq n \}
\bigr)
=
\lim_{n\to\infty} 0 = 0 .
\end{equation*}
Para o limite superior, escrevemos
\begin{equation*}
\limsup_{n\to\infty} x_n =
\lim_{n\to\infty}
\bigl(
\sup \{ x_k : k \geq n \}
\bigr) .
\end{equation*}
Não é difícil ver que
\begin{equation*}
\sup \{ x_k : k \geq n \} =
\begin{cases}
\frac{n+1}{n}   & \text{se } n \text{ é ímpar,} \\
\frac{n+2}{n+1} & \text{se } n \text{ é par.}
\end{cases}
\end{equation*}
Deixamos para o leitor mostrar que o limite é 1. Isto é,
\begin{equation*}
\limsup_{n\to\infty} x_n = 1 .
\end{equation*}
Note que a sequência $\{ x_n \}_{n=1}^\infty$ não é convergente.
\begin{myfigureht}
\myincludegraphics{sequence-limsupliminf_an_bn-example}{%
A mesma marcação da figura anterior, mas agora para uma
sequência específica. A linha tracejada inferior (o lim inf) é o eixo horizontal,
e a linha tracejada superior fica aproximadamente no meio da figura.
O primeiro ponto está bem acima da linha tracejada superior,
e há um círculo na mesma
posição, enquanto o losango está sobre o eixo horizontal.
Na verdade, todos os losangos
estão sobre o eixo horizontal.
O segundo ponto está sobre o eixo horizontal, mas o círculo está acima
da linha tracejada superior e abaixo do primeiro ponto.
O terceiro ponto está na mesma altura que o segundo círculo,
e o terceiro círculo ocupa essa mesma posição.
O padrão se repete: os pontos correspondentes a n ímpar
e os círculos se aproximam cada vez mais da linha tracejada superior
por cima.}
\caption{Os primeiros 20 termos da sequência do \exampleref{example:liminfsupex}.
A marcação é a mesma da \figureref{sequence-limsupliminf_an_bn}.
\label{sequence-limsupliminf_an_bn-example}}
\end{myfigureht}
\end{example}

Associamos certas subsequências a $\limsup$ e $\liminf$.
É importante observar que $\{ a_n \}_{n=1}^\infty$ e
$\{ b_n \}_{n=1}^\infty$ não são
subsequências de $\{ x_n \}_{n=1}^\infty$, nem precisam
ser formadas pelos mesmos números.
Por exemplo, se a sequência é $\{ \nicefrac{1}{n} \}_{n=1}^\infty$, então
$b_n = 0$ para todo $n \in \N$.

\begin{thm} \label{subseqlimsupinf:thm}
Se $\{ x_n \}_{n=1}^\infty$ é uma sequência limitada, então existe uma subsequência
$\{ x_{n_k} \}_{k=1}^\infty$ tal que
\begin{equation*}
\lim_{k\to \infty} x_{n_k} = \limsup_{n \to \infty} x_n .
\end{equation*}
De modo semelhante, existe uma subsequência (possivelmente diferente)
$\{ x_{m_k} \}_{k=1}^\infty$ tal que
\begin{equation*}
\lim_{k\to \infty} x_{m_k} = \liminf_{n \to \infty} x_n .
\end{equation*}
\end{thm}

\begin{proof}
Definamos $a_n \coloneqq \sup \{ x_k : k \geq n \}$.
Seja
$x \coloneqq \limsup_{n\to\infty} x_n = \lim_{n\to\infty} a_n$.
Definimos a subsequência indutivamente.
Seja $n_1 \coloneqq 1$ e suponhamos que $n_1,n_2,\ldots,n_{k-1}$ já estejam definidos para
algum $k \geq 2$. Escolhamos um $m \geq n_{k-1} + 1$
tal que
\begin{equation*}
a_{(n_{k-1}+1)} - x_m < \frac{1}{k} .
\end{equation*}
Esse $m$ existe
porque $a_{(n_{k-1}+1)}$ é o supremo do
conjunto $\{ x_\ell : \ell \geq n_{k-1} + 1 \}$ e, portanto, há termos
da sequência arbitrariamente próximos (ou até iguais) ao supremo.
Definamos $n_{k} \coloneqq m$. Assim, a subsequência $\{ x_{n_k} \}_{k=1}^\infty$ está definida. Agora
precisamos demonstrar que ela converge para $x$.

Para todo $k \geq 2$, temos
$a_{(n_{k-1}+1)} \geq a_{n_k}$ (por quê?) e $a_{n_{k}} \geq x_{n_k}$.
Portanto, para todo $k \geq 2$,
\begin{equation*}
\begin{split}
\sabs{a_{n_k} - x_{n_k}} & =
a_{n_k} - x_{n_k}
\\
& \leq
a_{(n_{k-1}+1)} - x_{n_k}
\\
& < \frac{1}{k} .
\end{split}
\end{equation*}

Mostremos que $\{ x_{n_k} \}_{k=1}^\infty$ converge para $x$.
Observe que a subsequência não precisa ser monótona. Seja $\epsilon > 0$ dado.
Como $\{ a_n \}_{n=1}^\infty$ converge para $x$, a subsequência
$\{ a_{n_k} \}_{k=1}^\infty$ converge para $x$.
Assim, existe um $M_1 \in \N$
tal que, para todo $k \geq M_1$, vale
\begin{equation*}
\sabs{a_{n_k} - x} < \frac{\epsilon}{2} .
\end{equation*}
Encontremos um $M_2 \in \N$ tal que
\begin{equation*}
\frac{1}{M_2} \leq \frac{\epsilon}{2}.
\end{equation*}
Seja $M \coloneqq \max \{M_1 , M_2 , 2 \}$. Para todo $k \geq M$,
\begin{equation*}
\begin{split}
\sabs{x- x_{n_k}} & =
\sabs{a_{n_k} - x_{n_k} + x - a_{n_k}}
\\
& \leq \sabs{a_{n_k} - x_{n_k}} + \sabs{x - a_{n_k}}
\\
& < \frac{1}{k} + \frac{\epsilon}{2}
\\
& \leq \frac{1}{M_2} + \frac{\epsilon}{2} \leq \frac{\epsilon}{2} +
\frac{\epsilon}{2} = \epsilon .
\end{split}
\end{equation*}

Deixamos para o leitor demonstrar a afirmação correspondente ao $\liminf$.
\end{proof}

\subsection{Uso dos limites inferior e superior}

A vantagem de $\liminf$ e $\limsup$ é que podemos sempre escrevê-los
para qualquer sequência (limitada).
Se conseguíssemos calculá-los, também poderíamos calcular o limite da
sequência, caso exista, ou demonstrar que ela diverge.
Trabalhar com $\liminf$ e $\limsup$ é um pouco
como trabalhar com limites, embora haja diferenças sutis.

\begin{prop} \label{liminfsupconv:prop}
Seja $\{ x_n \}_{n=1}^\infty$ uma sequência limitada. Então
$\{ x_n \}_{n=1}^\infty$ converge
se, e somente se,
\begin{equation*}
\liminf_{n\to \infty} x_n =
\limsup_{n\to \infty} x_n.
\end{equation*}
Além disso, se $\{ x_n \}_{n=1}^\infty$ converge, então
\begin{equation*}
\lim_{n\to \infty} x_n =
\liminf_{n\to \infty} x_n =
\limsup_{n\to \infty} x_n.
\end{equation*}
\end{prop}

\begin{proof}
Sejam $a_n$ e $b_n$ como na \defnref{liminflimsup:def}.
Em particular, para todo $n \in \N$,
\begin{equation*}
b_n \leq x_n \leq a_n .
\end{equation*}
Suponhamos primeiro que
$\liminf_{n\to\infty} x_n = \limsup_{n\to\infty} x_n$.
Então $\{ a_n \}_{n=1}^\infty$ e $\{ b_n \}_{n=1}^\infty$
convergem para o mesmo limite.
Pelo lema do confronto
(\lemmaref{squeeze:lemma}), $\{ x_n \}_{n=1}^\infty$ converge e
\begin{equation*}
\lim_{n\to \infty} b_n
=
\lim_{n\to \infty} x_n
=
\lim_{n\to \infty} a_n .
\end{equation*}

Agora suponhamos que $\{ x_n \}_{n=1}^\infty$ convirja para $x$.
Pelo \thmref{subseqlimsupinf:thm},
existe uma subsequência $\{ x_{n_k} \}_{k=1}^\infty$
que converge para $\limsup_{n\to\infty} x_n$.
Como $\{ x_n \}_{n=1}^\infty$ converge para $x$,
toda subsequência converge para $x$ e
portanto $\limsup_{n\to\infty} x_n = \lim_{k\to\infty} x_{n_k} = x$.
De modo semelhante, $\liminf_{n\to\infty} x_n = x$.
\end{proof}

O limite superior e o limite inferior comportam-se bem
com subsequências.

\begin{prop} \label{prop:subseqslimsupinf}
Suponhamos que $\{ x_n \}_{n=1}^\infty$ seja uma sequência limitada e
que $\{ x_{n_k} \}_{k=1}^\infty$ seja uma subsequência. Então
\begin{equation*}
\liminf_{n\to\infty} x_n \leq
\liminf_{k\to\infty} x_{n_k} \leq
\limsup_{k\to\infty} x_{n_k} \leq
\limsup_{n\to\infty} x_n .
\end{equation*}
\end{prop}

\begin{proof}
Já demonstramos a desigualdade do meio. Vamos demonstrar a terceira
desigualdade e deixar a primeira como exercício.

Queremos demonstrar que
$\limsup_{k\to\infty} x_{n_k} \leq \limsup_{n\to\infty} x_n$. Definamos
$a_n \coloneqq \sup \{ x_k : k \geq n \}$
como de costume.
Definamos também
$c_n \coloneqq \sup \{ x_{n_k} : k \geq n \}$.
Não é verdade que $\{ c_n \}_{n=1}^\infty$ seja necessariamente uma subsequência
de $\{ a_n \}_{n=1}^\infty$.
No entanto, como $n_k \geq k$ para todo $k$, temos
$\{ x_{n_k} : k \geq n \} \subset \{ x_k : k \geq n \}$.
O supremo de um subconjunto é menor ou igual ao supremo do
conjunto que o contém; portanto,
\begin{equation*}
c_n \leq a_n \qquad \text{para todo $n$}.
\end{equation*}
A \lemmaref{limandineq:lemma} implica que
\begin{equation*}
\lim_{n\to\infty} c_n \leq \lim_{n\to\infty} a_n ,
\end{equation*}
que é a conclusão desejada.
\end{proof}

O limite superior e o limite inferior
são os maiores e menores
limites subsequenciais. Se a subsequência $\{ x_{n_k} \}_{k=1}^\infty$ da proposição anterior
for convergente, então
$\liminf_{k\to\infty} x_{n_k} = \lim_{k\to\infty} x_{n_k} =
\limsup_{k\to\infty} x_{n_k}$. Portanto,
\begin{equation*}
\liminf_{n\to\infty} x_n \leq
\lim_{k\to\infty} x_{n_k} \leq
\limsup_{n\to\infty} x_n .
\end{equation*}

De modo semelhante, obtemos o seguinte critério útil de convergência
para sequências limitadas. Deixamos a demonstração como exercício.

\begin{prop} \label{seqconvsubseqconv:prop}
Uma sequência limitada $\{ x_n \}_{n=1}^\infty$ converge para $x$
se, e somente se,
toda subsequência convergente
$\{ x_{n_k} \}_{k=1}^\infty$ converge para $x$.
\end{prop}

\subsection{Teorema de Bolzano--Weierstrass}

Embora nem toda sequência limitada seja convergente, o
teorema de Bolzano--Weierstrass nos diz que podemos, ao menos, encontrar uma subsequência
convergente.
A versão do teorema de Bolzano--Weierstrass
apresentada nesta seção é a versão para
sequências de números reais.

\begin{thm}[Bolzano--Weierstrass]\index{teorema de Bolzano--Weierstrass}\label{thm:bwseq}
Suponhamos que uma sequência $\{ x_n \}_{n=1}^\infty$ de números reais seja limitada.
Então existe uma subsequência convergente $\{ x_{n_i} \}_{i=1}^\infty$.
\end{thm}

\begin{proof}
O \thmref{subseqlimsupinf:thm} afirma que existe
uma subsequência cujo limite é $\limsup_{n\to\infty} x_n$.
\end{proof}

Neste momento, o leitor talvez reclame que
o \thmref{subseqlimsupinf:thm} é estritamente mais forte que o
teorema de Bolzano--Weierstrass apresentado acima. Isso é verdade.
No entanto,
o \thmref{subseqlimsupinf:thm} se aplica apenas à reta real, enquanto
o teorema de Bolzano--Weierstrass vale em contextos mais gerais (isto é, em $\R^n$)
com praticamente o mesmo enunciado.

Como esse teorema é tão importante para a análise, apresentamos uma demonstração
explícita.
A ideia da demonstração a seguir também se generaliza para outros contextos.

\begin{proof}[Demonstração alternativa de Bolzano--Weierstrass]
Como a sequência é limitada, existem dois números $a_1 < b_1$
tais que $a_1 \leq x_n \leq b_1$ para todo $n \in \N$.
Definiremos uma subsequência $\{ x_{n_i} \}_{i=1}^\infty$ e duas
sequências $\{ a_i \}_{i=1}^\infty$ e $\{ b_i \}_{i=1}^\infty$ tais que
$\{ a_i \}_{i=1}^\infty$ seja monótona crescente, $\{ b_i \}_{i=1}^\infty$ seja monótona decrescente,
$a_i \leq x_{n_i} \leq b_i$ e $\lim_{i\to\infty} a_i =
\lim_{i\to\infty} b_i$. Então, pelo
\hyperref[squeeze:lemma]{lema do confronto}, a subsequência
$\{ x_{n_i} \}_{i=1}^\infty$ converge.

Definimos as sequências indutivamente. Vamos construí-las de modo que,
para todo $i$, tenhamos $a_i < b_i$
e $x_n \in [a_i,b_i]$ para infinitos $n \in \N$.
Já definimos $a_1$ e $b_1$. Tomemos $n_1 \coloneqq 1$, isto é,
$x_{n_1} = x_1$.
Suponhamos que, até algum $k \in \N$,
tenhamos definido a subsequência $x_{n_1}, x_{n_2}, \ldots,
x_{n_k}$ e as sequências $a_1,a_2,\ldots,a_k$
e $b_1,b_2,\ldots,b_k$.
Seja $y \coloneqq \frac{a_k+b_k}{2}$.
É claro que
$a_k < y < b_k$. Se houver infinitos $j \in \N$
tais que $x_j \in [a_k,y]$, definamos $a_{k+1} \coloneqq a_k$, $b_{k+1}
\coloneqq y$
e escolhamos $n_{k+1} > n_{k}$
tal que $x_{n_{k+1}} \in [a_k,y]$. Se não houver infinitos
$j$ tais que
$x_j \in [a_k,y]$, então haverá infinitos $j \in
\N$ tais que
$x_j \in [y,b_k]$. Nesse caso, definamos $a_{k+1} \coloneqq y$, $b_{k+1}
\coloneqq b_k$
e escolhamos $n_{k+1} > n_{k}$
tal que $x_{n_{k+1}} \in [y,b_k]$.

Assim, as sequências estão definidas. Resta demonstrar que
$\lim_{i\to\infty} a_i = \lim_{i\to\infty} b_i$. Os limites existem, pois as sequências
são monótonas. Na construção,
a diferença $b_i - a_i$ é reduzida pela metade a cada passo. Portanto,
$b_{i+1} - a_{i+1} = \frac{b_i-a_i}{2}$. Por
\hyperref[induction:thm]{indução},
\begin{equation*}
b_i - a_i = \frac{b_1-a_1}{2^{i-1}} .
\end{equation*}

Seja $x \coloneqq \lim_{i\to\infty} a_i$. Como $\{ a_i \}_{i=1}^\infty$ é monótona,
\begin{equation*}
x = \sup \{ a_i : i \in \N \} .
\end{equation*}
Seja $y \coloneqq \lim_{i\to\infty} b_i = \inf \{ b_i : i \in \N \}$. Como $a_i < b_i$ para
todo $i$, temos $x \leq y$.
Como as sequências são monótonas, temos (por quê?) que, para todo $i$,
\begin{equation*}
y-x \leq b_i-a_i = \frac{b_1-a_1}{2^{i-1}} .
\end{equation*}
Como $\frac{b_1-a_1}{2^{i-1}}$ pode ser arbitrariamente pequeno e $y-x \geq 0$,
temos $y-x = 0$. Conclua usando o \hyperref[squeeze:lemma]{lema do confronto}.
\end{proof}

Outra demonstração do teorema de Bolzano--Weierstrass consiste em provar a
seguinte afirmação,
que fica como exercício desafiador:
\emph{Toda sequência tem uma subsequência monótona}.

\subsection{Limites infinitos}

Assim como no caso de ínfimos e supremos, é possível permitir que certos
limites sejam infinitos. Isto é, escrevemos $\lim_{n\to\infty} x_n = \infty$ ou
$\lim_{n\to\infty} x_n = -\infty$ para certas sequências divergentes.

\begin{defn}
\index{limite infinito!de uma sequência}\index{limite!infinito}
Dizemos que
$\{ x_n \}_{n=1}^\infty$ \emph{\myindex{diverge para infinito}}%
\footnote{Às vezes, diz-se que $\{ x_n \}_{n=1}^\infty$ \emph{converge para infinito}.}
se, para todo $K \in
\R$, existe um $M \in \N$ tal que, para todo $n \geq M$, vale $x_n >
K$. Nesse caso, escrevemos
\begin{equation*}
\lim_{n \to \infty} x_n \coloneqq \infty .
\end{equation*}
De modo semelhante,
se, para todo $K \in \R$, existe um $M \in \N$ tal que,
para todo $n \geq M$, vale $x_n < K$, dizemos que $\{ x_n \}_{n=1}^\infty$
\emph{\myindex{diverge para menos infinito}} e escrevemos
\begin{equation*}
\lim_{n \to \infty} x_n \coloneqq -\infty .
\end{equation*}
\end{defn}

Com essa definição e permitindo $\infty$ e $-\infty$,
podemos escrever $\lim_{n\to\infty} x_n$ para qualquer sequência monótona.

\begin{prop} \label{prop:unboundedmonotone}
Suponhamos que $\{ x_n \}_{n=1}^\infty$ seja uma sequência monótona não limitada. Então
\begin{equation*}
\lim_{n \to \infty} x_n =
\begin{cases}
\infty  & \text{se } \{ x_n \}_{n=1}^\infty \text{ é crescente,} \\
-\infty & \text{se } \{ x_n \}_{n=1}^\infty \text{ é decrescente.}
\end{cases}
\end{equation*}
\end{prop}

\begin{proof}
O caso monótono crescente segue do
item c) do \exerciseref{exercise:infseqlimlims} abaixo.
Suponhamos que $\{x_n\}_{n=1}^\infty$ seja decrescente e não limitada.
O fato de a sequência não ser limitada significa que,
para todo $K \in \R$, existe um $M \in \N$ tal que $x_M < K$.
Pela monotonicidade, $x_n \leq x_M < K$ para todo $n \geq M$.
Portanto, $\lim_{n\to\infty} x_n = -\infty$.
\end{proof}

\begin{example}
\begin{equation*}
\lim_{n\to \infty} n = \infty,
\qquad
\lim_{n\to \infty} n^2 = \infty,
\qquad
\lim_{n\to \infty} -n = -\infty.
\end{equation*}
Deixamos a verificação para o leitor.
\end{example}

Também podemos permitir que $\liminf$ e $\limsup$ assumam
os valores $\infty$ e $-\infty$, de modo que
possamos aplicar $\liminf$ e $\limsup$ a qualquer sequência, e não apenas
a uma sequência limitada. Infelizmente, as sequências $\{ a_n \}_{n=1}^\infty$ e
$\{ b_n \}_{n=1}^\infty$
não são sequências de números reais, mas de números reais estendidos. Em
particular, $a_n$ pode ser igual a $\infty$ para algum $n$, e $b_n$ pode ser igual a
$-\infty$. Assim, não temos uma definição para os limites.
Mas, como os reais estendidos ainda formam um conjunto ordenado, podemos
tomar supremos e ínfimos.

\begin{defn}
Seja $\{ x_n \}_{n=1}^\infty$ uma sequência não limitada de números reais. Definimos
sequências de números reais estendidos por
$a_n \coloneqq \sup \{ x_k : k \geq n \}$ e
$b_n \coloneqq \inf \{ x_k : k \geq n \}$.
Definimos\glsadd{not:limsupseq}\glsadd{not:liminfseq}
\index{limite inferior}\index{limite superior}\index{limsup}\index{liminf}
\begin{equation*}
\limsup_{n \to \infty} x_n \coloneqq \inf \{ a_n : n \in \N \}, \qquad \text{e} \qquad
\liminf_{n \to \infty} x_n \coloneqq \sup \{ b_n : n \in \N \}.
\end{equation*}
\end{defn}

Essa definição coincide com a definição para sequências limitadas.

\begin{prop}
Seja $\{ x_n \}_{n=1}^\infty$ uma sequência não limitada. Definamos
$\{ a_n \}_{n=1}^\infty$ e $\{ b_n \}_{n=1}^\infty$ como acima.
Então $\{ a_n \}_{n=1}^\infty$ é decrescente, e $\{ b_n \}_{n=1}^\infty$ é crescente.
Se $a_n$ for um número real para todo $n$, então
$\limsup_{n\to\infty} x_n = \lim_{n\to\infty} a_n$.
Se $b_n$ for um número real para todo $n$, então
$\liminf_{n\to\infty} x_n = \lim_{n\to\infty} b_n$.
\end{prop}

\begin{proof}
Como antes,
$a_n = \sup \{ x_k : k \geq n \} \geq \sup \{ x_k : k \geq n+1 \} =
a_{n+1}$. Portanto, $\{ a_n \}_{n=1}^\infty$ é decrescente. De modo semelhante,
$\{ b_n \}_{n=1}^\infty$ é crescente.

Se $\{ a_n \}_{n=1}^\infty$ for uma sequência de números reais, então
$\lim_{n\to\infty} a_n = \inf \{ a_n : n \in \N \}$. Isso segue do
\thmref{thm:monotoneconv} se $\{ a_n \}_{n=1}^\infty$ for limitada e da
\propref{prop:unboundedmonotone} se $\{a_n \}_{n=1}^\infty$
não for limitada. Procedemos de maneira semelhante com $\{ b_n \}_{n=1}^\infty$.
\end{proof}

A definição se comporta como esperado para
$\limsup$ e $\liminf$; veja os exercícios \ref{exercise:infseqlimex}
e \ref{exercise:infseqlimlims}.

\begin{example}
Suponhamos que
$x_n \coloneqq 0$ para $n$ ímpar e $x_n \coloneqq n$ para $n$ par.
Então $a_n = \infty$ para todo $n$, pois, para todo $M$,
existe um $k$ par tal que $x_k = k \geq M$.
Por outro lado, $b_n = 0$ para todo $n$, pois,
para todo $n$, o conjunto
$\{ x_k : k \geq n \}$ contém $0$ e números positivos.
Assim,
\begin{equation*}
\lim_{n\to \infty} x_n \quad \text{não existe},
\qquad
\limsup_{n\to \infty} x_n = \infty ,
\qquad
\liminf_{n\to \infty} x_n = 0.
\end{equation*}
\end{example}

\subsection{Exercícios}

\begin{exercise}
Suponha que $\{ x_n \}_{n=1}^\infty$ seja uma sequência limitada. Defina $a_n$ e
$b_n$ como na \defnref{liminflimsup:def}. Mostre que $\{ a_n \}_{n=1}^\infty$
e $\{ b_n \}_{n=1}^\infty$ são limitadas.
\end{exercise}

\begin{exercise}
Suponha que $\{ x_n \}_{n=1}^\infty$ seja uma sequência limitada.
Defina $b_n$ como na \defnref{liminflimsup:def}. Mostre que
$\{ b_n \}_{n=1}^\infty$ é uma sequência crescente.
\end{exercise}

\begin{exercise}
Conclua a demonstração da \propref{prop:subseqslimsupinf}. Isto é,
suponha que $\{ x_n \}_{n=1}^\infty$ seja uma sequência limitada e
que $\{ x_{n_k} \}_{k=1}^\infty$ seja uma subsequência. Demonstre que
$\displaystyle \liminf_{n\to\infty} x_n \leq
\liminf_{k\to\infty} x_{n_k}$.
\end{exercise}

\begin{exercise}
Demonstre \propref{seqconvsubseqconv:prop}.
\end{exercise}

\begin{exercise}
\leavevmode
\begin{enumerate}[a)]
\item
Seja $x_n \coloneqq \dfrac{{(-1)}^n}{n}$.
Determine $\limsup\limits_{n\to\infty} x_n$ e $\liminf\limits_{n\to\infty} x_n$.
\item
Seja $x_n \coloneqq \dfrac{(n-1){(-1)}^n}{n}$.
Determine $\limsup\limits_{n\to\infty} x_n$ e $\liminf\limits_{n\to\infty} x_n$.
\end{enumerate}
\end{exercise}

\begin{exercise}
Sejam $\{ x_n \}_{n=1}^\infty$ e $\{ y_n \}_{n=1}^\infty$ sequências limitadas tais que
$x_n \leq y_n$ para todo $n$. Mostre que
\begin{equation*}
\limsup_{n\to\infty} x_n \leq
\limsup_{n\to\infty} y_n
\qquad \text{e} \qquad
\liminf_{n\to\infty} x_n \leq
\liminf_{n\to\infty} y_n .
\end{equation*}
\end{exercise}

\begin{exercise}
Sejam $\{ x_n \}_{n=1}^\infty$ e $\{ y_n \}_{n=1}^\infty$ sequências limitadas.
\begin{enumerate}[a)]
\item
Mostre que $\{ x_n + y_n \}_{n=1}^\infty$ é limitada.
\item
Mostre que
\begin{equation*}
\Bigl(\liminf_{n\to \infty} x_n\Bigr)
+
\Bigl(\liminf_{n\to \infty} y_n\Bigr)
\leq
\liminf_{n\to \infty} \, (x_n+y_n) .
\end{equation*}
Dica: uma demonstração consiste em encontrar uma subsequência $\{ x_{n_m}+y_{n_m} \}_{m=1}^\infty$ de
$\{ x_n + y_n \}_{n=1}^\infty$ que converge.
Em seguida, encontre uma subsequência $\{ x_{n_{m_i}} \}_{i=1}^\infty$ de
$\{ x_{n_m} \}_{m=1}^\infty$ que converge.
\item
Encontre sequências explícitas $\{ x_n \}_{n=1}^\infty$ e $\{ y_n \}_{n=1}^\infty$ tais que
\begin{equation*}
\Bigl(\liminf_{n\to \infty} x_n\Bigr)
+
\Bigl(\liminf_{n\to \infty} y_n\Bigr)
<
\liminf_{n\to \infty} \, (x_n+y_n) .
\end{equation*}
Dica: procure exemplos que não tenham limite.
\end{enumerate}
\end{exercise}

\begin{samepage}
\begin{exercise}
Sejam $\{ x_n \}_{n=1}^\infty$ e $\{ y_n \}_{n=1}^\infty$ sequências limitadas (pelo exercício anterior,
$\{ x_n + y_n \}_{n=1}^\infty$ é limitada).
\begin{enumerate}[a)]
\item
Mostre que
\begin{equation*}
\Bigl(\limsup_{n\to \infty} x_n\Bigr)
+
\Bigl(\limsup_{n\to \infty} y_n\Bigr)
\geq
\limsup_{n\to \infty} \, (x_n+y_n) .
\end{equation*}
Dica: veja o exercício anterior.
\item
Encontre sequências explícitas $\{ x_n \}_{n=1}^\infty$ e $\{ y_n \}_{n=1}^\infty$ tais que
\begin{equation*}
\Bigl(\limsup_{n\to \infty} x_n\Bigr)
+
\Bigl(\limsup_{n\to \infty} y_n\Bigr)
>
\limsup_{n\to \infty} \, (x_n+y_n) .
\end{equation*}
Dica: veja o exercício anterior.
\end{enumerate}
\end{exercise}
\end{samepage}

\begin{exercise}
Se $S \subset \R$ é um conjunto, então $x \in \R$ é um \emph{\myindex{ponto de
acumulação}}
se, para todo $\epsilon > 0$, o conjunto $(x-\epsilon,x+\epsilon) \cap S
\setminus \{ x \}$ não é vazio. Isto é, se há pontos de $S$
arbitrariamente próximos de $x$.
Por exemplo, $S \coloneqq \{ \nicefrac{1}{n} : n \in \N \}$ tem um único ponto de acumulação, $0$, mas $0 \notin S$.
Demonstre a seguinte versão do teorema de Bolzano--Weierstrass:

\medskip

\noindent
\emph{\textbf{Teorema.} Se $S \subset \R$ é um conjunto infinito e limitado,
então existe pelo menos um ponto de acumulação de $S$.}

\medskip

Dica: Se $S$ é infinito, então $S$ contém um subconjunto infinito enumerável.
Isto é, existe uma sequência $\{ x_n \}_{n=1}^\infty$ de números distintos em $S$.
\end{exercise}

\begin{samepage}
\begin{exercise}[Desafiador]
\leavevmode
\begin{enumerate}[a)]
\item
Demonstre que toda sequência contém uma subsequência monótona.
Dica: Chamamos $n \in \N$ de \emph{pico}
da sequência
$\{ x_n \}_{n=1}^\infty$
se $x_m \leq x_n$ para todo $m \geq n$.
Há duas possibilidades: ou a sequência tem apenas um número finito de
picos,
ou tem infinitos picos.
\item
Conclua o teorema de Bolzano--Weierstrass.
\end{enumerate}
\end{exercise}
\end{samepage}

\begin{exercise}
\pagebreak[2]
Demonstre uma versão mais forte da \propref{seqconvsubseqconv:prop}.
Suponha que $\{ x_n \}_{n=1}^\infty$ seja uma sequência tal que toda subsequência
$\{ x_{n_m} \}_{m=1}^\infty$ tenha uma subsequência
$\{ x_{n_{m_i}} \}_{i=1}^\infty$ que convirja para $x$.
\begin{enumerate}[a)]
\item
Primeiro, mostre que $\{ x_n \}_{n=1}^\infty$ é
limitada.
\item
Em seguida, mostre que $\{ x_n \}_{n=1}^\infty$ converge para $x$.
\end{enumerate}
\end{exercise}

\begin{exercise}
Seja $\{x_n\}_{n=1}^\infty$ uma sequência limitada.
\begin{enumerate}[a)]
\item
Demonstre que existe um $s$ tal que, para todo $r > s$, existe
um $M \in \N$ tal que, para todo $n \geq M$, vale
$x_n < r$.
\item
Se $s$ é um número como no item a), demonstre que $\limsup\limits_{n\to\infty} x_n \leq s$.
\item
Seja $S$ o conjunto de todos os valores de $s$ como no item a). Mostre que
$\limsup\limits_{n\to\infty} x_n = \inf \, S$.
\end{enumerate}
\end{exercise}

\begin{exercise}[Fácil] \label{exercise:infseqlimex}
\pagebreak[2]
Suponha que $\{ x_n \}_{n=1}^\infty$ seja tal que $\liminf\limits_{n\to\infty} x_n = -\infty$
e $\limsup\limits_{n\to\infty} x_n = \infty$.
\begin{enumerate}[a)]
\item
Mostre que $\{ x_n \}_{n=1}^\infty$ não converge e que também não é verdade
que $\lim\limits_{n\to\infty} x_n = \infty$
nem que $\lim\limits_{n\to\infty} x_n = -\infty$.
\item
Dê um exemplo de tal sequência.
\end{enumerate}
\end{exercise}

\begin{exercise} \label{exercise:infseqlimlims}
Seja $\{ x_n \}_{n=1}^\infty$ uma sequência.
\begin{enumerate}[a)]
\item
Mostre que
$\lim\limits_{n\to\infty} x_n = \infty$ se, e somente se,
$\liminf\limits_{n\to\infty} x_n = \infty$.
\item
Em seguida, mostre que $\lim\limits_{n\to\infty} x_n = - \infty$ se, e somente se,
$\limsup\limits_{n\to\infty} x_n = -\infty$.
\item
Se $\{ x_n \}_{n=1}^\infty$ é monótona crescente, mostre que ou
$\lim_{n\to\infty} x_n$ existe e é finito, ou $\lim_{n\to\infty} x_n = \infty$. Em qualquer
caso, $\lim_{n\to\infty} x_n = \sup \{ x_n : n \in \N \}$.
\end{enumerate}
\end{exercise}

\begin{exercise} \label{exercise:strongerratiotest2}
Demonstre a seguinte versão mais forte do \lemmaref{seq:ratiotest}, o teste
da razão.
Suponha que $\{ x_n \}_{n=1}^\infty$ seja uma sequência tal que $x_n \neq 0$ para todo
$n$.
\begin{enumerate}[a)]
\item
Demonstre que, se
\begin{equation*}
\limsup_{n \to \infty} \frac{\sabs{x_{n+1}}}{\sabs{x_n}} < 1 ,
\end{equation*}
então $\{ x_n \}_{n=1}^\infty$ converge para $0$.
\item
Demonstre que, se
\begin{equation*}
\liminf_{n \to \infty} \frac{\sabs{x_{n+1}}}{\sabs{x_n}} > 1 ,
\end{equation*}
então $\{ x_n \}_{n=1}^\infty$ não é limitada.
\end{enumerate}
\end{exercise}

\begin{exercise}
Suponha que $\{ x_n \}_{n=1}^\infty$ seja uma sequência limitada e que $a_n \coloneqq \sup \{ x_k : k \geq n \}$
como antes. Suponha que, para algum $\ell \in \N$,
$a_\ell \notin \{ x_k : k \geq \ell \}$. Mostre que $a_j = a_\ell$ para todo $j \geq \ell$ e, portanto,
$\limsup\limits_{n\to\infty} x_n = a_\ell$.
\end{exercise}

\begin{exercise}
Suponha que $\{ x_n \}_{n=1}^\infty$ seja uma sequência,
e $a_n \coloneqq \sup \{ x_k : k \geq n \}$ e
$b_n \coloneqq \inf \{ x_k : k \geq n \}$ como antes.
\begin{enumerate}[a)]
\item
Demonstre que, se $a_\ell = \infty$ para algum $\ell \in \N$, então
$\limsup\limits_{n\to\infty} x_n = \infty$.
\item
Demonstre que, se $b_\ell = -\infty$ para algum $\ell \in \N$, então
$\liminf\limits_{n\to\infty} x_n = -\infty$.
\end{enumerate}
\end{exercise}

\begin{exercise}
Suponha que $\{ x_n \}_{n=1}^\infty$ seja uma sequência
tal que tanto $\liminf_{n\to\infty} x_n$ quanto
$\limsup_{n\to\infty} x_n$ sejam finitos. Demonstre que $\{ x_n \}_{n=1}^\infty$ é limitada.
\end{exercise}

\begin{exercise}
Suponha que $\{ x_n \}_{n=1}^\infty$ seja uma sequência limitada e que $\epsilon > 0$ seja dado.
Demonstre que existe um $M$ tal que, para todo $k \geq M$,
\begin{equation*}
x_k - \Bigl(\limsup_{n\to\infty} x_n\Bigr) < \epsilon \qquad \text{e} \qquad
\Bigl(\liminf_{n\to\infty} x_n\Bigr) - x_k < \epsilon .
\end{equation*}
\end{exercise}

\begin{exercise} \label{exercise:extendsubseqlimsupinf}
Estenda o \thmref{subseqlimsupinf:thm} a sequências não limitadas:
suponha que $\{ x_n \}_{n=1}^\infty$ seja uma sequência.
Se $\limsup_{n\to\infty} x_n = \infty$, demonstre que
existe uma subsequência $\{ x_{n_i} \}_{i=1}^\infty$ que converge para $\infty$.
Em seguida, demonstre o mesmo resultado para $-\infty$ e, depois, demonstre as duas afirmações para
$\liminf$.
\end{exercise}


%%%%%%%%%%%%%%%%%%%%%%%%%%%%%%%%%%%%%%%%%%%%%%%%%%%%%%%%%%%%%%%%%%%%%%%%%%%%%%

\sectionnewpage

\section{Sequências de Cauchy}
\label{sec:cauchy}

\sectionnotes{menos de uma aula}

Muitas vezes, queremos descrever um número por meio de uma sequência que
converge para ele. Nesse caso, não podemos usar o próprio número na demonstração
de que a sequência converge.
Seria conveniente verificar a convergência sem conhecer o limite.

\begin{defn}
Uma sequência $\{ x_n \}_{n=1}^\infty$ é uma \emph{\myindex{sequência de Cauchy}}%
\footnote{%
Nomeada em homenagem ao matemático francês
\href{https://en.wikipedia.org/wiki/Cauchy}{Augustin-Louis Cauchy} (1789--1857).} se, para todo $\epsilon > 0$, existe $M \in \N$ tal que,
para todo $n \geq M$ e todo $k \geq M$,
\begin{equation*}
\sabs{x_n - x_k} < \epsilon .
\end{equation*}
\end{defn}

Intuitivamente, ser de Cauchy significa que, a partir de certo índice, os termos
da sequência ficam arbitrariamente próximos uns dos outros. Podemos esperar que
uma sequência assim seja convergente, e isso de fato ocorre porque $\R$ possui a
\hyperref[defn:lub]{propriedade do supremo}. Antes de demonstrar esse fato,
vejamos alguns exemplos.

\begin{example}
A sequência $\{ \nicefrac{1}{n} \}_{n=1}^\infty$ é de Cauchy.

Demonstração: dado $\epsilon > 0$, encontre $M$ tal que
$M > \nicefrac{2}{\epsilon}$. Então, para $n,k \geq M$,
vale $\nicefrac{1}{n} < \nicefrac{\epsilon}{2}$
e
$\nicefrac{1}{k} < \nicefrac{\epsilon}{2}$. Portanto, para $n,k \geq M$,
\begin{equation*}
\abs{\frac{1}{n} - \frac{1}{k}}
\leq
\abs{\frac{1}{n}} + \abs{\frac{1}{k}}
< \frac{\epsilon}{2} + \frac{\epsilon}{2} = \epsilon.
\end{equation*}
\end{example}

\begin{example}
A sequência $\bigl\{ {(-1)}^n \bigr\}_{n=1}^\infty$ não é de Cauchy.

Demonstração:
Dado $M \in \N$, tome um $n \geq M$ par
e defina $k \coloneqq n+1$. Então
\begin{equation*}
\begin{split}
\babs{{(-1)}^n - {(-1)}^k}
& =
\babs{{(-1)}^n - {(-1)}^{n+1}}
\\
& =
\babs{1 - (-1)} = 2 .
\end{split}
\end{equation*}
Portanto, para qualquer $\epsilon \leq 2$, a definição não pode ser satisfeita e
a sequência não é de Cauchy.
\end{example}

%\begin{example}
%The sequence $\{ \frac{n+1}{n} \}_{n=1}^\infty$ is a Cauchy sequence.
%
%Proof:  Given $\epsilon > 0$, find $M$ such that
%$M > \nicefrac{2}{\epsilon}$.  Then for $n,k \geq M$,
%we have $\nicefrac{1}{n} < \nicefrac{\epsilon}{2}$
%and
%$\nicefrac{1}{k} < \nicefrac{\epsilon}{2}$.  Therefore, for $n, k \geq M$,
%we have
%\begin{equation*}
%\begin{split}
%\abs{\frac{n+1}{n} - \frac{k+1}{k}}
%& =
%\abs{\frac{k(n+1)-n(k+1)}{nk}}
%\\
%& =
%\abs{\frac{kn+k-nk-n}{nk}}
%\\
%& =
%\abs{\frac{k-n}{nk}}
%\\
%& \leq
%\abs{\frac{k}{nk}}
%+
%\abs{\frac{-n}{nk}}
%\\
%& = \frac{1}{n} + \frac{1}{k}
%< \frac{\epsilon}{2} + \frac{\epsilon}{2} = \epsilon .
%\end{split}
%\end{equation*}
%\end{example}

\begin{prop}
Toda sequência de Cauchy é limitada.
\end{prop}

\begin{proof}
Suponha que $\{ x_n \}_{n=1}^\infty$ seja de Cauchy. Escolha $M$ tal que, para
$n,k \geq M$, tenhamos $\sabs{x_n-x_k} < 1$. Em particular,
para todo $n \geq M$,
\begin{equation*}
\sabs{x_n - x_M} < 1 .
\end{equation*}
Pela desigualdade triangular reversa,
$\sabs{x_n} - \sabs{x_M} \leq \sabs{x_n - x_M} < 1$.  Logo, para $n \geq M$,
\begin{equation*}
\sabs{x_n} < 1 + \sabs{x_M}.
\end{equation*}
Defina
\begin{equation*}
B \coloneqq \max \bigl\{ \sabs{x_1}, \sabs{x_2}, \ldots,
\sabs{x_{M-1}}, 1+ \sabs{x_M} \bigr\} .
\end{equation*}
Então $\sabs{x_n} \leq B$ para todo $n \in \N$.
\end{proof}

\begin{thm}
Uma sequência de números reais é de Cauchy se, e somente se, é convergente.
\end{thm}

\begin{proof}
Suponha que $\{ x_n \}_{n=1}^\infty$ converge para $x$
e seja dado $\epsilon > 0$.
Então existe $M$ tal que, para $n \geq M$,
\begin{equation*}
\sabs{x_n - x} < \frac{\epsilon}{2} .
\end{equation*}
Logo, para $n \geq M$ e $k \geq M$,
\begin{equation*}
\sabs{x_n - x_k} = 
\sabs{x_n - x + x - x_k}
\leq \sabs{x_n-x} + \sabs{x-x_k} < \frac{\epsilon}{2} + \frac{\epsilon}{2} =
\epsilon .
\end{equation*}

Essa direção foi simples. Suponha agora que $\{ x_n \}_{n=1}^\infty$ seja de Cauchy.
Já mostramos que $\{ x_n \}_{n=1}^\infty$ é limitada.
Para uma sequência limitada, existem o limite inferior e o limite superior; é aqui
que usamos a
\hyperref[defn:lub]{propriedade do supremo}.
Se mostrarmos que
\begin{equation*}
\liminf_{n\to \infty} x_n = \limsup_{n\to\infty} x_n ,
\end{equation*}
então, pela \propref{liminfsupconv:prop}, $\{ x_n \}_{n=1}^\infty$ será convergente.


Defina $a \coloneqq \limsup_{n\to\infty} x_n$ e
$b \coloneqq \liminf_{n\to\infty} x_n$.
Pelo \thmref{subseqlimsupinf:thm}, existem subsequências
$\{ x_{n_i} \}_{i=1}^\infty$ e
$\{ x_{m_i} \}_{i=1}^\infty$ tais que
\begin{equation*}
\lim_{i\to\infty} x_{n_i} = a
\qquad \text{e} \qquad
\lim_{i\to\infty} x_{m_i} = b.
\end{equation*}
Dado $\epsilon > 0$,
existe $M_1$ tal que
$\sabs{x_{n_i} - a} < \nicefrac{\epsilon}{3}$ para todo $i \geq M_1
e $M_2$ tal que
$\sabs{x_{m_i} - b} < \nicefrac{\epsilon}{3}$ para todo $i \geq M_2$.
Existe ainda $M_3$
tal que
$\sabs{x_n-x_k} < \nicefrac{\epsilon}{3}$
para todo $n,k \geq M_3$.
Seja $M \coloneqq \max \{ M_1, M_2, M_3 \}$.
Se $i \geq M$, então $n_i \geq M$ e $m_i \geq M$. Logo,
\begin{equation*}
\begin{split}
\sabs{a-b} & =
\sabs{a-x_{n_i}+x_{n_i}
-x_{m_i}+x_{m_i}
-b} \\
& \leq
\sabs{a-x_{n_i}}
+ \sabs{x_{n_i} -x_{m_i}}
+ \sabs{x_{m_i} -b} \\
& <
\frac{\epsilon}{3}
+
\frac{\epsilon}{3}
+
\frac{\epsilon}{3}
= \epsilon .
\end{split}
\end{equation*}
Como $\sabs{a-b} < \epsilon$ para todo $\epsilon > 0$, temos $a=b$ e
a sequência converge.
\end{proof}

\begin{remark}
Às vezes, usa-se a afirmação deste teorema para definir a propriedade de
completude dos números reais. Dizemos que um conjunto é
\emph{\myindex{completo de Cauchy}} (ou simplesmente \emph{\myindex{completo}})
se toda sequência de Cauchy converge para um elemento do conjunto.
Acima, provamos que, como $\R$ possui a \hyperref[defn:lub]{propriedade do supremo},
$\R$ é completo de Cauchy.
É possível construir $\R$ \myquote{completando} $\Q$ ao \myquote{acrescentar} pontos
suficientes para que toda sequência de Cauchy convirja (omitimos os detalhes).
O corpo resultante possui a propriedade do supremo.
A vantagem de definir completude por meio de sequências de Cauchy é que essa
noção se generaliza para contextos mais abstratos, como espaços métricos; veja
\chapterref{ms:chapter}.
\end{remark}

O critério de Cauchy é mais forte do que exigir que
$\sabs{x_{n+1}-x_n}$ (ou $\sabs{x_{n+j}-x_n}$, para um $j$ fixo) tenda a zero quando
$n$ tende ao infinito. Ao chegarmos às somas parciais da série harmônica
(veja o \exampleref{example:harmonicseries} na próxima seção), encontraremos
uma sequência tal que $x_{n+1}-x_n = \frac{1}{n+1}$,
mas $\{ x_n \}_{n=1}^\infty$ é divergente.
Na verdade, para essa sequência,
$\lim_{n\to\infty} \sabs{x_{n+j}-x_n} = 0$ para
todo $j \in \N$ (compare com o \exerciseref{exercise:badnocauchy}).
O ponto essencial na definição de Cauchy é que $n$ e $k$
variam independentemente e podem estar arbitrariamente distantes.

\subsection{Exercícios}

\begin{exercise}
Demonstre diretamente pela definição de sequência de Cauchy que $\bigl\{ \frac{n^2-1}{n^2} \bigr\}_{n=1}^\infty$ é de Cauchy.
\end{exercise}

\begin{exercise}
Seja $\{ x_n \}_{n=1}^\infty$ uma sequência tal que existe $C>0$, com $C<1$, e, para todo $n\geq 2$,
\begin{equation*}
\sabs{x_{n+1} - x_n} \leq C \sabs{x_{n}-x_{n-1}} .
\end{equation*}
Demonstre que $\{ x_n \}_{n=1}^\infty$ é de Cauchy.
Dica: você pode usar livremente a fórmula (para $C \neq 1$)
\begin{equation*}
1+ C+ C^2 + \cdots + C^n = \frac{1-C^{n+1}}{1-C}.
\end{equation*}
\end{exercise}

\begin{exercise}[Desafiador]
Suponha que $F$ seja um corpo ordenado que contém os números racionais
$\Q$, e que $\Q$ seja denso em $F$,
isto é: para quaisquer $x,y \in F$ tais que $x < y$,
existe $q \in \Q$ tal que $x < q < y$.
Dizemos que uma sequência $\{ x_n \}_{n=1}^\infty$ de números racionais é de Cauchy
se, para todo $\epsilon \in \Q$ com $\epsilon > 0$, existe
$M$ tal que, para todo $n,k \geq M$, temos $\sabs{x_n-x_k} < \epsilon$.
Suponha que toda sequência de Cauchy de números racionais tenha limite em $F$.
Demonstre que $F$ possui a \hyperref[defn:lub]{propriedade do supremo}.
\end{exercise}

\begin{exercise}
Sejam $\{ x_n \}_{n=1}^\infty$ e $\{ y_n \}_{n=1}^\infty$ sequências tais
que $\lim_{n\to\infty} y_n =0$. Suponha que, para todo $k \in \N$
e todo $m \geq k$, tenhamos
\begin{equation*}
\sabs{x_m-x_k} \leq y_k .
\end{equation*}
Mostre que $\{ x_n \}_{n=1}^\infty$ é de Cauchy.
\end{exercise}

\begin{exercise}
Suponha que uma sequência de Cauchy $\{ x_n \}_{n=1}^\infty$ tenha a seguinte propriedade: para todo $M \in \N$,
existem $k \geq M$ e $n \geq M$ tais que
$x_k < 0$ e $x_n > 0$. Usando apenas as definições de sequência de Cauchy
e de sequência convergente, mostre que
a sequência converge para $0$.
\end{exercise}

\begin{exercise}
Suponha que $\sabs{x_n-x_k} \leq \nicefrac{n}{k^2}$ para todos $n$ e $k$.
Mostre que $\{ x_n \}_{n=1}^\infty$ é de Cauchy.\footnote{Depois de provar que
a sequência é de Cauchy, prove, como exercício extra, que ela é constante.}
\end{exercise}

\begin{exercise}
Suponha que $\{ x_n \}_{n=1}^\infty$ seja uma sequência de Cauchy tal que, para infinitos
$n$, $x_n = c$. Usando apenas a definição de sequência de Cauchy, prove
que $\lim\limits_{n\to\infty} x_n = c$.
\end{exercise}

\begin{exercise}
Verdadeiro ou falso? Demonstre ou dê um contraexemplo: se $\{ x_n \}_{n=1}^\infty$ é uma sequência de Cauchy,
então existe $M$
tal que, para todo $n \geq M$, vale
$\sabs{x_{n+1}-x_n}
\leq
\sabs{x_{n}-x_{n-1}}$.
\end{exercise}

\begin{exercise}
Suponha que
$\{ x_n \}_{n=1}^\infty$ seja uma sequência tal que, para todo $\ell$ e $k$,
vale
$\sabs{x_\ell-x_k} \leq y_\ell + z_k$, em que
$\{ y_n \}_{n=1}^\infty$ e
$\{ z_n \}_{n=1}^\infty$ são sequências que convergem para $0$.
Mostre que
$\{ x_n \}_{n=1}^\infty$ é de Cauchy.
\end{exercise}


%%%%%%%%%%%%%%%%%%%%%%%%%%%%%%%%%%%%%%%%%%%%%%%%%%%%%%%%%%%%%%%%%%%%%%%%%%%%%%

\sectionnewpage

\section{Séries}
\label{sec:series}

%mbxINTROSUBSECTION

\sectionnotes{2 aulas}

Um objeto fundamental da matemática é a série. De fato, durante o desenvolvimento
dos fundamentos da análise, a motivação era compreender as séries. Seu estudo é
importante nas aplicações da análise. Por exemplo, soluções de equações diferenciais
são frequentemente dadas por séries, e as equações diferenciais são a base para
compreender quase toda a ciência moderna.

\subsection{Definição}

\begin{defn}
Dada uma sequência $\{ x_n \}_{n=1}^\infty$, escrevemos o objeto formal
\begin{equation*}
\sum_{n=1}^\infty x_n
\end{equation*}
\glsadd{not:series}
e o chamamos de \emph{\myindex{série}}. Uma série
\emph{converge}\index{convergente!série} se a sequência
$\{ s_k \}_{k=1}^\infty$
definida por
\glsadd{not:finitesum}
\begin{equation*}
s_k \coloneqq \sum_{n=1}^k x_n = x_1 + x_2 + \cdots + x_k
\end{equation*}
converge.
Os números $s_k$ são chamados de
\emph{\myindex{somas parciais}}.
Se a série converge, escrevemos
\begin{equation*}
\sum_{n=1}^\infty x_n =  \lim\limits_{k\to\infty} s_k .
\end{equation*}
Nesse caso, abusamos um pouco da notação e tratamos
$\sum_{n=1}^\infty x_n$ como um número.

Se a sequência $\{ s_k \}_{k=1}^\infty$ diverge,
dizemos que a série é \emph{divergente}\index{divergente!série}.
Nesse caso, $\sum_{n=1}^\infty x_n$ é apenas um objeto formal, e não um número.
\end{defn}

Em outras palavras, para uma série convergente, temos
\begin{equation*}
\sum_{n=1}^\infty x_n
=
\lim_{k\to\infty} 
\sum_{n=1}^k x_n .
\end{equation*}
Essa igualdade só vale se o limite à direita realmente existir. Se a série não
converge, o lado direito não faz sentido (o limite não existe).
Portanto, tenha cuidado: $\sum_{n=1}^\infty x_n$ pode significar duas coisas
diferentes (uma notação para a própria série ou o limite das somas parciais),
e o contexto deve esclarecer qual delas está em uso.

\begin{remark}
Às vezes, é conveniente começar
a série em um índice diferente de 1. Por exemplo, podemos escrever
\begin{equation*}
\sum_{n=0}^\infty r^n = \sum_{n=1}^\infty r^{n-1} .
\end{equation*}
O lado esquerdo é mais conveniente de escrever.
\end{remark}

\begin{remark}
É comum escrever a série $\sum_{n=1}^\infty x_n$ na forma
\begin{equation*}
x_1 + x_2 + x_3 + \cdots
\end{equation*}
entendendo que as reticências indicam uma série,
e não uma soma finita. Não usamos essa notação porque é informal e pode
levar a erros nas demonstrações.
\end{remark}

\begin{example}
A série
\begin{equation*}
\sum_{n=1}^\infty \frac{1}{2^n
\end{equation*}
converge e seu limite é 1. Isto é,
\begin{equation*}
\sum_{n=1}^\infty \frac{1}{2^n} = 
\lim_{k\to\infty} \sum_{n=1}^k \frac{1}{2^n} = 
1 .
\end{equation*}

Demonstração: precisamos da igualdade
\begin{equation*}
\left( \sum_{n=1}^k \frac{1}{2^n} \right)
+ \frac{1}{2^k}
= 1 .
\end{equation*}
A igualdade é imediata para $k=1$. A demonstração para um $k$ arbitrário
segue por \hyperref[induction:thm]{indução}; deixamos os detalhes a cargo do
leitor. A \figureref{figcutseries} ilustra a situação.
\begin{myfigureht}
\myincludepdft{figcutseries}{%
Diagrama da reta real. O ponto 0 está à esquerda e o ponto 1, à direita.
Marca-se a distância de um meio de 0 até um meio. Em seguida, marca-se a distância
de um quarto de um meio até o ponto três quartos. Depois, marca-se a distância
de um oitavo de três quartos até sete oitavos. Também está marcada a distância
um meio mais um quarto mais um oitavo, que reúne as três distâncias mencionadas.
Por fim, marca-se a distância restante de um oitavo entre sete oitavos e 1.}
\caption{A igualdade 
$\left( \sum_{n=1}^k \frac{1}{2^n} \right)
+ \frac{1}{2^k}
= 1$ ilustrada para $k=3$.\label{figcutseries}}
\end{myfigureht}

Seja $s_k$ a soma parcial. Escrevemos
\begin{equation*}
\sabs{
1 - s_k 
}
=
\abs{
1 - 
\sum_{n=1}^k \frac{1}{2^n}
}
=
\abs{\frac{1}{2^k}} = 
\frac{1}{2^k} .
\end{equation*}
A sequência $\bigl\{ \frac{1}{2^k} \bigr\}_{k=1}^\infty$ converge para zero e, portanto,
$\bigl\{ \sabs{1-s_k} \bigr\}_{k=1}^\infty$ também converge para zero.
Logo, $\{ s_k \}_{k=1}^\infty$ converge para $1$.
\end{example}

\begin{prop} \label{geometric:prop}
Suponha que
$-1 < r < 1$. Então a \emph{\myindex{série geométrica}}
$\sum_{n=0}^\infty r^n$ converge, e
\begin{equation*}
\sum_{n=0}^\infty r^n = \frac{1}{1-r} .
\end{equation*}
\end{prop}

Os detalhes da demonstração ficam como exercício.
A demonstração consiste em mostrar 
\begin{equation*}
\sum_{n=0}^{k-1} r^n = \frac{1-r^k}{1-r} ,
\end{equation*}
e então calcular o limite quando $k$ tende a $\infty$.
A série geométrica é uma das séries mais importantes e, de fato, uma das poucas
para as quais podemos encontrar o limite de forma tão explícita.

\medskip

\pagebreak[2]
Assim como para sequências, podemos falar da \emph{\myindex{cauda de uma série}}.

\begin{prop}
Seja $\sum_{n=1}^\infty x_n$ uma série e seja $M \in \N$. Então
\begin{equation*}
\sum_{n=1}^\infty x_n \quad \text{converge se, e somente se,} \quad
\sum_{n=M}^\infty x_n \quad \text{converge.}
\end{equation*}
\end{prop}

\begin{proof}
Consideremos as somas parciais das duas séries (para $k \geq M$)
\begin{equation*}
\sum_{n=1}^{k} x_n
=
\left(
\sum_{n=1}^{M-1} x_n
\right)
+
\sum_{n=M}^{k} x_n .
\end{equation*}
Observe que
$\sum_{n=1}^{M-1} x_n$ é um número fixo. Use
\propref{prop:contalg} para concluir a demonstração.
\end{proof}

\subsection{Séries de Cauchy}

\begin{defn}
Dizemos que uma série $\sum_{n=1}^\infty x_n$ é \emph{de Cauchy} ou uma
\emph{série de Cauchy}\index{série de Cauchy}
se a sequência de somas parciais $\{ s_n \}_{n=1}^\infty$ é uma sequência de Cauchy.
\end{defn}

Uma sequência de números reais é convergente se, e somente se, é
de Cauchy. Portanto, uma série é convergente se, e somente se, é de Cauchy.
A série $\sum_{n=1}^\infty x_n$ é de Cauchy se, e somente se, para todo $\epsilon > 0$,
existe $M \in \N$ tal que, para todo $n \geq M$
e $k \geq M$, vale
\begin{equation*}
\abs{ \left( \sum_{i=1}^k x_i \right) - \left( \sum_{i=1}^n x_i \right) }
< \epsilon .
\end{equation*}
Sem perda de generalidade, suponhamos que $n < k$. Então escrevemos
\begin{equation*}
\abs{ \left( \sum_{i=1}^k x_i \right) - \left( \sum_{i=1}^n x_i \right) }
=
\abs{ \sum_{i={n+1}}^k x_i }
< \epsilon .
\end{equation*}
Provamos a seguinte proposição simples.

\begin{prop} \label{prop:cachyser}
A série $\sum_{n=1}^\infty x_n$ é de Cauchy
se, e somente se, para todo $\epsilon > 0$,
existe $M \in \N$ tal que, para todo $n \geq M$
e todo $k > n$,
\begin{equation*}
\abs{ \sum_{i={n+1}}^k x_i }
< \epsilon .
\end{equation*}
\end{prop}

\subsection{Propriedades básicas}

\begin{prop}
Seja $\sum_{n=1}^\infty x_n$ uma série convergente. Então
a sequência $\{ x_n \}_{n=1}^\infty$ é convergente e
\begin{equation*}
\lim_{n\to\infty} x_n = 0.
\end{equation*}
\end{prop}

\begin{proof}
Seja dado $\epsilon > 0$. Como $\sum_{n=1}^\infty x_n$ converge, ela é de Cauchy.
Assim, encontramos $M$ tal que, para todo $n \geq M$,
\begin{equation*}
\epsilon > 
\abs{ \sum_{i={n+1}}^{n+1} x_i }
=
\sabs{ x_{n+1} } .
\end{equation*}
Logo, para todo $n \geq M+1$, temos $\sabs{x_{n}} < \epsilon$.
\end{proof}

\begin{example}
Se $r \geq 1$ ou $r \leq -1$, então a série geométrica $\sum_{n=0}^\infty r^n$
diverge.

Demonstração: $\sabs{r^n} = \sabs{r}^n \geq 1^n = 1$. Os termos não tendem a zero
e a série não pode convergir.
\end{example}

Portanto, se uma série converge, seus termos tendem a zero.
A implicação, porém, vale apenas em um sentido.
Vejamos um exemplo.

\begin{example} \label{example:harmonicseries}
A série $\sum_{n=1}^\infty \nicefrac{1}{n}$ diverge (embora
$\lim_{n\to\infty} \nicefrac{1}{n} = 0$).
Esta é a famosa \emph{\myindex{série harmônica}}%
\footnote{A divergência da série harmônica era conhecida muito
antes de a teoria das séries ser formalizada. A demonstração que
apresentamos é a mais antiga e foi dada por
\href{https://en.wikipedia.org/wiki/Oresme}{Nicole Oresme}
(1323?--1382).}.

Demonstração: mostraremos que a sequência das somas parciais é ilimitada e, portanto,
não pode convergir.
Escrevamos as somas parciais $s_n$ para $n = 2^k$ na forma:
\begin{align*}
 s_1 & = 1 , \\
 s_2 & = \left( 1 \right) + \left( \frac{1}{2} \right) , \\
 s_4 & = \left( 1 \right) + \left( \frac{1}{2} \right) +
        \left( \frac{1}{3} + \frac{1}{4} \right) , \\
 s_8 & = \left( 1 \right) + \left( \frac{1}{2} \right) +
        \left( \frac{1}{3} + \frac{1}{4} \right) +
        \left( \frac{1}{5} + \frac{1}{6} + \frac{1}{7} + \frac{1}{8} \right) , \\
& ~~ \vdots \\
 s_{2^k} & = 
1 + 
\sum_{i=1}^k
\left(
\sum_{m=2^{i-1}+1}^{2^i} \frac{1}{m}
\right) .
\end{align*}
Observe que $\nicefrac{1}{3} + \nicefrac{1}{4} \geq \nicefrac{1}{4} + \nicefrac{1}{4} =
\nicefrac{1}{2}$ e
$\nicefrac{1}{5} + \nicefrac{1}{6} + \nicefrac{1}{7} + \nicefrac{1}{8}
\geq \nicefrac{1}{8} + \nicefrac{1}{8} + \nicefrac{1}{8} + \nicefrac{1}{8} =
\nicefrac{1}{2}$. Mais geralmente
\begin{equation*}
\sum_{m=2^{k-1}+1}^{2^k} \frac{1}{m}
\geq
\sum_{m=2^{k-1}+1}^{2^k} \frac{1}{2^k}
=
(2^{k-1}) \frac{1}{2^k} = \frac{1}{2} .
\end{equation*}
Portanto,
\begin{equation*}
s_{2^k} = 
1 + 
\sum_{i=1}^k
\left(
\sum_{m=2^{i-1}+1}^{2^i} \frac{1}{m}
\right) 
\geq
1 + \sum_{i=1}^k \frac{1}{2} = 1 + \frac{k}{2} .
\end{equation*}
Como $\{ \nicefrac{k}{2} \}_{k=1}^\infty$ é ilimitada pela
\hyperref[thm:arch:i]{propriedade arquimediana},
$\{ s_{2^k} \}_{k=1}^\infty$ é ilimitada
e, portanto, $\{ s_n \}_{n=1}^\infty$ também é ilimitada.
Logo, $\{ s_n \}_{n=1}^\infty$ diverge e, consequentemente, $\sum_{n=1}^\infty \nicefrac{1}{n}$ diverge.
\end{example}

Assim como as somas finitas, as séries convergentes obedecem às propriedades de linearidade.
Isto é, podemos multiplicá-las por constantes
e somá-las, realizando essas operações termo a termo.

\begin{prop}[Linearidade das séries]\index{linearidade!de séries}
Sejam $\alpha \in \R$ e $\sum_{n=1}^\infty x_n$, $\sum_{n=1}^\infty y_n$
séries convergentes. Então
\begin{enumerate}[(i)]
\item
$\sum_{n=1}^\infty \alpha x_n$ é uma série convergente e
\begin{equation*}
\sum_{n=1}^\infty \alpha x_n
=
\alpha \sum_{n=1}^\infty x_n .
\end{equation*}
\item
$\sum_{n=1}^\infty ( x_n + y_n )$ é uma série convergente e
\begin{equation*}
\sum_{n=1}^\infty ( x_n + y_n ) 
=
\left( \sum_{n=1}^\infty x_n \right)
+
\left( \sum_{n=1}^\infty y_n \right) .
\end{equation*}
\end{enumerate}
\end{prop}

\begin{proof}
Para o primeiro item,
escrevemos a $k$-ésima soma parcial
\begin{equation*}
\sum_{n=1}^k \alpha x_n
=
\alpha \left( \sum_{n=1}^k x_n \right) .
\end{equation*}
Observamos que o lado direito é um múltiplo constante de
uma sequência convergente e, portanto, é convergente. Assim, calculamos o limite dos dois lados e obtemos
o resultado.

Para o segundo item, também consideramos a
$k$-ésima soma parcial
\begin{equation*}
\sum_{n=1}^k ( x_n + y_n ) 
=
\left( \sum_{n=1}^k x_n \right)
+
\left( \sum_{n=1}^k y_n \right) .
\end{equation*}
Observamos que o lado direito é a soma de sequências convergentes
e, portanto, é convergente. Assim, calculamos o limite dos dois lados e obtemos
a proposição.
\end{proof}

Uma aplicação útil do primeiro item é a seguinte
fórmula. Se $\sabs{r} < 1$ e $i \in \N$, então
\begin{equation*}
\sum_{n=i}^\infty r^n = \frac{r^i}{1-r} .
\end{equation*}
A fórmula decorre da série geométrica e da multiplicação por
$r^i$:
\begin{equation*}
r^i \sum_{n=0}^\infty r^n =
\sum_{n=0}^\infty r^{n+i}
=
\sum_{n=i}^\infty r^n .
\end{equation*}

Multiplicar séries não é tão simples quanto somá-las; veja a próxima
seção.
Naturalmente, não podemos multiplicar
termo a termo. Essa estratégia nem sequer funciona para somas finitas:
$(a+b)(c+d) \neq ac+bd$.

\subsection{Convergência absoluta}

Como é mais fácil trabalhar com sequências monótonas do que com sequências arbitrárias,
em geral é mais fácil trabalhar com séries $\sum_{n=1}^\infty x_n$ em que $x_n \geq 0$ para
todo $n$. Nesse caso, a sequência das somas parciais é monótona não decrescente
e converge se for limitada superiormente.
Formalizemos essa afirmação em uma proposição.

\begin{prop}
Se $x_n \geq 0$ para todo $n$, então $\sum_{n=1}^\infty x_n$ converge se, e somente se,
a sequência das somas parciais é limitada superiormente.
\end{prop}

O limite de uma sequência monótona não decrescente é seu supremo.
Assim, quando $x_n \geq 0$ para todo $n$, vale a
desigualdade
\begin{equation*}
\sum_{n=1}^k x_n \leq
\sum_{n=1}^\infty x_n .
\end{equation*}
Se admitirmos limites infinitos, a desigualdade ainda
vale quando a série diverge para infinito, embora, nesse caso, não seja
muito útil.

Veremos que o seguinte critério comum de convergência de séries
tem consequências importantes para as operações que podemos realizar com elas.

\begin{defn}
Uma série $\sum_{n=1}^\infty x_n$
\emph{\myindex{converge absolutamente}}\index{convergência absoluta} se
a série $\sum_{n=1}^\infty \sabs{x_n}$ converge.
Se uma série converge, mas não converge absolutamente, dizemos
que ela \emph{\myindex{converge condicionalmente}}.\index{convergência condicional}
\end{defn}

\begin{prop}
Se a série $\sum_{n=1}^\infty x_n$ converge absolutamente, então ela converge.
\end{prop}

\begin{proof}
Uma série converge se, e somente se, é de Cauchy.
Suponha que $\sum_{n=1}^\infty \sabs{x_n}$ seja de Cauchy. Isto é, para todo $\epsilon > 0$,
existe $M$ tal que, para todo $k \geq M$ e todo $n > k, vale 
\begin{equation*}
\sum_{i=k+1}^n \sabs{x_i} 
=
\abs{ \sum_{i=k+1}^n \sabs{x_i} }
<
\epsilon .
\end{equation*}
Aplicamos a desigualdade triangular para somas finitas e obtemos
\begin{equation*}
\abs{ \sum_{i=k+1}^n x_i }
\leq
\sum_{i=k+1}^n \sabs{x_i}
<
\epsilon .
\end{equation*}
Logo, $\sum_{n=1}^\infty x_n$ é de Cauchy e, portanto, converge.
\end{proof}

Se $\sum_{n=1}^\infty x_n$ converge absolutamente, em geral os limites de
$\sum_{n=1}^\infty x_n$ e $\sum_{n=1}^\infty \sabs{x_n}$ são diferentes. Calcular um
deles não ajuda a calcular o outro. No entanto, o cálculo acima conduz a
uma desigualdade útil para séries absolutamente convergentes:
uma versão da desigualdade triangular para séries,
cuja demonstração fica como exercício:
\begin{equation*}
\abs{ \sum_{i=1}^\infty x_i }
\leq
\sum_{i=1}^\infty \sabs{x_i} .
\end{equation*}

As séries absolutamente convergentes têm muitas propriedades importantes.
Por exemplo, podemos reordenar seus termos arbitrariamente ou multiplicar essas
séries com facilidade. As séries condicionalmente convergentes,
por outro lado,
muitas vezes não se comportam como esperaríamos. Veja a próxima seção.

Deixamos como exercício mostrar que
\begin{equation*}
\sum_{n=1}^\infty \frac{{(-1)}^n}{n}
\end{equation*}
converge, embora o leitor deva terminar esta seção antes de tentar.
Por outro lado, provamos que
\begin{equation*}
\sum_{n=1}^\infty \frac{1}{n}
\end{equation*}
diverge. Portanto,
$\sum_{n=1}^\infty \frac{{(-1)}^n}{n}$ é uma série condicionalmente convergente.

\subsection{Critério da comparação e série \texorpdfstring{$p$}{p}}

Observamos acima que, para uma série de termos positivos convergir,
não basta que os termos tendam a zero: eles precisam tender a zero
\myquote{rápido o suficiente}. Se conhecemos a convergência de uma série,
podemos usar o seguinte critério da comparação para verificar se os termos de outra
série tendem a zero \myquote{rápido o suficiente}.

\begin{samepage}
\begin{prop}[Critério da comparação]\index{critério da comparação para séries}
Sejam $\sum_{n=1}^\infty x_n$ e $\sum_{n=1}^\infty y_n$ séries tais que $0 \leq x_n \leq y_n$
para todo $n \in \N$.
\begin{enumerate}[(i)]
\item Se $\sum_{n=1}^\infty y_n$ converge, então $\sum_{n=1}^\infty x_n$ também converge.
\item Se $\sum_{n=1}^\infty x_n$ diverge, então $\sum_{n=1}^\infty y_n$ também diverge.
\end{enumerate}
\end{prop}
\end{samepage}

\begin{proof}
Como todos os termos das séries são não negativos, as sequências das
somas parciais são monótonas não decrescentes.
Como $x_n \leq y_n$ para todo $n$, as somas parciais satisfazem, para todo $k$,
\begin{equation} \label{comptest:eq}
\sum_{n=1}^k x_n \leq \sum_{n=1}^k y_n .
\end{equation}
Se a série $\sum_{n=1}^\infty y_n$ converge, suas somas parciais são
limitadas. Portanto, o lado direito de \eqref{comptest:eq}
é limitado para todo $k$; existe $B \in \R$ tal que
$\sum_{n=1}^k y_n \leq B$ para todo $k$, e assim
\begin{equation*}
\sum_{n=1}^k x_n \leq \sum_{n=1}^k y_n \leq B.
\end{equation*}
Logo, as somas parciais de $\sum_{n=1}^\infty x_n$
também são limitadas. Como formam uma sequência monótona não decrescente,
convergem. Assim, provamos o primeiro item.

Por outro lado, se $\sum_{n=1}^\infty x_n$ diverge, a sequência das somas parciais
é ilimitada, pois é monótona não decrescente. Isto é, as somas
parciais de $\sum_{n=1}^\infty x_n$ acabam superando qualquer número real. Juntando isso
com \eqref{comptest:eq}, vemos que, para todo $B \in
\R$, existe $k$ tal que 
\begin{equation*}
B \leq \sum_{n=1}^k x_n \leq \sum_{n=1}^k y_n .
\end{equation*}
Logo, as somas parciais de $\sum_{n=1}^\infty y_n$ também são ilimitadas, e
$\sum_{n=1}^\infty y_n$ também diverge.
\end{proof}

Uma série útil para aplicar o critério da comparação é a
série $p$%
\footnote{Ainda não definimos $x^p$ para $x > 0$ e
um $p \in \R$ arbitrário. A definição é $x^p \coloneqq \exp ( p \ln x )$.
Definiremos o logaritmo e a exponencial na \sectionref{sec:logandexp}.
Por enquanto, basta considerar $p$ racional,
para o qual $x^{k/m} = {(x^{1/m})}^{k}$. Veja também o \exerciseref{exercise:realpower}.}.

\begin{prop}[Série $p$ ou critério $p$]%
\index{série $p$@$p$-série}\index{critério p@critério $p$}
Para $p \in \R$,
a série
\begin{equation*}
\sum_{n=1}^\infty \frac{1}{n^p}
\end{equation*}
converge se, e somente se, $p > 1$.
\end{prop}

\begin{proof}
Suponha primeiro que $p \leq 1$.
Como $n \geq 1$, temos
$\frac{1}{n^p} \geq \frac{1}{n}$. Como
$\sum_{n=1}^\infty \frac{1}{n}$ diverge,
$\sum_{n=1}^\infty \frac{1}{n^p}$ também deve divergir para todo $p \leq 1$, pelo critério da comparação.

Suponha agora que $p > 1$.
Procedemos como no caso da
série harmônica, mas, em vez de mostrar que a sequência
das somas parciais é ilimitada, mostraremos que ela é limitada.
Os termos da série são positivos, logo a sequência das somas parciais
é monótona não decrescente e converge se for limitada
superiormente.
Seja $s_n$ a $n$-ésima soma parcial.
\begin{align*}
 s_1 & = 1 , \\
 s_3 & = \left( 1 \right) + \left( \frac{1}{2^p} + \frac{1}{3^p} \right) , \\
 s_7 & = \left( 1 \right) + \left( \frac{1}{2^p} + \frac{1}{3^p} \right) +
        \left( \frac{1}{4^p} + \frac{1}{5^p} + \frac{1}{6^p} + \frac{1}{7^p} \right) , \\
& ~~ \vdots \\
 s_{2^k - 1} &= 
1 + 
\sum_{i=1}^{k-1}
\left(
\sum_{m=2^i}^{2^{i+1}-1} \frac{1}{m^p}
\right) .
\end{align*}
Em vez de estimar por baixo, estimamos por cima.
Como $p$ é positivo, $2^p < 3^p$ e, portanto,
$\frac{1}{2^p} + \frac{1}{3^p} <
\frac{1}{2^p} + \frac{1}{2^p}$. De modo semelhante,
$\frac{1}{4^p} + \frac{1}{5^p} +
\frac{1}{6^p} + \frac{1}{7^p} <
\frac{1}{4^p} + \frac{1}{4^p} +
\frac{1}{4^p} + \frac{1}{4^p}$. Portanto,
para todo $k \geq 2$,
\begin{equation*}
\begin{split}
s_{2^k-1}
& =
1+
\sum_{i=1}^{k-1}
\left(
\sum_{m=2^{i}}^{2^{i+1}-1} \frac{1}{m^p}
\right) 
\\
& <
1+
\sum_{i=1}^{k-1}
\left(
\sum_{m=2^{i}}^{2^{i+1}-1} \frac{1}{{(2^i)}^p}
\right) 
\\
& =
1+
\sum_{i=1}^{k-1}
\left(
\frac{2^i}{{(2^i)}^p}
\right) 
\\
& =
1+
\sum_{i=1}^{k-1}
{\left(
\frac{1}{2^{p-1}}
\right)}^i .
\end{split}
\end{equation*}
Como $p > 1$, temos $\frac{1}{2^{p-1}} < 1$.
A \propref{geometric:prop} afirma que
\begin{equation*}
\sum_{i=1}^\infty
{\left(
\frac{1}{2^{p-1}}
\right)}^i
\end{equation*}
converge. Assim,
\begin{equation*}
s_{2^k-1} < 
1+
\sum_{i=1}^{k-1}
{\left(
\frac{1}{2^{p-1}}
\right)}^i 
\leq 
1+
\sum_{i=1}^\infty
{\left(
\frac{1}{2^{p-1}}
\right)}^i .
\end{equation*}
Para todo $n$, existe $k \geq 2$ tal que $n \leq 2^k-1$; como
$\{ s_n \}_{n=1}^\infty$ é monótona, temos $s_n \leq s_{2^k-1}$.
Assim, para todo $n$,
\begin{equation*}
s_n < 
1+
\sum_{i=1}^\infty
{\left(
\frac{1}{2^{p-1}}
\right)}^i
\end{equation*}
Portanto, a sequência das somas parciais é limitada e a série converge.
\end{proof}

Nem o critério da série $p$ nem o critério da comparação
nos dizem a que valor converge a soma. Eles apenas garantem a existência
do limite das somas parciais. Por exemplo, embora saibamos que
$\sum_{n=1}^\infty \nicefrac{1}{n^2}$ converge, é muito mais difícil determinar
\footnote{A demonstração desse fato foi o que tornou famoso o matemático suíço
\href{https://en.wikipedia.org/wiki/Leonhard_Euler}{Leonhard Euler}
(1707--1783).}
que o limite é $\nicefrac{\pi^2}{6}$.
Se tratarmos $\sum_{n=1}^\infty \nicefrac{1}{n^p}$ como uma função de $p$,
obtemos a função zeta de Riemann. O estudo do
comportamento dessa função envolve
um dos problemas em aberto mais famosos da matemática atual e tem aplicações
em áreas aparentemente não relacionadas, como a criptografia moderna.

\begin{example}
A série $\sum_{n=1}^\infty \frac{1}{n^2+1}$ converge.

Demonstração: primeiro, $\frac{1}{n^2+1} < \frac{1}{n^2}$ para todo $n \in \N$.
A série $\sum_{n=1}^\infty \frac{1}{n^2}$ converge pelo critério da série $p$.
Portanto, pelo critério da comparação, $\sum_{n=1}^\infty \frac{1}{n^2+1}$ converge.
\end{example}

\subsection{Critério da razão}

Suponha que $r > 0$. A razão entre dois termos consecutivos da série geométrica
$\sum_{n=0}^\infty r^n$ é $\frac{r^{n+1}}{r^n} = r$, e a série converge
sempre que $r < 1$. Assim como ocorre com sequências, esse fato
pode ser generalizado para séries mais arbitrárias,
desde que essa razão exista \myquote{no limite}. Então comparamos
a cauda de uma série com a série geométrica.


\begin{prop}[Critério da razão]\index{critério da razão para séries}
Seja $\sum_{n=1}^\infty x_n$ uma série, com $x_n \neq 0$ para todo $n$, e tal que
\begin{equation*}
L \coloneqq \lim_{n\to\infty} \frac{\sabs{x_{n+1}}}{\sabs{x_n}}
\qquad \text{existe.}
\end{equation*}
\begin{enumerate}[(i)]
\item
Se $L < 1$, então $\sum_{n=1}^\infty x_n$ converge absolutamente.
\item
Se $L > 1$, então $\sum_{n=1}^\infty x_n$ diverge.
\end{enumerate}
\end{prop}

Embora o critério, tal como enunciado, seja frequentemente suficiente, ele pode ser um pouco
fortalecido; veja o \exerciseref{exercise:strongerratiotest}.

\begin{proof}
Se $L > 1$, então
\lemmaref{seq:ratiotest} afirma que a sequência $\{ x_n \}_{n=1}^\infty$
diverge. Como uma condição necessária para a convergência de séries
é que seus termos tendam a zero, sabemos que $\sum_{n=1}^\infty x_n$ deve divergir.

Suponhamos, então, que $L < 1$.
Mostraremos que $\sum_{n=1}^\infty \sabs{x_n}$ deve convergir.
A demonstração é semelhante à de \lemmaref{seq:ratiotest}. Naturalmente, $L \geq
0$.
Escolha
$r$ tal que $L < r < 1$. Como $r-L > 0$, existe $M \in \N$ tal que, para
todo $n \geq M$,
\begin{equation*}
\abs{\frac{\sabs{x_{n+1}}}{\sabs{x_n}} - L} < r-L .
\end{equation*}
Portanto,
\begin{equation*}
\frac{\sabs{x_{n+1}}}{\sabs{x_n}} < r .
\end{equation*}
Para $n > M$ (isto é, para $n \geq M+1$),
escrevemos
\begin{equation*}
\sabs{x_n} =
\sabs{x_M}
\frac{\sabs{x_{M+1}}}{\sabs{x_{M}}}
\frac{\sabs{x_{M+2}}}{\sabs{x_{M+1}}}
\cdots
\frac{\sabs{x_{n}}}{\sabs{x_{n-1}}}
<
\sabs{x_M}
r r \cdots r = \sabs{x_M} r^{n-M} =
\bigl(\sabs{x_M} r^{-M}\bigr) r^n .
\end{equation*}
Para $k > M$, escrevemos a soma parcial na forma
\begin{equation*}
\begin{split}
\sum_{n=1}^k \sabs{x_n}
& =
\left(\sum_{n=1}^{M} \sabs{x_n} \right)
+
\left(\sum_{n=M+1}^{k} \sabs{x_n} \right)
\\
& <
\left(\sum_{n=1}^{M} \sabs{x_n} \right)
+
\left(\sum_{n=M+1}^{k} 
\bigl(\sabs{x_M} r^{-M}\bigr) r^n
\right)
\\
& =
\left(\sum_{n=1}^{M} \sabs{x_n} \right)
+
\bigl(\sabs{x_M} r^{-M}\bigr)
\left( \sum_{n=M+1}^{k} r^n \right) .
\end{split}
\end{equation*}
Como $0 < r < 1$, a série geométrica
$\sum_{n=0}^{\infty} r^n$ converge; portanto,
$\sum_{n=M+1}^{\infty} r^n$ também converge. Calculamos o
limite, quando $k$ tende ao infinito, do lado direito acima e obtemos
\begin{equation*}
\begin{split}
\sum_{n=1}^k \sabs{x_n}
& <
\left(\sum_{n=1}^{M} \sabs{x_n} \right)
+
\bigl(\sabs{x_M} r^{-M}\bigr)
\left( \sum_{n=M+1}^{k} r^n \right) 
\\
& \leq
\left(\sum_{n=1}^{M} \sabs{x_n} \right)
+
\bigl(\sabs{x_M} r^{-M}\bigr)
\left( \sum_{n=M+1}^{\infty} r^n \right) .
\end{split}
\end{equation*}
O lado direito é um número que não depende de $k$.
Logo, a sequência das somas parciais de $\sum_{n=1}^\infty \sabs{x_n}$ é limitada,
e $\sum_{n=1}^\infty \sabs{x_n}$ converge. Portanto, $\sum_{n=1}^\infty x_n$ é
absolutamente convergente.
\end{proof}

\begin{example}
A série
\begin{equation*}
\sum_{n=1}^\infty \frac{2^n}{n!}
\end{equation*}
converge absolutamente.

Demonstração: escrevemos
\begin{equation*}
\lim_{n\to\infty} \frac{2^{(n+1)}/(n+1)!}{2^n / n!} =
\lim_{n\to\infty} \frac{2}{n+1} = 0 .
\end{equation*}
Portanto, a série converge absolutamente pelo critério da razão.
\end{example}

\subsection{Exercises}

\begin{exercise}
Suppose the $k$th partial sum of $\displaystyle \sum_{n=1}^\infty x_n$ is $s_k = \frac{k}{k+1}$.
Find the series (that is, find $x_n$), prove that the series converges, and
then find the limit.
\end{exercise}

\begin{exercise} \label{geometric:exr}
Prove \propref{geometric:prop}, that is, for $-1 < r < 1$, prove
\begin{equation*}
\sum_{n=0}^\infty r^n = \frac{1}{1-r} .
\end{equation*}
Hint:  See \exampleref{example:geometricsum}.
\end{exercise}

\begin{exercise}
Decide the convergence or divergence of the following series.

\medskip

\noindent
a)
$\displaystyle \sum_{n=1}^\infty \frac{3}{9n+1}$
\qquad
b)
$\displaystyle \sum_{n=1}^\infty \frac{1}{2n-1}$
\qquad
c)
$\displaystyle \sum_{n=1}^\infty \frac{{(-1)}^n}{n^2}$
\qquad
d)
$\displaystyle \sum_{n=1}^\infty \frac{1}{n(n+1)}$
\qquad
e)
$\displaystyle \sum_{n=1}^\infty n e^{-n^2}$
\end{exercise}

\begin{samepage}
\begin{exercise}
\leavevmode
\begin{enumerate}[a)]
\item Prove that if
$\displaystyle
\sum_{n=1}^\infty x_n
$
converges, then
$\displaystyle
\sum_{n=1}^\infty ( x_{2n} + x_{2n+1} )
$
also converges.
\item
Find an explicit example where the converse does not hold.
\end{enumerate}
\end{exercise}
\end{samepage}

\begin{exercise}
For $i=1,2,\ldots,n$, let $\{ x_{i,k} \}_{k=1}^\infty$ denote $n$
sequences.  Suppose that for each $i$,
\begin{equation*}
\sum_{k=1}^\infty x_{i,k}
\end{equation*}
is convergent.  Prove
\begin{equation*}
\sum_{i=1}^n \left( \sum_{k=1}^\infty x_{i,k} \right)
=
\sum_{k=1}^\infty \left( \sum_{i=1}^n x_{i,k} \right) .
\end{equation*}
\end{exercise}

\begin{exercise} \label{exercise:strongerratiotest}
Prove the following stronger version of the ratio test:
Let $\sum_{n=1}^\infty x_n$ be a series.
\begin{enumerate}[a)]
\item
If there is an $N$ and a $\rho < 1$ such that
$\frac{\sabs{x_{n+1}}}{\sabs{x_n}} < \rho$ for all $n \geq N$,
then the series converges absolutely.
(Remark: Equivalently the condition can be stated as
$\limsup\limits_{n\to\infty} \frac{\sabs{x_{n+1}}}{\sabs{x_n}} < 1$.)
\item
If there is an $N$ such that
$\frac{\sabs{x_{n+1}}}{\sabs{x_n}} \geq 1$
for all $n \geq N$,
then the series diverges. 
\end{enumerate}
\end{exercise}

\begin{exercise}[Challenging]
Suppose $\{ x_n \}_{n=1}^\infty$ is a decreasing sequence and
$\sum_{n=1}^\infty x_n$ converges.
Prove $\displaystyle \lim_{n\to\infty} n x_n = 0$.
\end{exercise}

\begin{exercise}
Show that $\displaystyle \sum_{n=1}^\infty \frac{{(-1)}^n}{n}$ converges.
Hint: Consider the sum of two subsequent entries.
\end{exercise}

\begin{exercise}
\leavevmode
\begin{enumerate}[a)]
\item Prove that if $\sum_{n=1}^\infty x_n$ and $\sum_{n=1}^\infty y_n$ converge absolutely, then
$\sum_{n=1}^\infty x_ny_n$ converges absolutely.
\item Find an explicit example where the converse does not hold.
\item Find an explicit example where all three series are absolutely convergent,
are not just finite sums,
and $\left(\sum_{n=1}^\infty x_n\right)\left(\sum_{n=1}^\infty y_n\right)
\neq \sum_{n=1}^\infty x_ny_n$.  That is, show that series are
not multiplied term-by-term.
\end{enumerate}
\end{exercise}

\begin{exercise}
Prove the triangle inequality for series:
If $\sum_{n=1}^\infty x_n$ converges absolutely, then
\begin{equation*}
\abs{\sum_{n=1}^\infty x_n} \leq
\sum_{n=1}^\infty \sabs{x_n} .
\end{equation*}
\end{exercise}

\begin{exercise}
Prove the \emph{\myindex{limit comparison test}}.  That is, prove that if
$a_n > 0$ and $b_n > 0$ for all $n$, and
\begin{equation*}
0 < \lim_{n\to\infty} \frac{a_n}{b_n} < \infty ,
\end{equation*}
then either $\sum_{n=1}^\infty a_n$ and $\sum_{n=1}^\infty b_n$ both converge or both diverge.
\end{exercise}

\begin{exercise} \label{exercise:badnocauchy}
Let $x_n \coloneqq \sum_{i=1}^n \nicefrac{1}{i}$.  Show that for every $k$,
we get
$\displaystyle \lim_{n\to\infty} \sabs{x_{n+k}-x_n} = 0$, yet
$\{ x_n \}_{n=1}^\infty$ is not Cauchy.
\end{exercise}

\begin{samepage}
\begin{exercise}
Let $s_k$ be the $k$th partial sum of $\sum_{n=1}^\infty x_n$.
\begin{enumerate}[a)]
\item
Suppose there exists an $m \in \N$ such that $\lim\limits_{k\to\infty}
s_{mk}$ exists and $\lim\limits_{n\to\infty} x_n = 0$.  Show that
$\sum_{n=1}^\infty x_n$ converges.
\item
Find an example where $\lim\limits_{k\to\infty} s_{2k}$ exists and
$\lim\limits_{n\to\infty} x_n \neq 0$ (and therefore $\sum_{n=1}^\infty x_n$ diverges).
\item
(Challenging) Find an example where $\lim_{n\to\infty} x_n = 0$, and there exists
a subsequence $\{ s_{k_i} \}_{i=1}^\infty$ such that $\displaystyle \lim_{i\to\infty} s_{k_i}$ exists,
but $\sum_{n=1}^\infty x_n$ still diverges.
\end{enumerate}
\end{exercise}
\end{samepage}

\begin{exercise} \label{exercise:squareseriesconv}
Suppose $\sum_{n=1}^\infty x_n$ converges and $x_n \geq 0$ for all $n$.
Prove that $\sum_{n=1}^\infty x_n^2$ converges.
\end{exercise}

\begin{exercise}[Challenging] \label{exercise:cauchycondensation}
Suppose $\{ x_n\}_{n=1}^\infty$ is a decreasing sequence of positive numbers.
The proof of convergence/divergence for the $p$-series generalizes.
Prove the
\emph{\myindex{Cauchy condensation principle}}:
\begin{equation*}
\sum_{n=1}^\infty x_n
\qquad \text{converges if and only if} \qquad
\sum_{n=1}^\infty 2^n x_{2^n} \qquad \text{converges}.
\end{equation*}
\end{exercise}

\begin{exercise}
Use the Cauchy condensation principle
(see \exerciseref{exercise:cauchycondensation})
to decide the convergence of

\medskip

\noindent
a) $\displaystyle \sum_{n=1}^\infty \frac{\ln n}{n^2}$
\qquad
b) $\displaystyle \sum_{n=2}^\infty \frac{1}{n \ln n}$
\qquad
c) $\displaystyle \sum_{n=2}^\infty \frac{1}{n {(\ln n)}^2}$
\qquad
d) $\displaystyle \sum_{n=2}^\infty \frac{1}{n (\ln n ){(\ln \ln n)}^2}$

\medskip

\noindent
Note that only the tails of some of these series satisfy the
hypotheses of the principle; you should argue why that is sufficient.

\noindent
Hint: Feel free to use the identity $\ln (2^n) = n \ln 2$.
\end{exercise}

\begin{exercise}[Challenging]
Prove \emph{\myindex{Abel's theorem}}:

\medskip

\noindent
\emph{\textbf{Theorem.} Suppose $\sum_{n=1}^\infty x_n$ is a series whose sequence of partial sums
is bounded, $\{ \lambda_n \}_{n=1}^\infty$ is a sequence with
$\lim_{n\to\infty} \lambda_n = 0$, and
$\sum_{n=1}^\infty \sabs{ \lambda_{n+1} - \lambda_n }$ is convergent.
Then $\sum_{n=1}^\infty \lambda_n x_n$ is convergent.}
\end{exercise}



%%%%%%%%%%%%%%%%%%%%%%%%%%%%%%%%%%%%%%%%%%%%%%%%%%%%%%%%%%%%%%%%%%%%%%%%%%%%%%

\sectionnewpage
\section{More on series}
\label{sec:moreonseries}

%mbxINTROSUBSECTION

\sectionnotes{up to 2--3 lectures (optional, can safely be skipped or covered partially)}

\subsection{Root test}

A test similar to the ratio test is the
\emph{\myindex{root test}}.  The proof of this test is similar.
Again, the idea is to generalize what happens for the geometric series.

\begin{prop}[Root test]
Let $\sum_{n=1}^\infty x_n$ be a series and let
\begin{equation*}
L \coloneqq \limsup_{n\to\infty} {\sabs{x_n}}^{1/n} .
\end{equation*}
\begin{enumerate}[(i)]
\item If $L < 1$, then $\sum_{n=1}^\infty x_n$ converges absolutely.
\item If $L > 1$, then $\sum_{n=1}^\infty x_n$ diverges.
\end{enumerate}
\end{prop}

\begin{proof}
If $L > 1$, then there exists\footnote{%
In case $L=\infty$, see \exerciseref{exercise:extendsubseqlimsupinf}.
Alternatively, note that if $L=\infty$, then $\bigl\{ {\sabs{x_{n}}}^{1/n}
\bigr\}_{n=1}^\infty$
is unbounded, and thus so is
$\{ x_n \}_{n=1}^\infty$.}
a subsequence $\{ x_{n_k} \}_{k=1}^\infty$ such that
$L = \lim_{k\to\infty} {\sabs{x_{n_k}}}^{1/n_k}$.  Let
$r$ be such that $L > r > 1$.  There exists an $M$ such
that for all $k \geq M$, we have 
${\sabs{x_{n_k}}}^{1/n_k} > r > 1$, or in other words
$\sabs{x_{n_k}} > r^{n_k} > 1$.
The subsequence 
$\bigl\{ \sabs{x_{n_k}} \bigr\}_{k=1}^\infty$,
and therefore also
$\bigl\{ \sabs{x_{n}} \bigr\}_{n=1}^\infty$,
cannot possibly converge to zero, and so the series diverges.

Now suppose $L < 1$.  Pick $r$ such that $L < r < 1$.
By definition of limit supremum,
there is an $M$ such that for all $n \geq M$,
\begin{equation*}
\sup \bigl\{ {\sabs{x_k}}^{1/k} : k \geq n \bigr\} < r .
\end{equation*}
Therefore, for all $n \geq M$,
\begin{equation*}
{\sabs{x_n}}^{1/n} < r , \qquad \text{or in other words} \qquad \sabs{x_n} < r^n .
\end{equation*}
Let $k > M$, and estimate the $k$th partial sum:
\begin{equation*}
\sum_{n=1}^k \sabs{x_n} = 
\left( \sum_{n=1}^M \sabs{x_n} \right) + 
\left( \sum_{n=M+1}^k \sabs{x_n} \right)
\leq
\left( \sum_{n=1}^M \sabs{x_n} \right) + 
\left( \sum_{n=M+1}^k r^n \right) .
\end{equation*}
As $0 < r < 1$,
the geometric series $\sum_{n=M+1}^\infty r^n$ converges to
$\frac{r^{M+1}}{1-r}$.  As everything is positive,
\begin{equation*}
\sum_{n=1}^k \sabs{x_n} 
\leq
\left( \sum_{n=1}^M \sabs{x_n} \right) + 
\frac{r^{M+1}}{1-r} .
\end{equation*}
Thus the sequence of partial sums of $\sum_{n=1}^\infty \sabs{x_n}$ is bounded, and
the series converges.  Therefore, $\sum_{n=1}^\infty x_n$ converges absolutely.
\end{proof}

\subsection{Alternating series test}

The tests we have seen so far only addressed absolute convergence.  The
following test gives a large supply of conditionally convergent series.

\begin{prop}[Alternating series]
Let $\{ x_n \}_{n=1}^\infty$ be a monotone decreasing sequence of positive real numbers such
that $\lim\limits_{n\to\infty} x_n = 0$.  Then
\begin{equation*}
\sum_{n=1}^\infty {(-1)}^n x_n \qquad \text{converges.}
\end{equation*}
\end{prop}

\begin{proof}
Let $s_m \coloneqq \sum_{n=1}^m {(-1)}^n x_n$ be the $m$th partial sum.  Then
\begin{equation*}
s_{2k} =
\sum_{n=1}^{2k} {(-1)}^n x_n
=
(-x_1 + x_2) + \cdots + (-x_{2k-1} + x_{2k})
=
\sum_{\ell=1}^{k} (-x_{2\ell-1} + x_{2\ell}) .
\end{equation*}
The sequence $\{ x_n \}_{n=1}^\infty$ is decreasing, so $(-x_{2\ell-1}+x_{2\ell}) \leq 0$
for all $\ell$.
Thus, the subsequence $\{ s_{2k} \}_{k=1}^\infty$ of partial sums
is a decreasing sequence.  Similarly, $(x_{2\ell}-x_{2\ell+1}) \geq 0$, and so
\begin{equation*}
s_{2k} = - x_1 + ( x_2 - x_3 ) + \cdots + ( x_{2k-2} - x_{2k-1} ) + x_{2k}
\geq -x_1 .
\end{equation*}
The intuition behind the bound
$0 \geq s_{2k} \geq -x_1$ is illustrated in \figureref{fig:seralternate}.
\begin{myfigureht}
\myincludegraphics{ser-alternate}{%
A diagram in the plane where the values of a sequence plus minus x sub n are shown
for n from 1 to 8.  In solid lines starting from the horizontal axis,
we see a "graph" of minus x sub 1, x sub 2, minus x sub 3, x sub 4, etc.
We can see that each successive solid line is shorter
as each successive x sub n is smaller than the last.  In dotted lines we
also have minus x sub 2, minus x sub 4, etc.
The differences
minus x sub 1 plus x sub 2,
minus x sub 3 plus x sub 4,
minus x sub 5 plus x sub 6,
and
minus x sub 7 plus x sub 8, are shown and correspond to the differences
between the corresponding minus x sub 1 and minus x sub 2.
We note that these differences take up just part of the total length of the
line for minus x sub 1.}
\caption{Showing that $0 \geq s_{2k} \geq -x_1$ where $k=4$ for an alternating series.\label{fig:seralternate}}
\end{myfigureht}

As $\{ s_{2k} \}_{k=1}^\infty$ is decreasing and bounded below, it converges.
Let $a \coloneqq \lim_{k\to\infty} s_{2k}$.
We wish to show that $\lim_{m\to\infty} s_m = a$ (and not just for the subsequence).
Given $\epsilon > 0$, pick $M$ such that $\sabs{s_{2k}-a} <
\nicefrac{\epsilon}{2}$ whenever $k \geq M$.
Since $\lim_{n\to\infty} x_n = 0$, we also
make $M$ possibly larger
to obtain
$x_{2k+1} < \nicefrac{\epsilon}{2}$ whenever $k \geq M$.  
Suppose $m \geq 2M+1$.  If $m=2k$, then $k \geq M+\nicefrac{1}{2} \geq M$ and
$\sabs{s_{m}-a}=\sabs{s_{2k}-a} < \nicefrac{\epsilon}{2} < \epsilon$.
If $m=2k+1$, then also $k \geq M$.  Notice
$s_{2k+1} = s_{2k} - x_{2k+1}$.
Thus
\begin{equation*}
\sabs{s_{m}-a} = 
\sabs{s_{2k+1}-a} = 
\sabs{s_{2k}-a - x_{2k+1}} \leq
\sabs{s_{2k}-a} + x_{2k+1} < 
\nicefrac{\epsilon}{2}+ \nicefrac{\epsilon}{2} = \epsilon .  \qedhere
\end{equation*}
\end{proof}

Notably, there exist conditionally convergent series
where the absolute values of the terms go to zero arbitrarily slowly.
The series
\begin{equation*}
\sum_{n=1}^\infty \frac{{(-1)}^n}{n^p}
\end{equation*}
converges for arbitrarily small $p > 0$, but it does not converge
absolutely when $p \leq 1$.

\subsection{Rearrangements}

Absolutely convergent series behave as we imagine they should.  For example,
absolutely convergent series can be summed in any order whatsoever.  Nothing
of the sort holds for conditionally convergent series
(see \exampleref{example:harmonsumanything}
and \exerciseref{exercise:seriesconvergestoanything}).

Consider a series
\begin{equation*}
\sum_{n=1}^\infty x_n .
\end{equation*}
Given a bijective function $\sigma \colon \N \to \N$, the corresponding
rearrangement\index{rearrangement of a series} is the 
series:
\begin{equation*}
\sum_{k=1}^\infty x_{\sigma(k)} .
\end{equation*}
We simply sum the series in a different order.

\begin{prop}
Let $\sum_{n=1}^\infty x_n$ be an absolutely convergent series converging to a number
$x$.  Let $\sigma \colon \N \to \N$ be a bijection.  Then
$\sum_{n=1}^\infty x_{\sigma(n)}$ is absolutely convergent and converges to $x$.
\end{prop}

In other words,
a rearrangement of an absolutely convergent series converges (absolutely)
to the same number.

\begin{proof}
Let $\epsilon > 0$ be given.  As $\sum_{n=1}^\infty x_n$ is absolutely convergent, take $M$ such that
\begin{equation*}
\abs{\left(\sum_{n=1}^M x_n \right) - x} < \frac{\epsilon}{2}
\qquad \text{and} \qquad
\sum_{n=M+1}^\infty \sabs{x_n} < \frac{\epsilon}{2} .
\end{equation*}
As $\sigma$ is a bijection,
there exists a number $K$ such that for each
$n \leq M$, there exists $k \leq K$ such that $\sigma(k) = n$.
In other words
$\{ 1,2,\ldots,M \} \subset \sigma\bigl(\{ 1,2,\ldots,K \} \bigr)$.

For $N \geq K$, let $Q \coloneqq \max \sigma\bigl(\{ 1,2,\ldots,N \}\bigr)$.
Compute
\begin{equation*}
\begin{split}
\abs{\left( \sum_{n=1}^N x_{\sigma(n)} \right) - x}
& =
\abs{ \left( \sum_{n=1}^M x_n
+
\sum_{\substack{n=1\\\sigma(n) > M}}^N x_{\sigma(n)} \right) - x}
\\
& \leq
\abs{ \left( \sum_{n=1}^M x_n \right) - x}
+
\sum_{\substack{n=1\\\sigma(n) > M}}^N \sabs{x_{\sigma(n)}}
\\
& \leq
\abs{ \left( \sum_{n=1}^M x_n \right) - x}
+
\sum_{n=M+1}^Q \sabs{x_{n}}
\\
& < \nicefrac{\epsilon}{2} + \nicefrac{\epsilon}{2} = \epsilon .
\end{split}
\end{equation*}
So 
$\sum_{n=1}^\infty x_{\sigma(n)}$ converges to $x$.  To see that the convergence
is absolute, we apply the argument above to $\sum_{n=1}^\infty \sabs{x_n}$ to show
that $\sum_{n=1}^\infty \sabs{x_{\sigma(n)}}$ converges.
\end{proof}

\begin{example} \label{example:harmonsumanything}
Let us show that the alternating harmonic series
$\sum_{n=1}^\infty \frac{{(-1)}^{n+1}}{n}$, which does not converge absolutely, can be
rearranged to converge to anything.
The odd terms and the even terms diverge to
plus infinity and minus infinity, respectively (prove this!):
\begin{equation*}
\sum_{m=1}^\infty \frac{1}{2m-1} = \infty, \qquad \text{and} \qquad
\sum_{m=1}^\infty \frac{-1}{2m} = -\infty .
\end{equation*}
Let $a_n \coloneqq \frac{{(-1)}^{n+1}}{n}$ for simplicity, 
let an arbitrary number $L \in \R$ be given, and set $\sigma(1) \coloneqq 1$.
Suppose we have
defined $\sigma(n)$ for all $n \leq N$.  If
\begin{equation*}
\sum_{n=1}^N a_{\sigma(n)} \leq L ,
\end{equation*}
then let $\sigma(N+1) \coloneqq k$ be the smallest odd $k \in \N$
that we have not used yet,
that is, $\sigma(n) \neq k$ for all $n \leq N$.
Otherwise, let $\sigma(N+1) \coloneqq k$ 
be the smallest even $k$ that we have not yet used.

By construction, $\sigma \colon \N \to \N$ is one-to-one.
It is also onto, because if we keep adding either odd (resp.\ even) terms,
eventually we pass $L$ and switch
to the evens (resp.\ odds).  So we switch infinitely many times.

Finally, let $N$ be the $N$ where we just pass $L$ and switch.
Suppose we have just switched from odd to even, so we start
subtracting.
Suppose we keep subtracting until some
$N' > N$ and we will switch back from even to odd in the next step.
Then
\begin{equation*}
L + \frac{1}{\sigma(N)} \geq \sum_{n=1}^{N} a_{\sigma(n)}
> \sum_{n=1}^{N'} a_{\sigma(n)} > L- \frac{1}{\sigma(N')}.
\end{equation*}
Similarly for switching in the other direction.
Therefore,
the sum up to $N'$ is within $\frac{1}{\min \{ \sigma(N), \sigma(N') \}}$
of $L$.  As
we switch infinitely many times, $\sigma(N) \to \infty$
and $\sigma(N') \to \infty$. Hence,
\begin{equation*}
\sum_{n=1}^\infty a_{\sigma(n)} = 
\sum_{n=1}^\infty \frac{{(-1)}^{\sigma(n)+1}}{\sigma(n)} = L .
\end{equation*}

Here is an example to illustrate the proof.  When $L=1.2$, the order
is
\begin{equation*}
1+\nicefrac{1}{3}-\nicefrac{1}{2}+\nicefrac{1}{5}+\nicefrac{1}{7}+\nicefrac{1}{9}-\nicefrac{1}{4}+\nicefrac{1}{11}+\nicefrac{1}{13}-\nicefrac{1}{6}
+\nicefrac{1}{15}+\nicefrac{1}{17}+\nicefrac{1}{19} - \nicefrac{1}{8} + \cdots .
\end{equation*}
At this point we are no more than $\nicefrac{1}{8}$ from the limit.
See \figureref{fig:serrearrange}.
\begin{myfigureht}
\myincludegraphics{ser-rearrange}{%
A diagram of a sequence given as dots.  A horizontal dashed line is shown.
The sequence starts below the dashed line, then is above it in the next
step.  Then it goes below the dashed line again, then it takes a few steps
before it is above the dashed line, and the next step is again below the
dashed line.  And so on and so forth.  The steps get shorter and shorter.}
\caption{The first 14 partial sums of the rearrangement converging
to $1.2$.\label{fig:serrearrange}}
\end{myfigureht}
\end{example}

\subsection{Multiplication of series}

As we have
already mentioned,
multiplication of series is somewhat harder than addition.
If at least one of the series converges
absolutely, then we can use the following theorem.  For this result, it is
convenient to start the series at 0, rather than at 1.

\begin{thm}[\myindex{Mertens' theorem}\footnote{Proved by
the German mathematician
\href{https://en.wikipedia.org/wiki/Franz_Mertens}{Franz Mertens}
(1840--1927).}]
Suppose $\sum_{n=0}^\infty a_n$ and $\sum_{n=0}^\infty b_n$ are two convergent series, converging
to $A$ and $B$, respectively.  Suppose at least one of the series
converges absolutely.  Define
\begin{equation*}
c_n \coloneqq a_0 b_n + a_1 b_{n-1} + \cdots + a_n b_0 = \sum_{i=0}^n a_i b_{n-i} .
\end{equation*}
Then the series $\sum_{n=0}^\infty c_n$
converges to $AB$.
\end{thm}

The series $\sum_{n=0}^\infty c_n$ is called the \emph{\myindex{Cauchy product}} of
$\sum_{n=0}^\infty a_n$ and $\sum_{n=0}^\infty b_n$.

\begin{proof}
Suppose $\sum_{n=0}^\infty a_n$ converges absolutely, and let $\epsilon > 0$ be
given.
In this proof, instead of picking complicated estimates just to make
the final estimate come out as less than $\epsilon$,
let us simply obtain an estimate that depends on $\epsilon$
and can be made arbitrarily small.

Write
\begin{equation*}
A_m \coloneqq \sum_{n=0}^m a_n , \qquad B_m \coloneqq \sum_{n=0}^m b_n .
\end{equation*}
We rearrange the $m$th partial sum of $\sum_{n=0}^\infty c_n$:
\begin{equation*}
\begin{split}
\abs{\left(\sum_{n=0}^m c_n \right) - AB}
& =
\abs{\left( \sum_{n=0}^m \sum_{i=0}^n a_i b_{n-i} \right) - AB}
\\
& =
\abs{\left( \sum_{n=0}^m
  B_n a_{m-n} \right) - AB}
\\
& =
\abs{\left( \sum_{n=0}^m
  ( B_n -  B ) a_{m-n} \right)
    + B A_m - AB}
\\
& \leq
\left(
\sum_{n=0}^m
  \sabs{ B_n -  B } \sabs{a_{m-n}}
\right)
+
\sabs{B}\sabs{A_m - A}
\end{split}
\end{equation*}
We can surely make the second term on the right-hand side small.
The trick is to handle the first term.
Pick $K$ such that for all $m \geq K$, we have 
$\sabs{A_m - A} < \epsilon$ and
also
$\sabs{B_m - B} < \epsilon$.
As $\sum_{n=0}^\infty a_n$ converges absolutely,
there is a $K$ large enough such that
for all $m \geq K$,
\begin{equation*}
\sum_{n=K}^m \sabs{a_n} < \epsilon .
\end{equation*}
As $\sum_{n=0}^\infty b_n$ converges,
$B_{\text{max}} \coloneqq \sup \bigl\{ \sabs{ B_n - B } : n = 0,1,2,\ldots
\bigr\}$
is finite.  Take $m \geq 2K$.
In particular, $m-K+1 > K$.  So
\begin{equation*}
\begin{split}
\sum_{n=0}^m
  \sabs{ B_n -  B } \sabs{a_{m-n}}
& =
\left(
\sum_{n=0}^{m-K}
  \sabs{ B_n -  B } \sabs{a_{m-n}}
\right)
+
\left(
\sum_{n=m-K+1}^m
  \sabs{ B_n -  B } \sabs{a_{m-n}}
\right)
\\
& \leq
\left(
\sum_{n=K}^m
\sabs{a_{n}}
\right)
B_{\text{max}}
+
\left(
\sum_{n=0}^{K-1}
  \epsilon \sabs{a_{n}}
\right)
\\
& \leq
\epsilon
B_{\text{max}}
+
\epsilon
\left(
\sum_{n=0}^\infty \sabs{a_{n}}
\right) .
\end{split}
\end{equation*}
Therefore, for $m \geq 2K$, we have
\begin{equation*}
\begin{split}
\abs{\left(\sum_{n=0}^m c_n \right) - AB}
& \leq
\left(
\sum_{n=0}^m
  \sabs{ B_n -  B } \sabs{a_{m-n}}
\right)
+
\sabs{B}\sabs{A_m - A}
\\
& \leq
\epsilon
B_{\text{max}}
+
\epsilon
\left(
\sum_{n=0}^\infty \sabs{a_{n}}
\right)
+
\sabs{B}\epsilon
=
\epsilon 
\left(
B_{\text{max}}
+
\left(
\sum_{n=0}^\infty \sabs{a_{n}}
\right)
+
\sabs{B}
\right) .
\end{split}
\end{equation*}
The expression in the parentheses on the right-hand side
is a fixed number.
Hence,
we can make the right-hand side arbitrarily small by picking a small enough
$\epsilon> 0$.  So $\sum_{n=0}^\infty c_n$ converges to $AB$.
\end{proof}

\begin{example}
If both series are only conditionally convergent, the Cauchy product series
need not even converge.
Suppose we take $a_n = b_n = {(-1)}^n \frac{1}{\sqrt{n+1}}$.
The series $\sum_{n=0}^\infty a_n = \sum_{n=0}^\infty b_n$
converges
by the alternating series test; however, it does not converge
absolutely as can be seen from the $p$-test.  Let us look
at the Cauchy product.
\begin{equation*}
%mbxlatex \begin{aligned}
c_n
%mbxlatex &
= 
{(-1)}^n
\left(
\frac{1}{\sqrt{n+1}} + 
\frac{1}{\sqrt{2n}} + 
\frac{1}{\sqrt{3(n-1)}} + \cdots +
%\frac{1}{\sqrt{2n}} + 
\frac{1}{\sqrt{n+1}}
\right)
%mbxlatex \\
%mbxlatex &
=
{(-1)}^n
\sum_{i=0}^n \frac{1}{\sqrt{(i+1)(n-i+1)}} .
%mbxlatex \end{aligned}
\end{equation*}
Therefore,
\begin{equation*}
\sabs{c_n} 
=
\sum_{i=0}^n \frac{1}{\sqrt{(i+1)(n-i+1)}} 
\geq
\sum_{i=0}^n \frac{1}{\sqrt{(n+1)(n+1)}} 
= 1 .
\end{equation*}
The terms do not go to zero and hence $\sum_{n=0}^\infty c_n$ cannot converge.
\end{example}

\subsection{Power series}

Fix $x_0 \in \R$.
A \emph{\myindex{power series}} about $x_0$
is a series of the form
\begin{equation*}
\sum_{n=0}^\infty a_n {(x-x_0)}^n .
\end{equation*}
A power series is really a function of $x$, and
many important functions in analysis can be written
as power series.  We use the convention that
$0^0 = 1$ (that is, if $x=x_0$ and $n=0$).

A power series is said to be
\emph{convergent}\index{convergent!power series} if
there is at least one $x \neq x_0$ that makes the series converge.
If $x=x_0$, then the series always
converges since all terms except the first are zero.
If the series does not converge for any point $x \neq x_0$, we say that
the series is \emph{divergent}\index{divergent!power series}.

\begin{example} \label{ps:expex}
The series
\begin{equation*}
\sum_{n=0}^\infty \frac{1}{n!} x^n
\end{equation*}
is absolutely convergent for all $x \in \R$ using the ratio test:
For any $x \in \R$
\begin{equation*}
\lim_{n \to \infty}
\frac{\bigl(1/(n+1)!\bigr) \, x^{n+1}}{(1/n!) \, x^{n}}
=
\lim_{n \to \infty}
\frac{x}{n+1}
=
0.
\end{equation*}
Recall from calculus that this series converges to $e^x$.
\end{example}

\begin{example} \label{ps:1kex}
The series
\begin{equation*}
\sum_{n=1}^\infty \frac{1}{n} x^n
\end{equation*}
converges absolutely for all $x \in (-1,1)$ via the ratio test:
\begin{equation*}
\lim_{n \to \infty}
\abs{
\frac{\bigl(1/(n+1) \bigr) \, x^{n+1}}{(1/n) \, x^{n}}
}
=
\lim_{n \to \infty}
\sabs{x} \frac{n}{n+1}
=
\sabs{x} < 1 .
\end{equation*}
The series converges at $x=-1$,
as
$\sum_{n=1}^\infty \frac{{(-1)}^n}{n}$ converges
by the alternating series
test.
But the power series does not converge absolutely at $x=-1$, because
$\sum_{n=1}^\infty \frac{1}{n}$ does not converge.
The series
diverges at $x=1$.
When $\sabs{x} > 1$, then the series diverges via the ratio test.
\end{example}

\begin{example} \label{ps:divergeex}
The series
\begin{equation*}
\sum_{n=1}^\infty n^n x^n
\end{equation*}
diverges for all $x \neq 0$.  Let us apply the root test
\begin{equation*}
\limsup_{n\to\infty}
\,
\sabs{n^n x^n}^{1/n}
=
\limsup_{n\to\infty}
\,
n \sabs{x}
= \infty .
\end{equation*}
Therefore, the series diverges for all $x \neq 0$.
\end{example}

Convergence of power series in general works analogously to
the three examples above.

\begin{prop} \label{prop:powerserrealradius}
Let $\sum_{n=0}^\infty a_n {(x-x_0)}^n$ be a power series.
If the series is convergent, then either it converges absolutely at
all $x \in \R$, or
there exists a number $\rho$, such that
the series converges absolutely on the interval
$(x_0-\rho,x_0+\rho)$ and diverges when $x < x_0-\rho$ or $x > x_0+\rho$.
\end{prop}

The number $\rho$ is called the \emph{\myindex{radius of convergence}} of the
power series.  We write $\rho = \infty$ if the series converges for
all $x$, and we write $\rho = 0$ if the series is divergent.
At the endpoints, that is, if $x = x_0+\rho$ or $x = x_0-\rho$,
the proposition says nothing,
and the series might or might not converge.
See \figureref{ps:convfig}.
In \exampleref{ps:1kex},
the radius of convergence is $\rho=1$,
in \exampleref{ps:expex}, the radius of convergence is $\rho=\infty$,
and in \exampleref{ps:divergeex}, the radius of convergence is $\rho=0$.

\begin{myfigureht}
\myincludepdft{ps-conv}{%
A diagra of the real line with three points marked, from left to
right, x sub 0 minus rho, x sub 0, and x sub 0 plus rho.
The part of the line to the left of x sub 0 minus rho and
the part of the line to the right of x sub 0 plus rho is labeled
as "diverges".  The interval between x sub 0 minus rho and
x sub 0 plus rho is shaded and labeled as "converges absolutely". 
The two points x sub 0 minus rho and x sub 0 plus rho are labeled
with question marks.}
\caption{Convergence of a power series.\label{ps:convfig}}
\end{myfigureht}

\begin{proof}
Write
\begin{equation*}
R \coloneqq \limsup_{n\to\infty} {\sabs{a_n}}^{1/n} .
\end{equation*}
We apply the root test,
\begin{equation*}
L = \limsup_{n\to\infty} {\babs{a_n {(x-x_0)}^n}}^{1/n} 
=
\sabs{x-x_0} \limsup_{n\to\infty} {\sabs{a_n}}^{1/n}
=
\sabs{x-x_0} R .
\end{equation*}
If $R = \infty$, then $L=\infty$ for every $x \neq x_0$, and
the series diverges by the root test.
On the other hand,
if $R = 0$, then $L=0$ for every $x$,
and the series converges absolutely for all $x$.

Suppose $0 < R < \infty$.
The series
converges absolutely if
$1 > L = R \sabs{x-x_0}$,
that is,
\begin{equation*}
\sabs{x-x_0} < \nicefrac{1}{R} .
\end{equation*}
The series diverges when
$1 < L = R \sabs{x-x_0}$,
or
\begin{equation*}
\sabs{x-x_0} > \nicefrac{1}{R} .
\end{equation*}
Letting $\rho \coloneqq \nicefrac{1}{R}$ completes the proof.
\end{proof}

It may be useful to restate what we have learned in the proof
as a separate proposition.

\begin{prop}
Let $\sum_{n=0}^\infty a_n {(x-x_0)}^n$ be a power series, and let
\begin{equation*}
R \coloneqq \limsup_{n\to\infty} {\sabs{a_n}}^{1/n} .
\end{equation*}
If $R = \infty$, the power series is divergent.  If
$R=0$, then the power series converges everywhere.   Otherwise,
the radius of convergence $\rho = \nicefrac{1}{R}$.
\end{prop}

Often, the radius of convergence is written as $\rho = \nicefrac{1}{R}$ in all
three cases, with
the understanding of what $\rho$ should be if $R = 0$ or $R =
\infty$.

\pagebreak[2]
Convergent power series can be added and multiplied together, and multiplied
by constants.
The proposition has a straightforward proof using what we know about series
in general, and power series in particular.  We leave the proof to the reader.

\begin{prop}
Let $\sum_{n=0}^\infty a_n {(x-x_0)}^n$ and
$\sum_{n=0}^\infty b_n {(x-x_0)}^n$ be two convergent power series
with radius of convergence at least $\rho > 0$ and $\alpha \in \R$.  Then
for all $x$ such that $\sabs{x-x_0} < \rho$, we have 
\begin{equation*}
\left(\sum_{n=0}^\infty a_n {(x-x_0)}^n\right)
+
\left(\sum_{n=0}^\infty b_n {(x-x_0)}^n\right)
=
\sum_{n=0}^\infty (a_n+b_n) {(x-x_0)}^n ,
\end{equation*}
\begin{equation*}
\alpha
\left(\sum_{n=0}^\infty a_n {(x-x_0)}^n\right)
=
\sum_{n=0}^\infty \alpha a_n {(x-x_0)}^n ,
\end{equation*}
and
\begin{equation*}
\left(\sum_{n=0}^\infty a_n {(x-x_0)}^n\right)
\,
\left(\sum_{n=0}^\infty b_n {(x-x_0)}^n\right)
=
\sum_{n=0}^\infty c_n {(x-x_0)}^n ,
\end{equation*}
where
$c_n = a_0b_n + a_1 b_{n-1} + \cdots + a_n b_0$.
\end{prop}

For all $x$ with $\sabs{x-x_0} < \rho$, we have two convergent series, so
their term-by-term addition and multiplication by constants
follows by the previous section.
As for such $x$ the series converges absolutely,
we can apply Mertens' theorem to find the product of two series.
Consequently, after performing the algebraic operations, the
radius of convergence of the resulting series is at least $\rho$.
The radius of convergence of the result could be strictly larger than
the radius of convergence of either of the series we started with.
See the exercises.

Let us look at some examples of power series.
Polynomials are simply finite power series:  A polynomial
is a power series where
the $a_n$ are zero for all $n$ large enough.  We expand
a polynomial as a power series about any point $x_0$ by writing
the polynomial as a polynomial in $(x-x_0)$.  For example,
$2x^2-3x+4$ as a power series around $x_0 = 1$ is
\begin{equation*}
2x^2-3x+4 = 3 + (x-1) + 2{(x-1)}^2 .
\end{equation*}

We can also expand
\emph{\myindex{rational functions}} (that is, ratios of polynomials)
as power series, although we will not completely prove this fact here.
Notice that a series for a rational function only defines the function
on an interval even if the function is defined elsewhere.  For example, for
the \emph{\myindex{geometric series}}, we have that for
$x \in (-1,1)$,
\begin{equation*}
\frac{1}{1-x} =
\sum_{n=0}^\infty x^n .
\end{equation*}
The series diverges when $\sabs{x} > 1$, even though $\frac{1}{1-x}$ is
defined for all $x \neq 1$.

We can use the geometric series together with rules for addition and
multiplication of power series to expand rational functions as power
series around $x_0$,
as long as the denominator is not zero at $x_0$.  We state without
proof that this is always possible, and we give an example of such
a computation using the geometric series.

\begin{example}
Let us expand $\frac{x}{1+2x+x^2}$ as a power series around the origin ($x_0 = 0$) and
find the radius of convergence.

Write $1+2x+x^2 = {(1+x)}^2 = {\bigl(1-(-x)\bigr)}^2$, and suppose
$\sabs{x} < 1$.  Compute
\begin{equation*}
\begin{split}
\frac{x}{1+2x+x^2}
&=
x \,
{\left(
\frac{1}{1-(-x)}
\right)}^2
\\
&=
x \,
{\left( 
\sum_{n=0}^\infty {(-1)}^n x^n 
\right)}^2
\\
&=
x \,
\left(
\sum_{n=0}^\infty c_n x^n 
\right)
\\
&=
\sum_{n=0}^\infty c_n x^{n+1} .
\end{split}
\end{equation*}
Using the formula for the product of series,
we obtain $c_0 = 1$, $c_1 = -1 -1 = -2$, $c_2 = 1+1+1 = 3$, etc.
Hence, for $\sabs{x} < 1$, 
\begin{equation*}
\frac{x}{1+2x+x^2}
=
\sum_{n=1}^\infty {(-1)}^{n+1} n x^n .
\end{equation*}
The radius of convergence is at least 1.  We leave it to the reader to
verify that the radius of convergence is exactly equal to 1.
\end{example}

You can use the method of partial fractions you know from calculus.
For example, to find the power series for $\frac{x^3+x}{x^2-1}$ at 0, write
\begin{equation*}
\frac{x^3+x}{x^2-1}
=
x + \frac{1}{1+x} - \frac{1}{1-x}
=
x + \sum_{n=0}^\infty {(-1)}^n x^n - \sum_{n=0}^\infty x^n .
\end{equation*}

\subsection{Exercises}

\begin{exercise}
Decide the convergence or divergence of the following series.

\medskip

\noindent
a)
$\displaystyle \sum_{n=1}^\infty \frac{1}{2^{2n+1}}$
\qquad
b)
$\displaystyle \sum_{n=1}^\infty \frac{{(-1)}^{n}(n-1)}{n}$
\qquad
c)
$\displaystyle \sum_{n=1}^\infty \frac{{(-1)}^n}{n^{1/10}}$
\qquad
d)
$\displaystyle \sum_{n=1}^\infty \frac{n^n}{{(n+1)}^{2n}}$
\end{exercise}

\begin{exercise}
Suppose both $\sum_{n=0}^\infty a_n$ and $\sum_{n=0}^\infty b_n$ 
converge absolutely.
Show that the product series, $\sum_{n=0}^\infty c_n$ where
$c_n = a_0 b_n + a_1 b_{n-1} + \cdots + a_n b_0$, also converges absolutely.
\end{exercise}

\begin{exercise}[Challenging] \label{exercise:seriesconvergestoanything}
Let $\sum_{n=1}^\infty a_n$ be conditionally convergent.
Show that given an arbitrary $x \in \R$
there exists a rearrangement of $\sum_{n=1}^\infty a_n$
such that the rearranged series converges to $x$.
Hint: See \exampleref{example:harmonsumanything}.
\end{exercise}

\begin{exercise}
\pagebreak[2]
\leavevmode
\begin{enumerate}[a)]
\item
Show that the alternating harmonic series $\sum_{n=1}^\infty \frac{{(-1)}^{n+1}}{n}$
has a rearrangement
such that for every interval $(x,y)$, there exists a partial sum $s_n$
of the rearranged series such that $s_n \in (x,y)$.
\item
Show that the rearrangement you found does not converge.
See \exampleref{example:harmonsumanything}.
\item
Show that for every $x \in \R$, there exists a subsequence of
partial sums $\{ s_{n_k} \}_{k=1}^\infty$ of your rearrangement such that 
$\lim\limits_{k\to\infty} s_{n_k} = x$.
\end{enumerate}
\end{exercise}

\begin{exercise}
For the following power series, find if they are convergent or not, and
if so find their radius of convergence.
In part f), do note the factorial in the exponent.

\medskip

\noindent
a)
$\displaystyle \sum_{n=0}^\infty 2^n x^n$
\qquad
b) $\displaystyle \sum_{n=0}^\infty n x^n$
\qquad 
c) 
$\displaystyle \sum_{n=0}^\infty n! \, x^n$
%mbxSTARTIGNORE
\qquad
%mbxENDIGNORE
%mbx <rabr/>
d) $\displaystyle \sum_{n=0}^\infty \frac{1}{(2n)!} {(x-10)}^n$
\qquad
e) $\displaystyle \sum_{n=0}^\infty x^{2n}$
\qquad
f) $\displaystyle \sum_{n=0}^\infty n! \, x^{n!}$
\end{exercise}

\begin{exercise}
Suppose $\sum_{n=0}^\infty a_n x^n$ converges for $x=1$.
\begin{enumerate}[a)]
\item
What can you say about the radius of convergence?
\item
If you further know that at $x=1$ the convergence is not absolute,
what can you say?
\end{enumerate}
\end{exercise}

\begin{exercise}
Expand
$\dfrac{x}{4-x^2}$ as a power series around $x_0 = 0$,
and compute its radius of convergence.
\end{exercise}

\begin{exercise}
\leavevmode
\begin{enumerate}[a)]
\item
Find an example where the radii of convergence of $\sum_{n=0}^\infty a_n x^n$ and
$\sum_{n=0}^\infty b_n x^n$ are both 1, but the radius of convergence of
the sum of the two series is infinite.
\item
(Trickier)
Find an example where the radii of convergence of $\sum_{n=0}^\infty a_n x^n$ and
$\sum_{n=0}^\infty b_n x^n$ are both 1, but the radius of convergence of
the product of the two series is infinite.
\end{enumerate}
\end{exercise}

\begin{exercise}
Figure out how to compute the radius of convergence using the ratio test.
That is, suppose $\sum_{n=0}^\infty a_n x^n$ is a power series and
$R \coloneqq \lim_{n\to\infty} \frac{\sabs{a_{n+1}}}{\sabs{a_n}}$ exists or is $\infty$.
Find the radius of convergence and prove your claim.
\end{exercise}

\begin{samepage}
\begin{exercise}
\leavevmode
\begin{enumerate}[a)]
\item
Prove that $\lim_{n\to\infty} n^{1/n} = 1$ using the following procedure:  Write $n^{1/n} = 1+b_n$ and
note $b_n > 0$ when $n > 1$.  Then show that ${(1+b_n)}^n \geq 
\frac{n(n-1)}{2}b_n^2$ and use this to show that $\lim\limits_{n\to\infty} b_n = 0$.
\item
Use the result of part a) to show that
if $\sum_{n=0}^\infty a_n x^n$ is a convergent power series with radius of convergence $R$,
then $\sum_{n=0}^\infty n a_n x^n$ is also convergent with the same radius of convergence.
\end{enumerate}
\end{exercise}
\end{samepage}

\begin{exnote}
There are different notions of summability (convergence)
of a series
than just the one we have seen.
A common one is \emph{\myindex{Ces\`aro summability}}%
\footnote{Named for the Italian mathematician
\href{https://en.wikipedia.org/wiki/Ernesto_Ces\%C3\%A0ro}{Ernesto Ces\`aro}
(1859--1906).}.  Let $\sum_{n=1}^\infty a_n$ be a series
and let $s_n$ be the $n$th partial sum.  The series is said to
be Ces\`aro summable to $a$ if
\begin{equation*}
a = \lim_{n\to \infty} \frac{s_1 + s_2 + \cdots + s_n}{n} .
\end{equation*}
\end{exnote}

\begin{exercise}[Challenging]
\pagebreak[2]
\leavevmode
\begin{enumerate}[a)]
\item
If $\sum_{n=1}^\infty a_n$ is convergent to $a$ (in the usual sense), show that
$\sum_{n=1}^\infty a_n$ is Ces\`aro summable (see above) to $a$.
\item
Show that in the sense of Ces\`aro $\sum_{n=1}^\infty {(-1)}^n$ is summable to
$\nicefrac{-1}{2}$.
\item
Let $a_n \coloneqq k$ when $n = k^3$ for some $k \in \N$,
$a_n \coloneqq -k$ when $n = k^3+1$ for some $k \in \N$,
otherwise
let $a_n \coloneqq 0$.  Show that $\sum_{n=1}^\infty a_n$ diverges in the usual sense
(in fact, both the sequence of terms and the partial sums are unbounded), but it is
Ces\`aro summable to 0 (seems a little paradoxical at first sight).
\end{enumerate}
\end{exercise}

\begin{exercise}[Challenging]
Show that the monotonicity in the alternating series test
is necessary.  That is, find a sequence of positive real numbers
$\{ x_n \}_{n=1}^\infty$ with $\lim_{n\to\infty} x_n = 0$ but such that
$\sum_{n=1}^\infty {(-1)}^n x_n$ diverges.
\end{exercise}

\begin{exercise}
Find a series
$\sum_{n=1}^\infty x_n$ that converges,
but $\sum_{n=1}^\infty x_n^2$ diverges.
Hint: Compare \exerciseref{exercise:squareseriesconv}.
\end{exercise}

\begin{exercise}
Suppose $\{ c_n \}_{n=1}^\infty$ is a sequence.  Prove that for every $r \in (0,1)$,
there exists a strictly increasing sequence $\{ n_k \}_{k=1}^\infty$ of natural numbers
($n_{k+1} > n_k$) such that
\begin{equation*}
\sum_{k=1}^\infty c_k x^{n_k}
\end{equation*}
converges absolutely for all $x \in [-r,r]$.
\end{exercise}

\begin{exercise}[Tonelli/Fubini for sums, challenging]
\label{exercise:tonellifubiniforsums}%
\index{Fubini for sums}\index{Tonelli for sums}%
Let $\{ x_{k,\ell} \}_{k=1,\ell=1}^\infty$ denote a doubly
indexed sequence and let $\sigma \colon \N \to \N^2$ be a bijection.
Consider the series
\begin{equation*}
\text{i)}~\sum_{i=1}^\infty x_{\sigma(i)},
\qquad
\text{ii)}~\sum_{k=1}^\infty \left( \sum_{\ell=1}^\infty x_{k,\ell} \right),
\qquad
\text{iii)}~\sum_{\ell=1}^\infty \left( \sum_{k=1}^\infty x_{k,\ell} \right) .
\end{equation*}
The expressions ii) and iii) are series of series and so we say they converge
if the inner series always converges and the outer series then also
converges.
\begin{enumerate}[a)]
\item
(Tonelli) Suppose $x_{k,\ell} \geq 0$ for all $k,\ell$.
Show that the three series i), ii), iii) either all diverge (to $\infty$)
or they all converge to the same number.
In the case of divergence, some of the \myquote{inner} series might be infinity
in which case we consider the entire sum to diverge.
\item
(Fubini) Suppose i) converges absolutely.  Show that ii) and iii)
converge and they both converge to the same number as i).
\end{enumerate}
\end{exercise}
